\documentclass[11pt,a4paper]{article}

% Standard packages
\usepackage{amsmath,amssymb,amsthm}
\usepackage{mathtools}
\usepackage{algorithm}
\usepackage{algpseudocode}
\usepackage{graphicx}
\usepackage{booktabs}
\usepackage{multirow}
\usepackage{geometry}
\usepackage{hyperref}
\usepackage{xcolor}
\usepackage{enumitem}
\usepackage{placeins}
\usepackage{float}
\usepackage[numbers,sort&compress]{natbib}
\usepackage{lineno}

\geometry{margin=1in}
\linenumbers

% Improve line breaking to reduce overfull/underfull hbox warnings
\emergencystretch=1em
\tolerance=1500
\hbadness=5000

% Theorem environments
\newtheorem{theorem}{Theorem}
\newtheorem{lemma}{Lemma}
\newtheorem{corollary}{Corollary}
\newtheorem{definition}{Definition}
\newtheorem{assumption}{Assumption}
\newtheorem{remark}{Remark}
\newtheorem{note}{Note}

% ========================================
% NOTATION STANDARDS (ML Journal Style)
% ========================================

% Sets and spaces (calligraphic)
\newcommand{\R}{\mathbb{R}}
\newcommand{\Prob}{\mathbb{P}}

% Operators (upright Roman)
\newcommand{\E}{\mathbb{E}}
\DeclareMathOperator{\Var}{Var}
\DeclareMathOperator{\Cov}{Cov}
\DeclareMathOperator{\tr}{tr}
\DeclareMathOperator{\diag}{diag}
\DeclareMathOperator*{\argmin}{arg\,min}
\DeclareMathOperator*{\argmax}{arg\,max}

% Derivatives
\newcommand{\pd}[2]{\frac{\partial #1}{\partial #2}}
\newcommand{\dd}[2]{\frac{d #1}{d #2}}

% Norms and inner products
\newcommand{\norm}[1]{\left\|#1\right\|}
\newcommand{\inner}[2]{\left\langle #1, #2 \right\rangle}


\title{Spectral Regularization and Probabilistic Calibration in Ensemble Filtering: A Unified Theoretical Framework}

\author{Rylan Spence}

\begin{document}
\maketitle

\begin{abstract}
Ensemble filtering in chaotic, partially observed systems is routinely performed with ensembles far smaller than the state dimension, where sampling noise and rank deficiency degrade both accuracy and probabilistic calibration. We develop a unified theoretical framework connecting spectral regularization in observation space to calibrated uncertainty quantification, and we analyze a deterministic alternative to stochastic observation perturbations: the Quantity-of-Interest Principal Component Analysis Ensemble Data Consistent Filter (QPCA-EnDCF). The method whitens forecast--observation residuals, identifies dominant mismatch modes via an empirical eigen-decomposition, and applies a rank-$\kappa$ truncated correction that leaves orthogonal noise-dominated directions unchanged, avoiding perturbation-induced variance injection. Our analysis separates population and finite-ensemble objects and combines covariance concentration with eigenspace perturbation bounds to obtain a bias--variance decomposition for the analysis mean.
The results show that stochastic EnKF variants incur an irreducible $\mathcal{O}(1/N)$ variance term from perturbed observations, whereas QPCA-EnDCF replaces this with variability from estimating a rank-$\kappa$ residual projector. This projector-estimation contribution is also $\mathcal{O}(1/N)$ but carries $\kappa$- and cutoff-gap-dependent prefactors through projector stability, and in effective-dimension regimes it yields rank-driven sampling behavior rather than observation-dimension-driven behavior. Numerical experiments on the Lorenz--96 system under severe undersampling (state dimension exceeding the per-observation-time observation dimension and the per-observation-time observation dimension exceeding ensemble size, with multi-step assimilation windows) confirm the theory: QPCA-EnDCF achieves near-ideal spread--skill ratios and strong temporal tracking between spread and error, while sequential and 4D stochastic EnKF remain strongly underdispersed despite inflation. The method simultaneously reduces RMSE, mitigates tuning requirements, and remains robust under non-ideal observation error assumptions, suggesting deterministic spectral regularization as a practical route to reliable uncertainty quantification in resource-constrained data assimilation.
\end{abstract}

\noindent\textbf{Keywords}: data assimilation, uncertainty quantification, probabilistic calibration, dynamical systems




% ================================================================================
% SECTION: Introduction
% ================================================================================


\section{Introduction}

Ensemble Kalman filtering provides a practical route to Bayesian state estimation in high-dimensional dynamical systems, but its uncertainty quantification can deteriorate sharply when ensemble sizes are far smaller than the state dimension. In chaotic, partially observed settings, sampling noise and rank deficiency can drive variance collapse and produce underdispersed ensembles, compromising probabilistic calibration even when point estimates remain competitive. This paper studies the interplay between regularization and calibration in ensemble data assimilation, and develops a unified theoretical and empirical account of deterministic spectral regularization in observation space.

We focus on the Quantity-of-Interest Principal Component Analysis Ensemble Data Consistent Filter (QPCA-EnDCF), a deterministic alternative to perturbed-observation stochastic EnKF variants. QPCA-EnDCF whitens forecast--observation residuals, identifies dominant mismatch modes through an empirical spectral decomposition, and applies a rank-$\kappa$ truncated correction. This construction confines updates to signal-dominated subspaces while leaving orthogonal directions unchanged, reducing perturbation-induced sampling variance and yielding a controlled bias--variance tradeoff governed by the retained spectral rank.

Our evaluation emphasizes probabilistic calibration: ensembles should be reliable as distributions, not merely accurate as point estimators. We therefore assess spread--skill behavior, rank-based reliability diagnostics, and temporal tracking between ensemble spread and realized error. We consider both sequential and four-dimensional assimilation settings. In four-dimensional formulations, observations are assimilated jointly over windows of length $LT_{\mathrm{obs}}$, and the window-endpoint analysis can incorporate information from all observations within the window. When $LT_{\mathrm{obs}}$ is shorter than the Lyapunov timescale $\tau_{\mathrm{Lyap}}$ (the inverse of the largest Lyapunov exponent, i.e., the characteristic time over which infinitesimal perturbations amplify by a factor of $e$), forecast trajectories remain correlated across the window, so later observations can constrain earlier or initial states through the joint conditioning, yielding a smoothing-like capability absent from purely sequential filters. Controlled experiments on the Lorenz--96 system demonstrate that QPCA-EnDCF maintains near-calibrated uncertainty under severe undersampling, whereas conventional stochastic filters remain strongly underdispersed despite inflation.

The remainder of the paper is organized as follows. Section~\ref{sec:problem} specifies the dynamical and observational model and describes the assimilation algorithms. Section~\ref{subsec:practical_considerations} introduces the probabilistic and spectral framework used to analyze QPCA-EnDCF, and Section~\ref{sec:concentration_spectral} establishes convergence and bias--variance results. Sections~\ref{sec:motivation}--\ref{sec:operational} present numerical experiments and operational considerations, culminating in conclusions in Section~\ref{sec:conclusion}.

\FloatBarrier





% ================================================================================
% SECTION: Problem Formulation
% ================================================================================


\section{Problem Formulation}\label{sec:problem}

This section establishes the mathematical framework within which we evaluate ensemble data assimilation methods and their uncertainty quantification properties. We begin by specifying the dynamical system that generates the true states to be estimated, the Lorenz–96 model, chosen for its chaotic dynamics that mirror challenges in operational atmospheric data assimilation. We then define the observation process through which partial, noisy measurements of the state become available, characterizing both the observation operator and the noise model. With the forward problem thus specified, we describe how ensemble methods represent posterior uncertainty through finite collections of state realizations, and we introduce the empirical quantities through which probabilistic calibration is assessed.
Finally, we specify the three data assimilation algorithms under comparison, highlighting how their distinct design choices, stochastic perturbations versus deterministic spectral regularization, affect both computational efficiency and probabilistic calibration.

 Throughout, we emphasize the connections between algorithm design, uncertainty quantification objectives, and the theoretical results developed in subsequent sections.
% ================================================================================


Observation times are indexed by $k$, with $t_k=kT_{\mathrm{obs}}$ and corresponding states $\mathbf{x}_{t_k}$ (or $\mathbf{x}_k$ when the time dependence is implicit). For four-dimensional methods, we group $L$ consecutive observation times into windows indexed by $w$, and use $\ell\in\{1,\ldots,L\}$ for within-window indices. The state dimension is $n$, the number of observed components is $m$, the stacked observation dimension is $d:=mL$, the ensemble size is $N$, and the spectral truncation rank is $\kappa$. Superscripts $f$ and $a$ denote forecast and analysis quantities, respectively; for windowed methods we use subscripts $0$ and $L$ to distinguish window-initial and window-endpoint ensembles when needed (a precise convention is stated in Remark~\ref{rem:timing_convention}).


% ================================================================================

\subsection{Observation Model}
\label{subsec:obs_problem}

Observations provide partial and noisy information about the system state at discrete assimilation times. In all experiments, we observe $m=20$ of the $n=40$ state components through a fixed linear operator $\mathbf{H}\in\mathbb{R}^{m\times n}$ that subsamples the state by retaining every second component. Specifically, $\mathbf{H}$ extracts the odd-indexed variables $x_1,x_3,\ldots,x_{39}$, leaving the even-indexed components unobserved. Observations occur at discrete times $t_k = kT_{\mathrm{obs}}$ for $k=1,2,\ldots$, and we adopt the notational convention that subscript $k$ alone (as in $\mathbf{x}_k$) and subscript $t_k$ (as in $\mathbf{x}_{t_k}$) both refer to the state at time $t_k$, with the latter form used when emphasizing the continuous-time dependence. The observation process at time $t_k$ is modeled as
\begin{equation}
\mathbf{z}_k
=
\mathbf{H}\mathbf{x}_{k}+\boldsymbol{\eta}_k,
\qquad
\boldsymbol{\eta}_k\sim\mathcal{N}(\mathbf{0},\mathbf{R}),
\label{eq:obs_model}
\end{equation}

\begin{equation}
\mathbf{z}_{(w,\ell)}
=
\mathbf{H}\mathbf{x}_{(w,\ell)}+\boldsymbol{\eta}_{(w,\ell)},
\qquad
\boldsymbol{\eta}_{(w,\ell)}\sim\mathcal{N}(\mathbf{0},\mathbf{R}),
\label{eq:obs_model_windowed}
\end{equation}
with observation error covariance $\mathbf{R}=\sigma_{\mathrm{obs}}^2\mathbf{I}_m$ and $\sigma_{\mathrm{obs}}=1.5$. The noise variables $\{\boldsymbol{\eta}_{k}\}$ and $\{\boldsymbol{\eta}_{(w,\ell)}\}$ are independent across time and independent of the state respectively, representing measurement error and unresolved representativeness effects.

This $50\%$ spatial subsampling makes the instantaneous inverse problem ill-posed. The null space $\ker(\mathbf{H})$ has dimension $20$, so infinitely many states are consistent with a single observation vector $\mathbf{z}_k$ in the noise-free case. With Gaussian observation noise \eqref{eq:obs_model}, the likelihood depends only on $\mathbf{H}\mathbf{x}$ and provides no information about components in $\ker(\mathbf{H})$.
Consequently, when combined with the forecast distribution used by an assimilation algorithm, represented in practice by the forecast ensemble and formalized in Section~\ref{subsec:practical_considerations}, that has full support (e.g., a Gaussian approximation, as opposed to the singular Lorenz–96 invariant measure), the posterior remains supported on all of $\mathbb{R}^n$, with mass in the unobserved directions determined by the prior marginal, and the unobserved components cannot be meaningfully inferred from a single observation time without additional structure.
 Temporal dynamics provide this structure: through the discrete-time evolution $\mathbf{x}_{k+1}=\mathcal{M}(\mathbf{x}_k)$, where subscript $k$ denotes the observation time index and $\mathcal{M}$ represents integration of \eqref{eq:lorenz96} over one interval $T_{\mathrm{obs}}$, observational information propagates forward in time through sequential filtering or is used to update earlier states in windowed formulations, indirectly constraining unobserved variables via their dynamical coupling to observed components. This mechanism underlies the effectiveness of sequential and windowed assimilation in partially observed chaotic systems.

The observation noise level $\sigma_{\mathrm{obs}}=1.5$, relative to the climatological standard deviation $\sigma_{\mathrm{clim}}\approx 3.6$ (defined as the temporal standard deviation of individual state components under the invariant measure of \eqref{eq:lorenz96}, averaged spatially over all $n=40$ components), yields a signal-to-noise ratio
\begin{equation}
\mathrm{SNR}
=
\frac{\sigma_{\mathrm{clim}}}{\sigma_{\mathrm{obs}}}
\approx 2.4.
\label{eq:snr}
\end{equation}
This intermediate regime places a premium on regularization. Insufficient regularization allows observational noise to contaminate the analysis and be rapidly amplified by the chaotic dynamics; excessive regularization discards informative data and produces underdispersed, overconfident posteriors. Effective methods must balance these effects. A central premise of this work is that spectral regularization achieves this balance adaptively, without reliance on ad hoc tuning.

To exploit temporal correlations, windowed methods condition jointly on observations from multiple consecutive times. We stack $L$ observations within a window into
\begin{equation}
\mathbf{z}^{(w)}
=
\begin{bmatrix}
\mathbf{z}_{(w-1)L+1}\\
\vdots\\
\mathbf{z}_{wL}
\end{bmatrix}
\in\mathbb{R}^{d},
\label{eq:stacked_obs}
\end{equation}
where $w$ indexes the window and $(w-1)L+1,\ldots,wL$ denote the global observation time indices within that window, and $d=mL$. Assuming temporally independent observation errors, the augmented observation covariance is block diagonal,
\begin{equation}
\mathbf{R}^{(L)}
=
\begin{bmatrix}
\mathbf{R} & \mathbf{0} & \cdots & \mathbf{0} \\
\mathbf{0} & \mathbf{R} & \cdots & \mathbf{0} \\
\vdots & \vdots & \ddots & \vdots \\
\mathbf{0} & \mathbf{0} & \cdots & \mathbf{R}
\end{bmatrix}
\in\mathbb{R}^{d\times d}.
\end{equation}

 To define a corresponding stacked observation operator $\mathbf{H}^{(L)}$ mapping a window-initial state to predicted stacked observations; a tangent-linear version of this operator is introduced in Definition~\ref{def:data_space}.











% ================================================================================

\subsection{Ensemble Representation}
\label{subsec:ensemble_uq}

Ensemble-based data assimilation represents posterior uncertainty through a finite collection of state realizations rather than through explicit manipulation of full probability densities or covariance operators. An ensemble consists of $N$ state vectors $\mathbf{x}^{(1)},\ldots,\mathbf{x}^{(N)}\in\mathbb{R}^n$, assembled into the ensemble matrix
\[
\mathbf{X} = [\mathbf{x}^{(1)},\ldots,\mathbf{x}^{(N)}]\in\mathbb{R}^{n\times N}.
\]
The ensemble mean
\begin{equation}
\bar{\mathbf{x}}
=
\frac{1}{N}\mathbf{X}\mathbf{1}
=
\frac{1}{N}\sum_{j=1}^N \mathbf{x}^{(j)},
\label{eq:ensemble_mean}
\end{equation}
with $\mathbf{1}\in\mathbb{R}^N$ denoting the vector of ones, provides a Monte Carlo approximation of the posterior mean and serves as the primary point estimate.

We adopt the time and window indexing conventions from Section~\ref{subsec:obs_problem}; in particular, superscripts $f$ and $a$ denote forecast and analysis ensembles, and for windowed methods we distinguish window-initial and window-endpoint ensembles when needed.

Uncertainty is encoded by fluctuations about the mean. Defining the (state) anomaly matrix
\begin{equation}
\mathbf{A}_{x}
=
\mathbf{X}-\bar{\mathbf{x}}(\mathbf{1})^\top
=
[\mathbf{x}^{(1)}-\bar{\mathbf{x}},\ldots,\mathbf{x}^{(N)}-\bar{\mathbf{x}}],
\label{eq:anomaly_matrix}
\end{equation}
the empirical covariance associated with $\mathbf{X}$ is
\begin{equation}
\widehat{\mathbf{P}}
=
\frac{1}{N-1}\mathbf{A}_{x}\mathbf{A}_{x}^\top.
\label{eq:sample_covariance}
\end{equation}
The $(N-1)^{-1}$ normalization provides the standard finite-sample correction; it is unbiased under independent sampling and remains customary in cycling experiments.
When time indexing is required, we write $\mathbf{X}_k^{f}$ and $\mathbf{X}_k^{a}$ for forecast and analysis ensembles at observation time $t_k$, and denote the corresponding anomalies and empirical covariances by $\mathbf{A}_{x,k}^{f}$, $\mathbf{A}_{x,k}^{a}$ and $\widehat{\mathbf{P}}_{k}^{\,f}$, $\widehat{\mathbf{P}}_{k}^{\,a}$.

Ensemble filters compute analysis increments through observation-space covariances derived from predicted observations. For a (possibly stacked) linear observation operator $\mathbf{H}^{(L)}$, define the corresponding ensemble of predicted observations
\[
\mathbf{Y}:=\mathbf{H}^{(L)}\mathbf{X}\in\mathbb{R}^{d\times N},
\qquad
\bar{\mathbf{y}}:=\frac{1}{N}\mathbf{Y}\mathbf{1}\in\mathbb{R}^{d},
\qquad
\mathbf{A}_{z}:=\mathbf{Y}-\bar{\mathbf{y}}\mathbf{1}^\top\in\mathbb{R}^{d\times N}.
\]
Empirical gains are then formed from cross-covariances between state anomalies and observation anomalies. For the EnDCF, these objects are constructed from windowed ensembles and stacked observations, as defined next.

\begin{definition}[EnDCF ensemble covariances and empirical gain]
\label{def:endcf_gain_ensemble}
Let $\mathbf{X}_0^{f}\in\mathbb{R}^{n\times N}$ denote the window-initial forecast ensemble, and let $\mathbf{X}_L^{f}\in\mathbb{R}^{n\times N}$ denote the forecast ensemble at the window endpoint (assimilation time). For four-dimensional methods, $\mathbf{X}_L^{f}$ is obtained by propagating $\mathbf{X}_0^{f}$ forward through $L$ observation times; for sequential methods, $\mathbf{X}_L^{f}=\mathbf{X}_0^{f}$. Let $\mathbf{Y}_{\mathrm{stack}}\in\mathbb{R}^{d\times N}$ denote the corresponding stacked forecast observations,
\[
\mathbf{Y}_{\mathrm{stack}}
:=
\bigl[
\mathbf{H}^{(L)}\mathbf{x}_0^{(1),f},
\ldots,
\mathbf{H}^{(L)}\mathbf{x}_0^{(N),f}
\bigr]
\in\mathbb{R}^{d\times N},
\]
where $\mathbf{H}^{(L)}$ is the linear stacked observation operator from Definition~\ref{def:data_space}.
Define the sample means
\[
\bar{\mathbf{x}}_L^{f}
:=
\frac{1}{N}\mathbf{X}_L^{f}\mathbf{1}\in\mathbb{R}^{n},
\qquad
\bar{\mathbf{y}}_{\mathrm{stack}}
:=
\frac{1}{N}\mathbf{Y}_{\mathrm{stack}}\mathbf{1}\in\mathbb{R}^{d},
\]
and the corresponding anomaly matrices
\[
\mathbf{A}_x
:=
\mathbf{X}_L^{f}-\bar{\mathbf{x}}_L^{f}\mathbf{1}^\top
\in\mathbb{R}^{n\times N},
\qquad
\mathbf{A}_z
:=
\mathbf{Y}_{\mathrm{stack}}-\bar{\mathbf{y}}_{\mathrm{stack}}\mathbf{1}^\top
\in\mathbb{R}^{d\times N}.
\]
The empirical cross-covariance and observation-space covariance are defined by
\[
\mathbf{P}_{xz}
:=
\frac{1}{N-1}\mathbf{A}_x\mathbf{A}_z^\top
\in\mathbb{R}^{n\times d},
\qquad
\mathbf{P}_{zz}
:=
\frac{1}{N-1}\mathbf{A}_z\mathbf{A}_z^\top
\in\mathbb{R}^{d\times d}.
\]
The EnDCF gain used in implementation is the empirical gain
\begin{equation}
\mathbf{K}^{\mathrm{DC}}
:=
\mathbf{P}_{xz}\,\mathbf{P}_{zz}^{\dagger}
\in\mathbb{R}^{n\times d},
\label{eq:endcf_gain_pinv}
\end{equation}
where $(\cdot)^{\dagger}$ denotes the Moore--Penrose pseudoinverse.
\end{definition}

The computational efficiency of ensemble methods comes with intrinsic limitations on uncertainty representation. In the small-ensemble regime $N\ll n$ typical of operational applications, ensemble covariances have rank at most $N-1$, which can lead to variance collapse and filter divergence if unmitigated (see Section~\ref{sec:inflation}). In particular, $\mathbf{P}_{zz}$ has rank at most $\min(d,N-1)$ and is therefore singular whenever $N-1<d$, which is the regime of primary interest in this work. In such cases, the gain \eqref{eq:endcf_gain_pinv} remains well defined, whereas the inverse $\mathbf{P}_{zz}^{-1}$ does not exist. More generally, one may replace $\mathbf{P}_{zz}^{\dagger}$ with a stabilized inverse, such as Tikhonov regularization $(\mathbf{P}_{zz}+\varepsilon\mathbf{I})^{-1}$; the analysis extends verbatim provided the implemented gain is uniformly controlled in operator norm (a mild moment condition is stated later in Section~\ref{subsec:regularity_assumptions}).



% ================================================================================

\section{Data Assimilation Algorithms}\label{sec:algorithms_uq}

We now describe the three ensemble-based data assimilation algorithms considered in this study. Each method addresses, in a distinct manner, the competing requirements of regularization and uncertainty preservation that arise in the small-ensemble regime. The central methodological distinction is whether observational information is incorporated through stochastic perturbations or through deterministic, spectrally regularized corrections. As we demonstrate later, this choice has significant consequences for computational efficiency, sampling error, and probabilistic calibration.

\subsection{Sequential Stochastic Ensemble Kalman Filter}

The sequential stochastic ensemble Kalman filter (EnKF) \citep{burgers1998,houtekamer1998data} serves as a baseline method against which the spectral approaches are evaluated. At each observation time, ensemble members are updated independently according to
\begin{equation}
\mathbf{x}^{(j),a}
=
\mathbf{x}^{(j),f}
+
\mathbf{K}_k\bigl(\mathbf{z}_k+\boldsymbol{\epsilon}_{k}^{(j)}-\mathbf{H}\mathbf{x}^{(j),f}\bigr) \in \mathbb{R}^{n},
\qquad j=1,\ldots,N,
\label{eq:enkf_update}
\end{equation}
where superscripts $f$ and $a$ denote forecast and analysis states, respectively. The random vectors $\boldsymbol{\epsilon}_{k}^{(j)}\sim\mathcal{N}(\mathbf{0},\mathbf{R})$ are independent observation perturbations sampled separately for each ensemble member. The Kalman gain is computed from empirical forecast statistics as
\begin{equation}
\mathbf{K}_k
=
\widehat{\mathbf{P}}^{\,f}_{k}\mathbf{H}^\top
\left(\mathbf{H}\widehat{\mathbf{P}}_{k}^{\,f}\mathbf{H}^\top+\mathbf{R}\right)^{-1}.
\label{eq:kalman_gain}
\end{equation}
By forming the gain in observation space, this expression avoids explicit inversion of the state-space covariance $\mathbf{P}^f$, reducing computational cost and improving numerical stability when $m<n$.

Observation perturbations play a dual role in the stochastic EnKF. First, they ensure second-moment consistency in expectation: taking the expectation over the observation perturbations $\{\boldsymbol{\epsilon}_{k}^{(j)}\}$ conditional on the forecast ensemble, the empirical analysis covariance satisfies
\[
\mathbb{E}_{\boldsymbol{\epsilon}}[\widehat{\mathbf{P}}_{k}^{\,a}]=(\mathbf{I}-\mathbf{K}_k\mathbf{H})\widehat{\mathbf{P}}_{k}^{\,f}(\mathbf{I}-\mathbf{K}_k\mathbf{H})^\top + \mathbf{K}_k\mathbf{R}\mathbf{K}_k^\top
\]
in agreement with the Kalman update formula, where $\widehat{\mathbf{P}}^{\,a}$ denotes the sample covariance of the analysis ensemble. Second, by presenting each ensemble member with a distinct realization of the observation, the perturbations prevent collapse of the ensemble onto a single point when the same observation is assimilated repeatedly. These benefits, however, come at the cost of additional sampling variability. The perturbation-induced contribution to the ensemble mean error has expected squared norm $\mathcal{O}(\|\mathbf{K}_k\mathbf{R}^{1/2}\|_F^2/N)$ where $\|\cdot \|_F$ denotes the Frobenius norm, which scales with the gain magnitude and observation error covariance. This sampling noise can substantially degrade both point estimates and uncertainty quantification for small ensemble sizes (see Section~\ref{subsec:spread_skill}).

Regularization in the stochastic EnKF arises implicitly through the additive observation covariance $\mathbf{R}$ in the innovation term. This mechanism inflates all eigenvalues of $\mathbf{H}\mathbf{P}^f\mathbf{H}^\top$ uniformly by $\sigma_{\mathrm{obs}}^2$, stabilizing the inversion but failing to distinguish between signal-dominated and noise-dominated subspaces. The resulting regularization is robust but conservative, prioritizing numerical stability over adaptivity.

\subsection{Four-Dimensional Stochastic Ensemble Kalman Filter}

The four-dimensional stochastic ensemble Kalman filter (4D-EnKF) \citep{hunt2004four,evensen2009book} extends the sequential formulation by conditioning jointly on observations collected over an assimilation window of length $L$. The ensemble is propagated forward through $L$ observation times, generating forecasts $\{\mathbf{x}_\ell^{(j),f}\}_{\ell=1}^L$ for each member $j$. The update targets the window-final state $\mathbf{x}_L^{(j),f}$ using information from all $L$ observations. Within each window, ensemble members are updated using stacked, perturbed observations
\begin{equation}
\tilde{\mathbf{z}}^{(w,j)}
=
\mathbf{z}^{(w)}+\boldsymbol{\epsilon}^{(w,j)}\in \mathbb{R}^{d},
\qquad
\boldsymbol{\epsilon}^{(w,j)}\sim\mathcal{N}(\mathbf{0},\mathbf{R}^{(L)}),
\label{eq:4d_perturb}
\end{equation}
where $\mathbf{z}^{(w)}$ is the stacked observation vector defined in \eqref{eq:stacked_obs} and $\mathbf{R}^{(L)}$ is the corresponding block-diagonal covariance. The Kalman gain is computed analogously to the sequential case, but now in the augmented observation space $\mathbb{R}^{d}$, relating the window-final state to the stacked observations through cross-covariances computed from the ensemble of forecasts at all $L$ times.

The principal advantage of four-dimensional formulations is their ability to exploit temporal correlations. Observations acquired later in the window can inform earlier states through the dynamical model, effectively smoothing the trajectory and indirectly constraining unobserved components. However, the increase in observation-space dimensionality amplifies sampling variance; the net benefit depends on the interplay between improved conditioning and increased sampling noise (see Section~\ref{subsec:operational_window_length} for detailed analysis).

\subsection{Spectral Regularization Through QPCA-EnDCF}

The Quantity-of-Interest Principal Component Analysis Ensemble Data Consistent Filter (QPCA-EnDCF) adopts a fundamentally different strategy, eliminating observation perturbations altogether in favor of deterministic, spectrally regularized corrections. The method proceeds through three stages: whitening of forecast--observation residuals, spectral identification of dominant mismatch modes, and construction of truncated corrections.

For clarity, we describe the construction here at the algorithmic level. The population and sample objects used in the theoretical analysis are introduced in Definitions~\ref{def:whitened_residual_covariances}--\ref{def:spectral_decompositions}, and the implementation details are summarized in Algorithm~\ref{alg:qpca_endcf} and discussed further in Section~\ref{subsec:practical_considerations}.

In the first stage, residuals are normalized by observation uncertainty to obtain whitened mismatches. For the stacked four-dimensional observation framework with diagonal observation error covariance $\mathbf{R}^{(L)}=\sigma_{\mathrm{obs}}^2\mathbf{I}_{d}$, this takes the form
\begin{equation}
\mathbf{E}
=
\sigma_{\mathrm{obs}}^{-1}
\bigl(\mathbf{Y}_{\mathrm{stack}}-\mathbf{z}^{(w)}\mathbf{1}^\top\bigr)
=
(\mathbf{R}^{(L)})^{-1/2}
\bigl(\mathbf{Y}_{\mathrm{stack}}-\mathbf{z}^{(w)}\mathbf{1}^\top\bigr),
\label{eq:whitening_problem}
\end{equation}
where $\mathbf{Y}_{\mathrm{stack}}\in\mathbb{R}^{d\times N}$ contains the stacked forecast observations for all ensemble members. The scalar form $\sigma_{\mathrm{obs}}^{-1}$ is used when $\mathbf{R}^{(L)}$ is diagonal; implementation details for non-diagonal covariances are provided in Section~\ref{subsec:correlated_errors}.

In the second stage, the covariance of the centered whitened residuals is decomposed spectrally,
\begin{equation}
\mathbf{C}_E
=
\frac{1}{N-1}\mathbf{E}_c\mathbf{E}_c^\top
=
\sum_{i=1}^{r}\hat{\lambda}_i\hat{\mathbf{v}}_i\hat{\mathbf{v}}_i^\top = \hat{\mathbf{V}}_r \hat{\boldsymbol{\Lambda}}_r \hat{\mathbf{V}}_r^\top,
\label{eq:spectral_decomp_problem}
\end{equation}
where $r=\min(d,N-1)$ denotes the maximum possible rank of $\mathbf{C}_E$, and $\mathbf{E}_c:=\mathbf{E}-\bar{\mathbf{e}}\mathbf{1}^\top$ denotes the centered whitened residual matrix, with $\mathbf{E}=[\mathbf{e}^{(1)},\ldots,\mathbf{e}^{(N)}]\in\mathbb{R}^{d\times N}$ and $\bar{\mathbf{e}}=\frac{1}{N}\sum_{j=1}^N\mathbf{e}^{(j)}$. 
In practice, the positive eigenvalues quantify the variance of forecast--observation mismatch along orthogonal directions $\hat{\mathbf{v}}_i$. Large eigenvalues correspond to coherent, dynamically meaningful discrepancies, whereas small (but nonzero) eigenvalues primarily reflect sampling and measurement noise.

In the final stage, corrections are constructed by projecting onto the leading $\kappa$ eigenmodes,
\begin{equation}
\mathbf{Q}_{\mathrm{PCA}}
=
-\hat{\mathbf{V}}_\kappa\hat{\mathbf{V}}_\kappa^\top\mathbf{E},
\label{eq:qpca_correction}
\end{equation}
where $\hat{\mathbf{V}}_\kappa=[\hat{\mathbf{v}}_1,\ldots,\hat{\mathbf{v}}_\kappa]$ and $\kappa\ll\min(d,N-1)$. Note that eigenvectors $\hat{\mathbf{v}}_i$ are computed from the centered residuals $\mathbf{E}_c$ (characterizing covariance structure), but the projection is applied to the uncentered residuals $\mathbf{E}$ (containing both mean mismatch and fluctuations). This allows correction of both systematic bias (captured by $\bar{\mathbf{e}}$) and random fluctuations about the mean, but only to the extent that these components lie in the retained $\kappa$-dimensional eigenspace. Components of the mean mismatch orthogonal to $\{\hat{\mathbf{v}}_1,\ldots,\hat{\mathbf{v}}_\kappa\}$ are not corrected, which can introduce bias if the mean mismatch does not align well with the leading covariance directions. This truncation implements hard spectral regularization in observation space, retaining dominant mismatch directions while suppressing noise-dominated components.

Because $\mathbf{Q}_{\mathrm{PCA}}$ is defined in whitened coordinates, we unwhiten to obtain the stacked observation-space increment
\begin{equation}
\boldsymbol{\Delta}_{\mathrm{obs}}
:=
(\mathbf{R}^{(L)})^{1/2}\mathbf{Q}_{\mathrm{PCA}}
=
-(\mathbf{R}^{(L)})^{1/2}\hat{\mathbf{V}}_\kappa\hat{\mathbf{V}}_\kappa^\top\mathbf{E}.
\label{eq:qpca_delta_obs}
\end{equation}
The increment $\boldsymbol{\Delta}_{\mathrm{obs}}$ is then mapped back to state space via the empirical cross-covariance gain matrix $\mathbf{K}^{\mathrm{DC}}=\mathbf{P}_{xz}\mathbf{P}_{zz}^{\dagger}$, which projects observation-space innovations onto the state-space ensemble covariance structure, and the resulting state-space increments are applied deterministically to the ensemble (see Definition~\ref{def:endcf_gain_ensemble} and Algorithm~\ref{alg:qpca_endcf} for details).

We now introduce the covariance objects that underlie the QPCA-EnDCF update. These objects are constructed from whitened forecast--observation residuals and admit both population-level definitions, corresponding to idealized infinite-ensemble limits, and sample-level counterparts computed from finite ensembles. Making this distinction explicit is essential for the theoretical analysis that follows, as the algorithm operates on sample quantities whose behavior must be related to their population analogues.

\begin{definition}[Whitened residuals and associated covariances]
\label{def:whitened_residual_covariances}
Let $\mathbf{R}^{(L)}\in\mathbb{R}^{d\times d}$ be symmetric positive definite. Consider a fixed realization of the observation vector $\mathbf{z}^{(w)}\in\mathbb{R}^{d}$, which we treat as deterministic throughout this definition. Let $\mathbf{e}^{(j)}\in\mathbb{R}^{d}$ denote the $j$th column of the whitened residual matrix $\mathbf{E}$ defined in \eqref{eq:whitening_problem}. The corresponding population mean and covariance, conditional on the realized observations $\mathbf{z}^{(w)}$, are given by
\begin{align*}
\boldsymbol{\mu}_E
&:=
\mathbb{E}[\mathbf{e}^{(j)}\mid\mathbf{z}^{(w)}]
=
(\mathbf{R}^{(L)})^{-1/2}
\bigl(\mathbf{H}^{(L)}\boldsymbol{\mu}_0^{f}-\mathbf{z}^{(w)}\bigr),\\
\boldsymbol{\Sigma}_E
&:=
\mathrm{Cov}(\mathbf{e}^{(j)}\mid\mathbf{z}^{(w)})
=
(\mathbf{R}^{(L)})^{-1/2}
\mathbf{H}^{(L)}\mathbf{P}^{f}(\mathbf{H}^{(L)})^\top
(\mathbf{R}^{(L)})^{-1/2}.
\end{align*}
Here the expectation is taken over the forecast distribution $P_{\mathcal{X}}$ (i.e., over ensemble draws), with the observations held fixed at their realized values. Note that $\boldsymbol{\Sigma}_E$ does not depend on $\mathbf{z}^{(w)}$ under this conditioning, since the covariance structure is determined solely by the forecast distribution. The sample mean $\bar{\mathbf{e}}$ and sample covariance $\mathbf{C}_E$ are defined from $\{\mathbf{e}^{(j)}\}_{j=1}^N$ as in \eqref{eq:spectral_decomp_problem}.
\end{definition}
In the theoretical development that follows, we work conditionally on the realized observation vector $\mathbf{z}^{(w)}$ (equivalently, we treat $\mathbf{z}^{(w)}$ as fixed and take expectations only over ensemble sampling and any algorithmic randomness). To reduce notation, we typically suppress the explicit conditioning ``$\mid\mathbf{z}^{(w)}$'' in expressions such as $\mathbb{E}[\cdot]$ and $\mathrm{Cov}(\cdot)$, except where it is needed for clarity.
The transformation $(\mathbf{R}^{(L)})^{-1/2}$ rescales the observation space so that, in whitened coordinates, the (stacked) observation errors are isotropic with unit covariance. This normalization is standard in data assimilation and greatly simplifies both the geometric interpretation of residuals and the derivation of moment and concentration bounds for sample covariance operators.

\begin{definition}[Population and sample spectral decompositions]
\label{def:spectral_decompositions}
Let $d:=mL$. The population covariance $\boldsymbol{\Sigma}_E$ admits a spectral decomposition
\[
\boldsymbol{\Sigma}_E
=
\sum_{i=1}^{d}\lambda_i\,\mathbf{v}_i\mathbf{v}_i^\top,
\qquad
\lambda_1\ge\cdots\ge\lambda_d\ge 0.
\]
Let $\{(\hat{\lambda}_i,\hat{\mathbf{v}}_i)\}_{i=1}^{d}$ denote the eigenpairs of the sample covariance $\mathbf{C}_E$ from \eqref{eq:spectral_decomp_problem}, ordered so that $\hat{\lambda}_1\ge\cdots\ge\hat{\lambda}_d\ge 0$, with the convention that $\hat{\lambda}_i=0$ for $i>\min(d,N{-}1)$. For any truncation level $1\le\kappa\le \min(d,N-1)$, define the truncated covariances
\[
\boldsymbol{\Sigma}_E^{\kappa}
:=
\sum_{i=1}^{\kappa}\lambda_i\,\mathbf{v}_i\mathbf{v}_i^\top,
\qquad
\mathbf{C}_E^{\kappa}
:=
\sum_{i=1}^{\kappa}\hat{\lambda}_i\,\hat{\mathbf{v}}_i\hat{\mathbf{v}}_i^\top.
\]
These truncated operators play a central role in the construction and analysis of the spectrally regularized EnDCF update.
\end{definition}



\subsection{Practical Considerations}
\label{subsec:practical_considerations}

Because the update is deterministic, it introduces no additional sampling noise from perturbed observations. Each ensemble member is corrected by a deterministic function of its whitened residual, and the correction acts only in the $\kappa$-dimensional subspace spanned by the leading empirical eigenmodes, leaving the complementary $(d-\kappa)$-dimensional subspace unchanged. In typical configurations, this implies that the vast majority of observation-space directions remain unaffected by the update, concentrating corrections on the dominant mismatch mode while avoiding noise-dominated directions. This stands in sharp contrast to stochastic filters, in which random perturbations can distort ensemble geometry in directions unrelated to posterior uncertainty. The adaptive nature of the spectral truncation is central to the method’s regularization properties. By selecting modes based on the realized eigenvalue spectrum, QPCA-EnDCF balances data fidelity against noise suppression without introducing fixed inflation or localization parameters. Although the truncation level $\kappa$ must be chosen, we show in subsequent sections that performance is relatively insensitive to its precise value over a broad range, and that simple spectral heuristics provide effective guidance. The theoretical analysis that follows formalizes these geometric and statistical considerations, establishing rigorous results on variance preservation, sampling-error reduction, and asymptotic calibration for spectrally regularized ensemble methods.

Windowed assimilation also amplifies sampling effects in stochastic ensemble methods. Observation perturbations must be drawn in the higher-dimensional space $\mathbb{R}^{d}$. Under the assumption that perturbations $\{\boldsymbol{\epsilon}^{(w,j)}\}_{j=1}^N$ are i.i.d.\ with $\boldsymbol{\epsilon}^{(w,j)}\sim\mathcal{N}(\mathbf{0},\mathbf{R}^{(L)})$ and $\mathbf{R}^{(L)}=\sigma_{\mathrm{obs}}^2\mathbf{I}_{d}$, consider the (unbiased) sample covariance
\[
\widehat{\mathbf{R}}^{(L)}
:=
\frac{1}{N-1}\sum_{j=1}^N
(\boldsymbol{\epsilon}^{(w,j)}-\bar{\boldsymbol{\epsilon}}^{(w)})
(\boldsymbol{\epsilon}^{(w,j)}-\bar{\boldsymbol{\epsilon}}^{(w)})^\top,
\qquad
\bar{\boldsymbol{\epsilon}}^{(w)}:=\frac{1}{N}\sum_{j=1}^N \boldsymbol{\epsilon}^{(w,j)}.
\]
Then, with $d=mL$, in the Gaussian case one has
\[
\left(\mathbb{E}\bigl[\|\widehat{\mathbf{R}}^{(L)}-\mathbf{R}^{(L)}\|_F^2\bigr]\right)^{1/2}
=
\sigma_{\mathrm{obs}}^2\sqrt{\frac{d(d+1)}{N-1}}
=
\mathcal{O}\!\left(\sigma_{\mathrm{obs}}^2\,\frac{d}{\sqrt{N}}\right),
\]
and hence, by Jensen's inequality,
$\mathbb{E}\|\widehat{\mathbf{R}}^{(L)}-\mathbf{R}^{(L)}\|_F=\mathcal{O}(\sigma_{\mathrm{obs}}^2\,d/\sqrt{N})$. Relative to the sequential case ($d=m$), the prefactor grows like $L$, so for moderate ensemble sizes this $L$ amplification can offset gains from temporal coupling. Quantifying and mitigating this tradeoff between information gain and sampling noise motivates the comparative numerical analyses presented in Section~\ref{sec:motivation}.



This section introduces the mathematical framework and notation underlying the analysis of QPCA-EnDCF. We distinguish between population-level quantities, which are deterministic and defined with respect to the forecast distribution, and sample-level quantities, which are computed from a finite ensemble and are therefore inherently random. Population quantities represent the idealized limit of an infinite ensemble; sample quantities fluctuate around their population analogues due to finite-$N$ sampling effects. A primary goal of the theoretical results developed later is to quantify this discrepancy by establishing concentration bounds. All random quantities are defined on an underlying probability space $(\Omega,\mathcal{F},\mathbb{P})$, and all expectations are taken with respect to $\mathbb{P}$ unless otherwise stated. In the small-ensemble regime, the empirical observation-space covariance is necessarily singular, requiring a stabilized inverse (see Subsection~\ref{subsec:ensemble_uq}). We proceed by formalizing the forecast and observation spaces and the stacked QoI map, then defining whitened residual covariances and their spectra, and finally introducing the empirical EnDCF gain and QPCA projectors used by the algorithm.

\par\smallskip\noindent\refstepcounter{remark}\textbf{Remark~\theremark}
\label{rem:timing_convention}
In four-dimensional methods, two distinct time indices within a window must be distinguished:
\begin{enumerate}[label=(\roman*),nosep]
\item \textbf{Window-initial states} $\mathbf{x}_0^{(j),f}$: These are the forecast states at the beginning of a window ($\ell=0$). The stacked observation operator $\mathbf{H}^{(L)}$ (Definition~\ref{def:data_space}) is applied to $\mathbf{x}_0^{(j),f}$ to produce predicted observations at times $\ell=1,\ldots,L$ via the tangent-linear propagation. The whitened residuals $\mathbf{e}^{(j)}$ are functions of these window-initial states.
\item \textbf{Window-endpoint states} $\mathbf{x}_L^{(j),f}$: These are the forecast states at the end of the window ($\ell=L$), obtained by propagating $\mathbf{x}_0^{(j),f}$ through $L$ applications of $\mathcal{M}$ (i.e., over $L$ observation intervals of length $T_{\mathrm{obs}}$). The EnDCF gain $\mathbf{K}^{\mathrm{DC}}$ (Definition~\ref{def:endcf_gain_ensemble}) maps observation-space corrections to state-space increments that are applied to $\mathbf{x}_L^{(j),f}$ to produce the analysis $\mathbf{x}_L^{(j),a}$.
\end{enumerate}
Throughout this paper, we use the subscript convention explicitly: $\mathbf{x}_0^{(j),f}$ always denotes the window-initial state (used for residual and $\mathbf{Y}_{\mathrm{stack}}$ computation), while $\mathbf{x}_L^{(j),f}$ denotes the window-endpoint state (to which the gain is applied). When $\mathbf{x}^{(j),f}$ appears without a time subscript in theoretical statements, it refers to the window-initial state unless the gain matrix $\mathbf{K}$ appears in the same expression, in which case it refers to the window-endpoint state.
\par\smallskip
\noindent We use the corresponding notation for forecast means: $\boldsymbol{\mu}_0^{f}:=\mathbb{E}[\mathbf{x}_0^{(j),f}]$ and $\boldsymbol{\mu}_L^{f}:=\mathbb{E}[\mathbf{x}_L^{(j),f}]$.
\par\smallskip


We now formalize the spectral projection operators that implement dimensional truncation in QPCA-EnDCF. These projectors isolate the dominant modes of forecast--observation mismatch and provide the mechanism through which spectral regularization is imposed in observation space.

\begin{definition}[Spectral projectors and QPCA truncation]
\label{def:projectors_truncation}
Let $\mathbf{C}\in\mathbb{R}^{d\times d}$ be a symmetric matrix with orthonormal eigenvectors $\{\mathbf{v}_i(\mathbf{C})\}_{i=1}^{d}$, ordered according to nonincreasing eigenvalues. For any integer $1\le\kappa\le d$, define the rank-$\kappa$ orthogonal projector
\[
\mathbf{P}_\kappa(\mathbf{C})
:=
\sum_{i=1}^{\kappa}\mathbf{v}_i(\mathbf{C})\mathbf{v}_i(\mathbf{C})^\top.
\]
When $\mathbf{C}=\boldsymbol{\Sigma}_E$ we write $\mathbf{P}_\kappa:=\mathbf{P}_\kappa(\boldsymbol{\Sigma}_E)$, and when $\mathbf{C}=\mathbf{C}_E$ we write $\hat{\mathbf{P}}_\kappa:=\mathbf{P}_\kappa(\mathbf{C}_E)$. In QPCA-EnDCF, the empirical projector $\hat{\mathbf{P}}_\kappa$ is used to restrict observation-space increments to a low-dimensional subspace spanned by the leading empirical eigenmodes.
\end{definition}
If $\mathbf{C}$ has repeated eigenvalues at or near the cutoff (e.g., $\lambda_\kappa=\lambda_{\kappa+1}$), the ordered eigenvectors are not unique, so the rank-$\kappa$ subspace may be nonunique. In such cases, $\mathbf{P}_\kappa(\mathbf{C})$ is understood as an orthogonal projector onto some $\kappa$-dimensional invariant subspace associated with the $\kappa$ largest eigenvalues; any orthonormal basis for that subspace yields the same projector. In the main results, we impose a cutoff gap $\lambda_\kappa-\lambda_{\kappa+1}\ge \delta_\kappa>0$ (Theorem~\ref{thm:bias_variance}), which ensures the leading $\kappa$-dimensional invariant subspace, and hence $\mathbf{P}_\kappa$ and $\hat{\mathbf{P}}_\kappa$ are well defined and stable under perturbations.
The effect of spectral truncation depends critically on the alignment between the mean whitened innovation and the retained subspace. Quantities such as $\|(\mathbf{I}-\mathbf{P}_\kappa)\boldsymbol{\mu}_E\|^2$ (or their analogues for approximate maps $Q_s$) measure the component of the mean innovation discarded by truncation. In general, this error cannot be controlled solely by an eigenvalue-tail bound of the form $\sum_{i>\kappa}\lambda_i$ without additional structural assumptions linking the mean $\boldsymbol{\mu}_E$ to the covariance $\boldsymbol{\Sigma}_E$. This distinction plays an important role in the theoretical analysis of truncation bias presented in later sections.




% -----------------------------------------------------------------------------

% ================================================================================


% ================================================================================
% SECTION 4: Illustrative Example
% ================================================================================
\section{Illustrative Example}\label{sec:motivation}


\subsection{Lorenz--96 Dynamical System}
\label{subsec:lorenz96_problem}

We use the Lorenz--96 model \citep{lorenz1996predictability} as the canonical testbed for all numerical experiments. Originally introduced as a low-dimensional surrogate for midlatitude atmospheric dynamics, the system has become a standard benchmark in data assimilation because it combines sustained chaos with modest computational cost. This balance enables systematic investigation of uncertainty quantification under long sequences of forecast--analysis cycles.

The model consists of $n=40$ state variables $\mathbf{x}=(x_1,\ldots,x_{40})^\top\in\mathbb{R}^{40}$ governed by
\begin{equation}
\frac{\mathrm{d}x_i}{\mathrm{d}t}
=
(x_{i+1}-x_{i-2})x_{i-1}-x_i+F,
\qquad i=1,\ldots,40,
\label{eq:lorenz96}
\end{equation}
with periodic indexing modulo $40$ and constant forcing $F=8.0$. The dynamics reflect the interaction of three components: a quadratic advection term $(x_{i+1}-x_{i-2})x_{i-1}$ that transfers energy across scales and induces strong nonlinearity, linear damping $-x_i$ that prevents unbounded growth, and constant forcing that maintains a statistically stationary, nonequilibrium regime.

Time integration of \eqref{eq:lorenz96} is performed using a classical fourth-order Runge--Kutta scheme with timestep $\Delta t=0.01$. This choice renders numerical discretization error negligible relative to observation noise and ensemble sampling error. Consequently, performance differences observed in the experiments can be attributed to algorithmic and statistical effects rather than time-integration artifacts.

This section specifies the experimental setup used to evaluate ensemble data assimilation methods and introduces diagnostics for probabilistic calibration. The configuration is chosen to stress uncertainty quantification under severe ensemble undersampling, so that performance is assessed not only by point accuracy but also by distributional fidelity.

\subsection{Experimental Configuration}\label{subsec:benchmark}

Unless stated otherwise, experiments use ensemble size $N=10$, observe $m=20$ of $n=40$ state variables, and adopt observation noise level $\sigma_{\mathrm{obs}}=1.5$. This extreme undersampling ($N\ll n$) provides a stringent test of whether algorithms preserve probabilistic calibration despite structural rank limitations.

Observations are generated by
\[
\mathbf{z}_k=\mathbf{H}\mathbf{x}^{\mathrm{true}}_k+\boldsymbol{\eta}_k,
\qquad 
\boldsymbol{\eta}_k\sim\mathcal{N}(\mathbf{0},\sigma_{\mathrm{obs}}^2\mathbf{I}_m),
\]
where $\mathbf{H}\in\mathbb{R}^{m\times n}$ extracts every second state component. Observations are assimilated every $T_{\mathrm{obs}}=0.1$ time units. Relative to climatological variability $\sigma_{\mathrm{clim}}\approx3.6$, the resulting signal-to-noise ratio is $\mathrm{SNR}\approx2.4$.

Windowed methods assimilate windows of $L=5$ consecutive observations, spanning $LT_{\mathrm{obs}}=0.5$ time units. This corresponds to approximately $0.83\,\tau_{\mathrm{Lyap}}$ with $\tau_{\mathrm{Lyap}}\approx0.6$, yielding a temporal autocorrelation of state variables at lag $0.5$ of approximately $0.61$ (computed under the Lorenz--96 invariant measure at $F=8$). We perform $W=50$ windows, for a total of $K=WL=250$ observation times. For four-dimensional methods, analysis updates occur once per window (at window endpoints), yielding $W=50$ analysis states. For sequential methods, updates occur at every observation time, yielding $K=WL=250$ analysis states. To enable fair comparison of time-series calibration, spread and RMSE are evaluated at window endpoints only. We index these evaluation times by $w\in\{1,\ldots,W\}$, corresponding to global observation-time indices $k_w:=wL\in\{L,2L,\ldots,WL\}=\{5,10,\ldots,250\}$. Equations~\eqref{eq:spread_framework}--\eqref{eq:rmse_framework} therefore aggregate over $w=1,\ldots,W$. This convention ensures that (i) 4D methods are evaluated at the times when their analysis is actually computed, and (ii) sequential methods are evaluated at comparable temporal spacing, avoiding artifacts from dense versus sparse evaluation. Algorithm-specific parameters are fixed across experiments. QPCA-EnDCF uses truncation level $\kappa=1$ (retaining only the leading eigenmode) and no multiplicative inflation ($\lambda_{\mathrm{infl}}=1.00$), consistent with eigenvalue spectra in which the dominant mode captures $60$--$80\%$ of residual variance. Stochastic methods employ multiplicative inflation $\lambda_{\mathrm{infl}}=1.05$ and perturbed observations drawn from $\mathcal{N}(\mathbf{0},\mathbf{R})$ for sequential updates or $\mathcal{N}(\mathbf{0},\mathbf{R}^{(L)})$ for four-dimensional updates.



\subsection{Probabilistic Calibration Results}\label{subsec:calibration}

This section presents the central empirical investigation of this work: the assessment of probabilistic calibration in ensemble-based data assimilation under severe undersampling. The preceding sections developed the theoretical framework predicting that spectral regularization should maintain calibrated uncertainty estimates despite finite ensemble size, and established the experimental configuration designed to stress-test these predictions. We now examine whether QPCA-EnDCF achieves the distributional properties necessary for reliable uncertainty quantification, contrasting its performance with standard stochastic ensemble methods.

The fundamental question we address is whether ensemble uncertainty estimates accurately reflect true estimation error, both in magnitude and in temporal dynamics. This question transcends point estimation accuracy. An ensemble method that produces accurate state estimates but severely underestimates uncertainty provides misleading probabilistic information, rendering ensemble forecasts unsuitable for risk-informed decision making. Conversely, a method exhibiting larger point errors but well-calibrated uncertainty quantifies estimation reliability honestly, enabling rational decisions under uncertainty. Probabilistic calibration thus represents a distinct dimension of performance, complementary to but independent of mean-squared error.

We structure the analysis around two complementary diagnostic perspectives, each revealing different aspects of distributional calibration. Section~\ref{subsec:spread_skill} examines spread--skill relationships, which probe whether ensemble variance exhibits the magnitude and temporal correlation with mean-squared error expected under second-moment calibration. Section~\ref{subsec:rank_histograms} employs rank histogram analysis to assess full distributional calibration without assuming Gaussianity, testing whether the ensemble satisfies the exchangeability property implied by proper probabilistic representation. These diagnostics provide complementary information: spread--skill relationships directly connect uncertainty estimates to error but assume elliptical distributions; rank histograms relax parametric assumptions but aggregate information across state components. Together, they provide a comprehensive assessment of whether ensemble distributions reliably quantify estimation uncertainty.

All experiments reported in this section employ five independent Monte Carlo realizations under the baseline configuration specified in Section~\ref{sec:motivation}. Diagnostics are evaluated at each window endpoint, yielding $W=50$ evaluation points per trial (and $5W=250$ window-endpoint samples per algorithm when pooling across trials). When reporting summary scalars (e.g., Table~\ref{tab:calibration_metrics}), we compute each metric within a trial by aggregating over its $W$ windows, then report mean $\pm$ standard deviation across the five trials.


% ================================================================================

\subsection{Spread--Skill Relationships}\label{subsec:spread_skill}

Spread--skill diagnostics assess whether ensemble-estimated uncertainty is consistent with actual estimation error. A calibrated ensemble should satisfy two criteria: its spread should match the RMSE in magnitude (scale consistency) and moves together with it over time (dynamic tracking).

We quantify posterior uncertainty by the time-averaged ensemble spread
\begin{equation}
\sigma_{\mathrm{ens}}
=
\left[
\frac{1}{W}\sum_{w=1}^W \frac{1}{n}\,\tr\!\bigl(\widehat{\mathbf{P}}_{wL}^{\,a}\bigr)
\right]^{1/2},
\label{eq:spread_framework}
\end{equation}
where $\widehat{\mathbf{P}}_{wL}^{\,a} = \frac{1}{N-1}{\mathbf{A}^{a}_{wL}}{\mathbf{A}^{a}_{wL}}^{\!\top}$ is the empirical analysis covariance at the endpoint $k=wL$, formed from ensemble anomalies $\mathbf{A}^{a}_{wL} = \mathbf{X}_{wL}^a - \bar{\mathbf{x}}_{wL}^a\mathbf{1}^{\top}$. Estimation accuracy is quantified by the root-mean-square error
\begin{equation}
\mathrm{RMSE}
=
\left[
\frac{1}{W}\sum_{w=1}^W \frac{1}{n}\|\bar{\mathbf{x}}_{wL}^a - \mathbf{x}_{wL}^{\mathrm{true}}\|^2
\right]^{1/2},
\label{eq:rmse_framework}
\end{equation}
where $\bar{\mathbf{x}}_{wL}^a$ is the ensemble mean analysis at the endpoint $k=wL$ and $\mathbf{x}_{wL}^{\mathrm{true}}$ is the corresponding true state.

Dynamic tracking is assessed by examining the cycle-wise (per-window) spread and RMSE sequences. Define the per-cycle spread
\[
\sigma_w := \left[\frac{1}{n}\,\tr\!\bigl(\widehat{\mathbf{P}}_{wL}^{\,a}\bigr)\right]^{1/2}
\]
and the per-cycle RMSE
\[
\mathrm{RMSE}_w := \left[\frac{1}{n}\|\bar{\mathbf{x}}_{wL}^a - \mathbf{x}_{wL}^{\mathrm{true}}\|^2\right]^{1/2}.
\]
Calibration is summarized by the cycle-wise spread--skill ratios
\begin{equation}
\gamma_w := \frac{\sigma_w}{\mathrm{RMSE}_w},
\qquad
\bar{\gamma} := \frac{1}{W}\sum_{w=1}^W \gamma_w,
\label{eq:gamma_framework}
\end{equation}
for which ideal calibration corresponds to $\gamma_w \approx 1$ (and hence $\bar{\gamma}\approx 1$); values $\gamma_w \ll 1$ and $\gamma_w \gg 1$ indicate overconfidence and underconfidence, respectively.
The Pearson correlation $\rho$ between the sequences $\{\sigma_w\}_{w=1}^W$ and $\{\mathrm{RMSE}_w\}_{w=1}^W$ measures whether ensemble uncertainty tracks estimation error across time: calibrated ensembles exhibit $\rho$ close to 1.

Table~\ref{tab:calibration_metrics} reports, for each algorithm, calibration statistics summarized across five independent trials under the baseline configuration, including mean ensemble spread, RMSE, average spread--skill ratio $\bar{\gamma}$, and the Pearson correlation $\rho$ between the window-wise spread and RMSE sequences.

\begin{table}[htbp]
\centering
\caption{Probabilistic calibration metrics under the baseline configuration ($N=10$, $L=5$, $\sigma_{\mathrm{obs}}=1.5$). Values are mean $\pm$ standard deviation across five independent trials, where each trial-level statistic is computed by aggregating over its $W=50$ windows.}
\label{tab:calibration_metrics}
\begin{tabular}{lcccc}
\toprule
\textbf{Algorithm} & \textbf{Spread $\sigma_{\mathrm{ens}}$} & \textbf{RMSE} & \textbf{Ratio $\bar{\gamma}$} & \textbf{Correlation $\rho$} \\
\midrule
Sequential EnKF & $0.34 \pm 0.06$ & $4.51 \pm 1.11$ & $0.095 \pm 0.108$ & $0.009 \pm 0.111$ \\
4D-EnKF         & $0.45 \pm 0.09$ & $4.42 \pm 1.15$ & $0.120 \pm 0.096$ & $0.219 \pm 0.099$ \\
QPCA-EnDCF      & $2.84 \pm 0.51$ & $3.55 \pm 0.70$ & $0.811 \pm 0.101$ & $0.820 \pm 0.085$ \\
\bottomrule
\end{tabular}
\end{table}

QPCA-EnDCF achieves $\bar{\gamma}=0.811\pm0.101$, close to ideal calibration with only mild residual underdispersion, and exhibits strong temporal tracking with $\rho=0.820\pm0.085$. These calibration gains accompany a $20\%$ reduction in RMSE relative to stochastic methods, indicating simultaneous improvement in point accuracy and uncertainty quantification. In contrast, the Sequential EnKF is severely underdispersed ($\bar{\gamma}=0.095\pm0.108$) with no meaningful temporal coupling ($\rho=0.009\pm0.111$), reflecting variance collapse under extreme undersampling ($N=10\ll n=40$). The 4D-EnKF yields only modest improvement ($\bar{\gamma}=0.120\pm0.096$, $\rho=0.219\pm0.099$), remaining far from calibrated despite an increase in effective degrees of freedom from $m=20$ to $mL=100$.

Temporal behavior is illustrated in Figure~\ref{fig:temporal_calibration}, which shows spread and RMSE across the $50$ assimilation windows.

\begin{figure}[htbp]
\centering
\includegraphics[width=\textwidth]{figures/fig04a_spread_vs_rmse_temporal.png}
\caption{Temporal evolution of ensemble spread (solid) and RMSE (dashed) over the 50-window sequence; shaded regions denote the 25th--75th percentile range (IQR band) across five trials. QPCA-EnDCF exhibits coherent covariation between spread and RMSE throughout. Sequential EnKF shows persistently low spread with no tracking of RMSE, while 4D-EnKF displays weak positive correlation but remains underdispersed.}
\label{fig:temporal_calibration}
\end{figure}

QPCA-EnDCF maintains coherent tracking, with spread and RMSE fluctuating together in the range $2$--$4$. Sequential EnKF exhibits the characteristic signature of ensemble collapse: spread remains below $0.5$ while RMSE varies between $3$ and $6$. The 4D-EnKF shows slightly larger spread ($0.4$--$0.6$) with weak coupling.

Figure~\ref{fig:reliability} provides complementary calibration diagnostics.

\begin{figure}[htbp]
\centering
\includegraphics[width=\textwidth]{figures/fig04_combined_calibration_analysis.png}
\caption{Probabilistic calibration diagnostics. \textbf{(A)} Spread--skill ratio $\gamma_w$ across assimilation windows; the dashed line indicates perfect calibration. \textbf{(B)} Reliability diagram plotting window-level ensemble spread $\sigma_w$ against window-level $\mathrm{RMSE}_w$ for all windows and trials, with the diagonal denoting $\sigma_w=\mathrm{RMSE}_w$.}
\label{fig:reliability}
\end{figure}

QPCA-EnDCF clusters near the identity line across RMSE values $2$--$5$, confirming accurate uncertainty quantification over varying error regimes. Sequential EnKF exhibits near-vertical clustering at $\sigma_{\mathrm{ens}}\approx0.3$, consistent with near-zero spread--skill correlation. The 4D-EnKF shows slightly larger spreads with weak positive association, but remains systematically above the diagonal (underdispersed).
The spread--skill analysis demonstrates that QPCA-EnDCF yields ensembles whose second moments are well calibrated in both magnitude and temporal dynamics, while stochastic methods exhibit order-of-magnitude underdispersion stemming from fundamental design limitations: observation perturbations introduce additional noise, while sampling error is distributed across all degrees of freedom rather than concentrated in dynamically dominant modes.
Spread--skill diagnostics assess second-moment calibration but do not test whether the full ensemble distribution is consistent with the truth. Rank histograms provide a distributional calibration test without Gaussian assumptions, based on exchangeability: for a calibrated ensemble, the true state is statistically indistinguishable from ensemble members, so all ranks occur with equal probability.
Figure~\ref{fig:rank_histograms} in Appendix~\ref{subsec:rank_histograms} summarizes the rank histogram analysis used to assess full distributional calibration. Overall, the rank histogram analysis demonstrates that QPCA-EnDCF substantially improves distributional calibration beyond variance matching alone. Its flatness metric, which is an order of magnitude smaller than those of the EnKF variants, indicates that the truth is placed appropriately across the ensemble support.
 While not perfectly uniform, deviations are minor relative to the severe systematic distortions observed in stochastic methods. The improvement is attributable to the deterministic update and spectral regularization of QPCA-EnDCF, which avoid observation perturbation noise and concentrate corrections into dynamically dominant modes, preserving ensemble diversity and yielding a more faithful approximation of the posterior distribution. See Appendix~\ref{subsec:rank_histograms} for further details and additional diagnostics.

% -----------------------------------------------------------------------------

% ================================================================================



\FloatBarrier


% ================================================================================
% SECTION 5: CONCENTRATION AND SPECTRAL PERTURBATION ANALYSIS
% ================================================================================
\section{Concentration and Spectral Perturbation Analysis for QPCA-EnDCF}
\label{sec:concentration_spectral}

This section presents the main convergence results for QPCA-EnDCF and provides the theoretical foundation for its uncertainty quantification properties. The analysis is organized around a decomposition that makes explicit how spectral truncation and finite-ensemble sampling jointly determine performance.

The development proceeds in three stages. First, we establish baseline identities and moment-based concentration estimates for the sample covariance of the whitened forecast--observation residuals, quantifying how empirical second-order statistics fluctuate around their population counterparts. Second, we invoke deterministic matrix perturbation theory to translate these covariance deviations into stability results for eigenvalues and eigenspaces, thereby controlling the accuracy of empirical spectral projectors. Third, we derive a bias--variance decomposition for the mean-squared error of the analysis mean, which cleanly separates truncation effects, which are governed by the truncation level $\kappa$ and the alignment of the mean innovation with the retained subspace from sampling effects, which depend on the ensemble size $N$ and on the stability of the empirical projector.

The EnDCF gain used in practice is an empirical quantity constructed from finite ensembles. In the regime where $N-1<d$, the observation-space sample covariance is singular, requiring a stabilized inverse (Definition~\ref{def:endcf_gain_ensemble}, Subsection~\ref{subsec:ensemble_uq}). Whenever expectations involve this random gain, we invoke the moment condition in Assumption~\ref{ass:gain_moment}.


% -----------------------------------------------------------------------------

% ================================================================================


\subsection{Regularity Assumptions}
\label{subsec:regularity_assumptions}

We collect here the regularity conditions required for the convergence analysis of QPCA-EnDCF. These assumptions are deliberately mild and are designed to isolate the essential probabilistic and spectral inputs needed for the bias--variance decomposition and the associated concentration estimates.

\begin{assumption}[Regularity conditions]
\label{ass:regularity_bias_variance}
Let $d:=mL$. We assume the following:
\begin{enumerate}[label=(\roman*)]
\item The forecast distribution $P_{\mathcal{X}}$ has a finite second moment, that is,
\[
\mathbb{E}\|\mathbf{x}\|^{2}<\infty .
\]

\item The stacked observation error covariance $\mathbf{R}^{(L)}\in\mathbb{R}^{d\times d}$ is symmetric positive definite and satisfies
\[
\lambda_{\min}\!\bigl(\mathbf{R}^{(L)}\bigr)\ge r_{\min}>0.
\]

\item The forecast ensemble $\{\mathbf{x}^{(j),f}\}_{j=1}^N$ consists of independent and identically distributed samples from $P_{\mathcal{X}}$. The corresponding whitened residuals $\{\mathbf{e}^{(j)}\}_{j=1}^N$, defined in Definition~\ref{def:whitened_residual_covariances}, satisfy
\[
\mathbb{E}\|\mathbf{e}^{(1)}\|^{4}<\infty,
\]
or, equivalently, $\mathbb{E}\|\mathbf{e}^{(1)}-\boldsymbol{\mu}_E\|^{4}<\infty$.

\item The population covariance $\boldsymbol{\Sigma}_E$ introduced in Definition~\ref{def:whitened_residual_covariances} admits the spectral decomposition
\[
\boldsymbol{\Sigma}_E=\sum_{i=1}^{d}\lambda_i\,\mathbf{v}_i\mathbf{v}_i^\top,
\qquad
\lambda_1\ge\lambda_2\ge\cdots\ge\lambda_d\ge 0.
\]
\end{enumerate}
\end{assumption}

\noindent
Assumption~\ref{ass:regularity_bias_variance}(i) ensures that all first and second-moment quantities appearing in the filter are well defined. Assumption~\ref{ass:regularity_bias_variance}(ii) guarantees that the whitening transformation is well posed. Assumption~\ref{ass:regularity_bias_variance}(iii) provides the probabilistic input for the covariance concentration estimates developed below; it is strictly weaker than Gaussianity and suffices for mean-square convergence results, while sharper fluctuation bounds may require stronger distributional assumptions. Finally, Assumption~\ref{ass:regularity_bias_variance}(iv) fixes notation for the population spectrum; no spectral gap or decay condition is imposed unless explicitly stated in later results.
The i.i.d.\ sampling assumption in (iii) is a mathematical idealization adopted for deriving the convergence bounds and concentration estimates in Sections~\ref{subsec:fundamental_concentration}--\ref{subsec:main_bv}. In practical cycling experiments (Sections~\ref{subsec:calibration}--\ref{sec:operational}), the analysis ensemble from window $w$ becomes the forecast ensemble for window $w+1$, inducing temporal dependence that violates strict i.i.d.\ conditions.
Nonetheless, the theoretical results provide qualitative guidance: under ergodic dynamics and sufficient mixing, consecutive windows become approximately independent in distribution, and the main conclusions, that sampling variance scales as $\mathcal{O}(1/N)$ and that the projector estimation contribution is governed by $\kappa$ and the cutoff gap $\delta_\kappa$ through stability bounds, remain empirically valid. In regimes where residual energy concentrates in low rank subspaces, the dominant sampling sensitive contribution exhibits rank driven behavior.
The experiments demonstrate that the predicted performance advantages persist under cycling, even though the formal assumptions are satisfied only asymptotically or approximately.
\par\smallskip

\begin{assumption}[Moment control for the empirical gain]
\label{ass:gain_moment}
Whenever expectations involve the random empirical gain, we assume that
\[
\mathbb{E}\!\left[\|\mathbf{K}^{\mathrm{DC}}\|_2^{2}\right]<\infty.
\]
This moment condition is mild in practice: when a stabilized inverse such as $(\mathbf{P}_{zz}+\varepsilon\mathbf{I})^{-1}$ is used, the inverse is uniformly bounded, and the condition follows from standard moment bounds on the ensemble anomalies (e.g., if the forecast distribution has exponential or sub-exponential tails, as in the Gaussian case). The pseudoinverse satisfies the condition under similar regularity conditions on the forecast distribution. This is the only property of the gain required to pass from conditional to unconditional mean-square bounds in the analysis.
\end{assumption}





% -----------------------------------------------------------------------------

% ================================================================================
\subsection{Fundamental Concentration Results}
\label{subsec:fundamental_concentration}

We begin by establishing the basic probabilistic identities and concentration estimates that form the backbone of the convergence analysis. The results in this subsection are deliberately elementary but essential: they separate deterministic bias from sampling variance and quantify how empirical covariance operators concentrate around their population counterparts. Together, these estimates provide the input required for subsequent spectral perturbation and bias--variance analyses.

\begin{lemma}[Bias--variance decomposition for the analysis mean]
\label{lem:mean_error_decomp}
Let
\(
\bar{\mathbf{x}}^{a}=\frac{1}{N}\sum_{j=1}^N\mathbf{x}^{(j),a}
\)
denote the analysis ensemble mean (at a fixed assimilation time; we suppress the time index). Then
\[
\bar{\mathbf{x}}^{a}-\mathbf{x}^{\mathrm{true}}
=
\bigl(\mathbb{E}[\bar{\mathbf{x}}^{a}]-\mathbf{x}^{\mathrm{true}}\bigr)
+
\bigl(\bar{\mathbf{x}}^{a}-\mathbb{E}[\bar{\mathbf{x}}^{a}]\bigr),
\]
and consequently
\[
\mathbb{E}\!\left[\|\bar{\mathbf{x}}^{a}-\mathbf{x}^{\mathrm{true}}\|^2\right]
=
\|\mathbb{E}[\bar{\mathbf{x}}^{a}]-\mathbf{x}^{\mathrm{true}}\|^2
+
\mathbb{E}\!\left[\|\bar{\mathbf{x}}^{a}-\mathbb{E}[\bar{\mathbf{x}}^{a}]\|^2\right].
\]
\end{lemma}

\begin{proof}
Define the deterministic bias vector
\(
\mathbf{b}:=\mathbb{E}[\bar{\mathbf{x}}^{a}]-\mathbf{x}^{\mathrm{true}}
\)
	and the zero-mean fluctuation
	\(
	\boldsymbol{\xi}:=\bar{\mathbf{x}}^{a}-\mathbb{E}[\bar{\mathbf{x}}^{a}],
	\)
	so that $\mathbb{E}[\boldsymbol{\xi}]=\mathbf{0}$. Writing
	\(
	\bar{\mathbf{x}}^{a}-\mathbf{x}^{\mathrm{true}}=\mathbf{b}+\boldsymbol{\xi},
	\)
	we compute
	\[
	\mathbb{E}\|\mathbf{b}+\boldsymbol{\xi}\|^2
	=
	\|\mathbf{b}\|^2
	+
	2\langle \mathbf{b},\mathbb{E}\boldsymbol{\xi}\rangle
	+
	\mathbb{E}\|\boldsymbol{\xi}\|^2
	=
	\|\mathbf{b}\|^2+\mathbb{E}\|\boldsymbol{\xi}\|^2,
	\]
	which yields the stated identity.
	\end{proof}

\begin{lemma}[Unbiasedness of the sample covariance]
\label{lem:unbiasedness_sample_cov}
Let $\boldsymbol{\Sigma}_E$ and $\mathbf{C}_E$ be as in Definition~\ref{def:whitened_residual_covariances}. Under Assumption~\ref{ass:regularity_bias_variance}(iii) and for all $N\ge 2$,
\[
\mathbb{E}[\mathbf{C}_E]=\boldsymbol{\Sigma}_E.
\]
\end{lemma}

\begin{proof}
Unbiasedness of the sample covariance with Bessel correction holds for independent samples with finite second moments. Let
\(
\tilde{\mathbf{e}}^{(j)}:=\mathbf{e}^{(j)}-\boldsymbol{\mu}_E
\)
and
\(
\bar{\tilde{\mathbf{e}}}:=\frac{1}{N}\sum_{j=1}^N\tilde{\mathbf{e}}^{(j)}.
\)
Then $\mathbf{e}^{(j)}-\bar{\mathbf{e}}=\tilde{\mathbf{e}}^{(j)}-\bar{\tilde{\mathbf{e}}}$ and
\[
\mathbf{C}_E
=
\frac{1}{N-1}\sum_{j=1}^N \tilde{\mathbf{e}}^{(j)}(\tilde{\mathbf{e}}^{(j)})^\top
-
\frac{N}{N-1}\bar{\tilde{\mathbf{e}}}\,\bar{\tilde{\mathbf{e}}}^\top.
\]
Taking expectations and using
\(
\mathbb{E}[\tilde{\mathbf{e}}^{(j)}(\tilde{\mathbf{e}}^{(j)})^\top]=\boldsymbol{\Sigma}_E
\)
together with
\(
\mathbb{E}[\bar{\tilde{\mathbf{e}}}\,\bar{\tilde{\mathbf{e}}}^\top]
=\mathrm{Cov}(\bar{\tilde{\mathbf{e}}})
=\boldsymbol{\Sigma}_E/N,
\)
we obtain $\mathbb{E}[\mathbf{C}_E]=\boldsymbol{\Sigma}_E$.
\end{proof}


\begin{lemma}[Mean-square Frobenius deviation of the sample covariance]
\label{lem:frobenius_concentration}
Let $\boldsymbol{\Sigma}_E$ and $\mathbf{C}_E$ be as in Definition~\ref{def:whitened_residual_covariances}. Under Assumption~\ref{ass:regularity_bias_variance}(iii), there exists a constant $C_{\mathrm{cov}}>0$, depending only on the dimension $d$ and on standardized fourth moments of $\mathbf{e}^{(1)}$, such that for all $N\ge 2$,
\[
\mathbb{E}\!\left[\|\mathbf{C}_E-\boldsymbol{\Sigma}_E\|_F^2\right]
\le
\frac{C_{\mathrm{cov}}}{N-1}\,
\mathrm{tr}\!\left(\boldsymbol{\Sigma}_E^2\right).
\]
In particular, $\mathbb{E}\|\mathbf{C}_E-\boldsymbol{\Sigma}_E\|_F^2=\mathcal{O}(N^{-1})$ as $N\to\infty$.
\end{lemma}

\begin{proof}
Let
\(
\tilde{\mathbf{e}}^{(j)}:=\mathbf{e}^{(j)}-\boldsymbol{\mu}_E
\)
and
\(
\bar{\tilde{\mathbf{e}}}:=\frac{1}{N}\sum_{j=1}^N\tilde{\mathbf{e}}^{(j)}.
\)
A direct expansion yields
\[
\mathbf{C}_E-\boldsymbol{\Sigma}_E
=
\frac{1}{N-1}\sum_{j=1}^N
\bigl(\tilde{\mathbf{e}}^{(j)}(\tilde{\mathbf{e}}^{(j)})^\top-\boldsymbol{\Sigma}_E\bigr)
-
\frac{N}{N-1}
\Bigl(\bar{\tilde{\mathbf{e}}}\,\bar{\tilde{\mathbf{e}}}^\top-\frac{1}{N}\boldsymbol{\Sigma}_E\Bigr).
\]
Applying the inequality $(a+b)^2\le 2a^2+2b^2$, using independence of the samples, and invoking the fourth-moment bound
\(
\mathbb{E}\|\bar{\tilde{\mathbf{e}}}\|^4\lesssim N^{-2}\mathbb{E}\|\tilde{\mathbf{e}}^{(1)}\|^4
\)
(which follows from Assumption~\ref{ass:regularity_bias_variance}(iii)), we obtain
\[
\mathbb{E}\!\left[\|\mathbf{C}_E-\boldsymbol{\Sigma}_E\|_F^2\right]
\le
\frac{C_{\mathrm{cov}}}{N-1}\,
\mathbb{E}\!\left[
\bigl\|
\tilde{\mathbf{e}}^{(1)}(\tilde{\mathbf{e}}^{(1)})^\top-\boldsymbol{\Sigma}_E
\bigr\|_F^2
\right],
\]
	for a constant $C_{\mathrm{cov}}>0$ depending only on $d$ and standardized fourth moments. Finally, since
	\(
	\|\tilde{\mathbf{e}}^{(1)}(\tilde{\mathbf{e}}^{(1)})^\top-\boldsymbol{\Sigma}_E\|_F^2
	\le
	2\|\tilde{\mathbf{e}}^{(1)}(\tilde{\mathbf{e}}^{(1)})^\top\|_F^2+2\|\boldsymbol{\Sigma}_E\|_F^2
	=
		2\|\tilde{\mathbf{e}}^{(1)}\|^4+2\,\mathrm{tr}(\boldsymbol{\Sigma}_E^2),
		\)
		and the standardized fourth-moment control implicit in the constant $C_{\mathrm{cov}}$ ensures that $\mathbb{E}\|\tilde{\mathbf{e}}^{(1)}\|^4$ is bounded by a dimension-dependent multiple of $\mathrm{tr}(\boldsymbol{\Sigma}_E^2)$. In particular, if $\tilde{\mathbf{e}}^{(1)}$ is Gaussian then
		\[
		\mathbb{E}\|\tilde{\mathbf{e}}^{(1)}\|^4
		=
		\bigl(\mathrm{tr}(\boldsymbol{\Sigma}_E)\bigr)^2 + 2\,\mathrm{tr}(\boldsymbol{\Sigma}_E^2)
		\le
		(d+2)\,\mathrm{tr}(\boldsymbol{\Sigma}_E^2),
		\]
		where the final inequality uses $\bigl(\mathrm{tr}(\boldsymbol{\Sigma}_E)\bigr)^2\le d\,\mathrm{tr}(\boldsymbol{\Sigma}_E^2)$. Absorbing this ratio into $C_{\mathrm{cov}}$ yields the stated bound.
\end{proof}

If, in addition, $\tilde{\mathbf{e}}^{(1)}$ is Gaussian, then $(N-1)\mathbf{C}_E$ is Wishart and
\[
\mathbb{E}\|\mathbf{C}_E-\boldsymbol{\Sigma}_E\|_F^2
=
\frac{\bigl(\mathrm{tr}(\boldsymbol{\Sigma}_E)\bigr)^2+\mathrm{tr}(\boldsymbol{\Sigma}_E^2)}{N-1}.
\]
In particular, $\mathbb{E}\|\mathbf{C}_E-\boldsymbol{\Sigma}_E\|_F^2\le \frac{d+1}{N-1}\mathrm{tr}(\boldsymbol{\Sigma}_E^2)$.

\begin{lemma}[Fourth-moment bound for the sample mean]
\label{lem:sample_mean_fourth_moment}
Let $\{\mathbf{e}^{(j)}\}_{j=1}^N$ be i.i.d.\ in $\mathbb{R}^d$ with mean $\boldsymbol{\mu}_E$ and finite fourth moment $\mathbb{E}\|\mathbf{e}^{(1)}-\boldsymbol{\mu}_E\|^4<\infty$. Define $\bar{\mathbf{e}}:=\frac{1}{N}\sum_{j=1}^N\mathbf{e}^{(j)}$. Then there exists a constant $C_{\mathrm{mean}}>0$, depending only on $d$, such that for all $N\ge 2$,
\[
\mathbb{E}\|\bar{\mathbf{e}}-\boldsymbol{\mu}_E\|^4
\le
\frac{C_{\mathrm{mean}}}{N^2}\left(\mathbb{E}\|\mathbf{e}^{(1)}-\boldsymbol{\mu}_E\|^4+\bigl(\mathbb{E}\|\mathbf{e}^{(1)}-\boldsymbol{\mu}_E\|^2\bigr)^2\right).
\]
In particular, $\bigl(\mathbb{E}\|\bar{\mathbf{e}}-\boldsymbol{\mu}_E\|^4\bigr)^{1/2}=\mathcal{O}(N^{-1})$ as $N\to\infty$.
\end{lemma}

\begin{proof}
Let $\mathbf{X}^{(j)}:=\mathbf{e}^{(j)}-\boldsymbol{\mu}_E$ and $\bar{\mathbf{X}}:=\frac{1}{N}\sum_{j=1}^N\mathbf{X}^{(j)}=\bar{\mathbf{e}}-\boldsymbol{\mu}_E$. For each coordinate $i\in\{1,\ldots,d\}$, the scalar random variables $\{X_i^{(j)}\}_{j=1}^N$ are i.i.d.\ with mean $0$, and the fourth-moment identity for sums yields
\[
\mathbb{E}\bigl[(\bar X_i)^4\bigr]
=
\frac{1}{N^4}\left(N\,\mathbb{E}\bigl[(X_i^{(1)})^4\bigr] + 3N(N-1)\bigl(\mathbb{E}[(X_i^{(1)})^2]\bigr)^2\right)
\le
\frac{\mathbb{E}\bigl[(X_i^{(1)})^4\bigr]}{N^3}
\;+\;
\frac{3\bigl(\mathbb{E}[(X_i^{(1)})^2]\bigr)^2}{N^2}.
\]
Next, using $(\sum_{i=1}^d a_i)^2\le d\sum_{i=1}^d a_i^2$ with $a_i=(\bar X_i)^2$ gives
\[
\|\bar{\mathbf{X}}\|^4
=
\left(\sum_{i=1}^d (\bar X_i)^2\right)^2
\le
d\sum_{i=1}^d (\bar X_i)^4.
\]
Taking expectations and summing the coordinate bounds yields
\[
\mathbb{E}\|\bar{\mathbf{X}}\|^4
\le
\frac{d}{N^3}\sum_{i=1}^d \mathbb{E}\bigl[(X_i^{(1)})^4\bigr]
\;+\;
\frac{3d}{N^2}\sum_{i=1}^d \bigl(\mathbb{E}[(X_i^{(1)})^2]\bigr)^2.
\]
Finally, $\sum_i \mathbb{E}[(X_i^{(1)})^4]\le \mathbb{E}\|\mathbf{X}^{(1)}\|^4$ and
$\sum_i (\mathbb{E}[(X_i^{(1)})^2])^2\le (\sum_i \mathbb{E}[(X_i^{(1)})^2])^2=(\mathbb{E}\|\mathbf{X}^{(1)}\|^2)^2$.
Absorbing the factor $d$ into $C_{\mathrm{mean}}$ gives the stated bound.
\end{proof}



% -----------------------------------------------------------------------------

% ================================================================================

\subsection{Spectral Perturbation Bounds}
\label{subsec:spectral_perturbation}

We next convert matrix-level covariance deviations into spectral stability estimates. This step is essential for QPCA-EnDCF, whose update is governed by the rank-$\kappa$ empirical projector $\hat{\mathbf{P}}_\kappa$. The results below are deterministic perturbation inequalities: randomness enters only through the perturbation $\mathbf{C}_E-\boldsymbol{\Sigma}_E$, which will later be controlled in expectation via the concentration estimates of Section~\ref{subsec:fundamental_concentration}.

\begin{lemma}[Spectral perturbation and projector stability]
\label{lem:covariance_concentration}
Let $\boldsymbol{\Sigma}_E$, $\mathbf{C}_E$, and their eigendecompositions be as in Definition~\ref{def:spectral_decompositions}. Then:
\begin{enumerate}[label=(\roman*)]
\item For each $i=1,\ldots,d$,
\[
|\hat{\lambda}_i-\lambda_i|
\le
\|\mathbf{C}_E-\boldsymbol{\Sigma}_E\|_2 .
\]

	\item Fix $i\in\{1,\ldots,d\}$ and assume that $\lambda_i$ is isolated, i.e.,
	\[
	\mathrm{gap}_i
	:=
\min\{\lambda_{i-1}-\lambda_i,\ \lambda_i-\lambda_{i+1}\}
\ge
\delta_i>0,
\qquad
(\lambda_0:=+\infty,\ \lambda_{d+1}:=0).
\]
	Choose the sign of $\hat{\mathbf{v}}_i$ so that $\langle \hat{\mathbf{v}}_i,\mathbf{v}_i\rangle\ge 0$. Then
	\[
	\sin\angle(\hat{\mathbf{v}}_i,\mathbf{v}_i)
	=
	\|(\mathbf{I}-\mathbf{v}_i\mathbf{v}_i^\top)\hat{\mathbf{v}}_i\|_2
	\le
	\frac{\|\mathbf{C}_E-\boldsymbol{\Sigma}_E\|_2}{\delta_i},
	\qquad
	\|\hat{\mathbf{v}}_i-\mathbf{v}_i\|_2
	\le
	\frac{\sqrt{2}\,\|\mathbf{C}_E-\boldsymbol{\Sigma}_E\|_2}{\delta_i}.
	\]

\item Fix $\kappa\in\{1,\ldots,d\}$ and assume a cutoff gap
\[
\lambda_\kappa-\lambda_{\kappa+1}\ge \delta_\kappa>0,
\qquad (\lambda_{d+1}:=0).
\]
Let
\(
\mathbf{P}_\kappa:=\sum_{i=1}^{\kappa}\mathbf{v}_i\mathbf{v}_i^\top
\)
and
\(
\hat{\mathbf{P}}_\kappa:=\sum_{i=1}^{\kappa}\hat{\mathbf{v}}_i\hat{\mathbf{v}}_i^\top.
\)
	Then
	\[
	\|\hat{\mathbf{P}}_\kappa-\mathbf{P}_\kappa\|_F
	\le
	\frac{\sqrt{2\kappa}}{\delta_\kappa}\,
	\|\mathbf{C}_E-\boldsymbol{\Sigma}_E\|_2
	\le
	\frac{\sqrt{2\kappa}}{\delta_\kappa}\,
	\|\mathbf{C}_E-\boldsymbol{\Sigma}_E\|_F.
	\]


\item With $\boldsymbol{\Sigma}_E^\kappa:=\mathbf{P}_\kappa\boldsymbol{\Sigma}_E\mathbf{P}_\kappa$
and $\mathbf{C}_E^\kappa:=\hat{\mathbf{P}}_\kappa\mathbf{C}_E\hat{\mathbf{P}}_\kappa$, one has
\[
\|\mathbf{C}_E^\kappa-\boldsymbol{\Sigma}_E^\kappa\|_F
\;\le\;
\|\mathbf{P}_\kappa(\mathbf{C}_E-\boldsymbol{\Sigma}_E)\mathbf{P}_\kappa\|_F
+
\|(\hat{\mathbf{P}}_\kappa-\mathbf{P}_\kappa)\mathbf{C}_E\|_F
+
\|\mathbf{C}_E(\hat{\mathbf{P}}_\kappa-\mathbf{P}_\kappa)\|_F .
\]
\end{enumerate}
\end{lemma}

\begin{proof}
The statements follow from standard deterministic perturbation results for symmetric matrices.

\smallskip
\noindent Weyl's inequality yields
\[
|\lambda_i(\mathbf{C}_E)-\lambda_i(\boldsymbol{\Sigma}_E)|
\le
\|\mathbf{C}_E-\boldsymbol{\Sigma}_E\|_2,
\qquad i=1,\ldots,d,
\]
proving (i).

\smallskip
\noindent Fix $i$ and assume the separation condition in (ii). The Davis--Kahan $\sin\Theta$ theorem gives
\[
\sin\angle(\hat{\mathbf{v}}_i,\mathbf{v}_i)
=
\|(\mathbf{I}-\mathbf{v}_i\mathbf{v}_i^\top)\hat{\mathbf{v}}_i\|_2
\le
\frac{\|\mathbf{C}_E-\boldsymbol{\Sigma}_E\|_2}{\delta_i}.
\]
Choosing the sign of $\hat{\mathbf{v}}_i$ so that $\langle \hat{\mathbf{v}}_i,\mathbf{v}_i\rangle\ge 0$ and using
$\|\hat{\mathbf{v}}_i-\mathbf{v}_i\|_2=2\sin(\angle(\hat{\mathbf{v}}_i,\mathbf{v}_i)/2)\le \sqrt{2}\,\sin\angle(\hat{\mathbf{v}}_i,\mathbf{v}_i)$ yields the Euclidean bound in (ii).

\smallskip
\noindent Under the cutoff-gap condition in (iii), Davis--Kahan in projector form yields
\[
\|\hat{\mathbf{P}}_\kappa-\mathbf{P}_\kappa\|_F
\le
\frac{\sqrt{2\kappa}}{\delta_\kappa}\,
\|\mathbf{C}_E-\boldsymbol{\Sigma}_E\|_2,
\]
and the final inequality follows from $\|\cdot\|_2\le \|\cdot\|_F$.

\smallskip
\noindent Using $\boldsymbol{\Sigma}_E^\kappa=\mathbf{P}_\kappa\boldsymbol{\Sigma}_E\mathbf{P}_\kappa$ and
$\mathbf{C}_E^\kappa=\hat{\mathbf{P}}_\kappa\mathbf{C}_E\hat{\mathbf{P}}_\kappa$, insert and subtract
$\mathbf{P}_\kappa\mathbf{C}_E\mathbf{P}_\kappa$ to obtain
\[
\mathbf{C}_E^\kappa-\boldsymbol{\Sigma}_E^\kappa
=
\mathbf{P}_\kappa(\mathbf{C}_E-\boldsymbol{\Sigma}_E)\mathbf{P}_\kappa
+
(\hat{\mathbf{P}}_\kappa-\mathbf{P}_\kappa)\mathbf{C}_E\hat{\mathbf{P}}_\kappa
+
\mathbf{P}_\kappa\mathbf{C}_E(\hat{\mathbf{P}}_\kappa-\mathbf{P}_\kappa).
\]
Taking Frobenius norms and applying the triangle inequality gives (iv).
\end{proof}

\begin{lemma}[Fourth-moment bound for the sample covariance error (Gaussian case)]
\label{lem:cov_frob_fourth_moment}
Let $\boldsymbol{\Sigma}_E$ and $\mathbf{C}_E$ be as in Definition~\ref{def:whitened_residual_covariances}. Assume that the centered whitened residuals
\[
\tilde{\mathbf{e}}^{(j)}:=\mathbf{e}^{(j)}-\boldsymbol{\mu}_E
\]
are i.i.d.\ Gaussian,
\[
\tilde{\mathbf{e}}^{(j)} \stackrel{\mathrm{i.i.d.}}{\sim} \mathcal{N}(\mathbf{0},\boldsymbol{\Sigma}_E)
\qquad\text{(equivalently, }\mathbf{e}^{(j)} \stackrel{\mathrm{i.i.d.}}{\sim} \mathcal{N}(\boldsymbol{\mu}_E,\boldsymbol{\Sigma}_E)\text{)}.
\]
Then there exists a constant $C_{\mathrm{W}}>0$, depending only on the observation-space dimension $d=mL$, such that for all $N\ge 2$,
\begin{equation}
\label{eq:cov_fourth_moment_bound}
\mathbb{E}\!\left[\|\mathbf{C}_E-\boldsymbol{\Sigma}_E\|_F^4\right]
\le
\frac{C_{\mathrm{W}}}{(N-1)^2}\,\bigl(\mathrm{tr}(\boldsymbol{\Sigma}_E^2)\bigr)^2.
\end{equation}
In particular,
\[
\Bigl(\mathbb{E}\|\mathbf{C}_E-\boldsymbol{\Sigma}_E\|_F^4\Bigr)^{1/2}
=
\mathcal{O}\!\left(\frac{1}{N}\right)
\quad\text{as }N\to\infty.
\]
\end{lemma}

\begin{proof}
Let $\tilde{\mathbf{e}}^{(j)}:=\mathbf{e}^{(j)}-\boldsymbol{\mu}_E$, so that
$\tilde{\mathbf{e}}^{(j)}\stackrel{\mathrm{i.i.d.}}{\sim}\mathcal{N}(\mathbf{0},\boldsymbol{\Sigma}_E)$.
Write the sample covariance as
\[
\mathbf{C}_E
=
\frac{1}{N-1}\sum_{j=1}^N(\tilde{\mathbf{e}}^{(j)}-\bar{\tilde{\mathbf{e}}})(\tilde{\mathbf{e}}^{(j)}-\bar{\tilde{\mathbf{e}}})^\top,
\qquad
\bar{\tilde{\mathbf{e}}}:=\frac{1}{N}\sum_{j=1}^N \tilde{\mathbf{e}}^{(j)}.
\]
Let $\mathbf{A}\in\mathbb{R}^{d\times N}$ be the data matrix with columns $\tilde{\mathbf{e}}^{(j)}$, and let
$\mathbf{J}_N:=\mathbf{I}_N-\frac{1}{N}\mathbf{1}\mathbf{1}^\top$ be the centering matrix. Then
\[
(N-1)\mathbf{C}_E=\mathbf{A}\mathbf{J}_N\mathbf{A}^\top.
\]

\smallskip
\noindent Since $\mathbf{J}_N$ is an orthogonal projector of rank $N-1$, there exists an orthogonal matrix
$\mathbf{U}\in\mathbb{R}^{N\times N}$ such that
\[
\mathbf{J}_N
=
\mathbf{U}
\begin{bmatrix}
\mathbf{I}_{N-1} & 0\\
0 & 0
\end{bmatrix}
\mathbf{U}^\top.
\]
Let $\mathbf{U}_{1:(N-1)}\in\mathbb{R}^{N\times (N-1)}$ denote the first $N-1$ columns of $\mathbf{U}$. Then
\[
(N-1)\mathbf{C}_E
=
\mathbf{A}\mathbf{U}_{1:(N-1)}
\bigl(\mathbf{A}\mathbf{U}_{1:(N-1)}\bigr)^\top.
\]

The key point is that right multiplication by an orthogonal matrix preserves the distribution of a Gaussian data matrix. Indeed, since the columns of $\mathbf{A}$ are i.i.d.\ $\mathcal{N}(\mathbf{0},\boldsymbol{\Sigma}_E)$, for any orthogonal $\mathbf{U}$,
\[
\mathrm{vec}(\mathbf{A}\mathbf{U})
=
(\mathbf{U}^\top\otimes \mathbf{I}_d)\,\mathrm{vec}(\mathbf{A})
\sim
\mathcal{N}\!\left(\mathbf{0},\,\mathbf{I}_N\otimes \boldsymbol{\Sigma}_E\right),
\]
so $\mathbf{A}\mathbf{U}$ and $\mathbf{A}$ are identically distributed. In particular, the columns of $\mathbf{A}\mathbf{U}$ are again i.i.d.\ $\mathcal{N}(\mathbf{0},\boldsymbol{\Sigma}_E)$, and thus the first $N-1$ columns
\[
\mathbf{B}:=\mathbf{A}\mathbf{U}_{1:(N-1)}\in\mathbb{R}^{d\times (N-1)}
\]
are i.i.d.\ $\mathcal{N}(\mathbf{0},\boldsymbol{\Sigma}_E)$. Therefore
\[
(N-1)\mathbf{C}_E=\mathbf{B}\mathbf{B}^\top
\]
is a central Wishart random matrix with $N-1$ degrees of freedom and scale $\boldsymbol{\Sigma}_E$.

\smallskip
\noindent Standard fourth-moment bounds for the Wishart distribution imply that there exists a constant $C_{\mathrm{W}}>0$, depending only on $d$, such that
\[
\mathbb{E}\!\left[\|\mathbf{C}_E-\boldsymbol{\Sigma}_E\|_F^4\right]
\le
\frac{C_{\mathrm{W}}}{(N-1)^2}\,\bigl(\mathrm{tr}(\boldsymbol{\Sigma}_E^2)\bigr)^2,
\]
which is \eqref{eq:cov_fourth_moment_bound}. Taking square roots yields the stated $\mathcal{O}(N^{-1})$ rate.
\end{proof}

Spectral separation is intrinsic to eigenspace stability: when eigenvalues are repeated or nearly repeated, arbitrarily small perturbations may induce large rotations within the associated invariant subspace even though eigenvalues remain stable. For QPCA-EnDCF, stability of the rank-$\kappa$ projector is therefore strongest when the cutoff gap $\lambda_\kappa-\lambda_{\kappa+1}$ is well separated.



% -----------------------------------------------------------------------------

% ================================================================================
\section{Bias--Variance Decomposition and Main Theorem}
\label{sec:bv_decomposition}


% -----------------------------------------------------------------------------

% ================================================================================

\subsection{Measure Spaces and Probability Measures}
\label{subsec:measure_spaces}

\begin{definition}[State space and forecast measure]
\label{def:parameter_space}
Let $\mathcal{X}=\mathbb{R}^n$ denote the state space, equipped with its Borel $\sigma$-algebra $\mathcal{B}_{\mathcal{X}}$ and Lebesgue measure $\mu_{\mathcal{X}}$. We assume the existence of a forecast (prior) density $\pi_{\mathcal{X}}^{f}:\mathcal{X}\to\mathbb{R}_{+}$ that is absolutely continuous with respect to $\mu_{\mathcal{X}}$ and satisfies
\[
\int_{\mathcal{X}}\pi_{\mathcal{X}}^{f}(\mathbf{x})\,\mathrm{d}\mu_{\mathcal{X}}(\mathbf{x})=1,
\qquad
\int_{\mathcal{X}}\|\mathbf{x}\|^{2}\,\pi_{\mathcal{X}}^{f}(\mathbf{x})\,\mathrm{d}\mu_{\mathcal{X}}(\mathbf{x})<\infty.
\]
The associated probability measure $P_{\mathcal{X}}$ is defined by
\[
P_{\mathcal{X}}(A)
=
\int_A \pi_{\mathcal{X}}^{f}(\mathbf{x})\,\mathrm{d}\mu_{\mathcal{X}}(\mathbf{x}),
\qquad
\forall A\in\mathcal{B}_{\mathcal{X}}.
\]
A forecast ensemble $\{\mathbf{x}^{(j),f}\}_{j=1}^N$ is assumed to consist of independent and identically distributed draws from $P_{\mathcal{X}}$.
\end{definition}


In practice, $P_{\mathcal{X}}$ represents the algorithmic forecast distribution (e.g., the Gaussian approximation constructed by the assimilation method) rather than the true invariant measure of the Lorenz--96 attractor, which is singular with respect to Lebesgue measure. The i.i.d.\ sampling assumption is an idealization: in sequential filtering, ensemble members become statistically dependent through repeated cycling. The assumption is invoked here to enable tractable analysis of sampling error and is understood as an approximation valid when ensemble diversity is actively maintained (e.g., through inflation or spectral regularization).
\par\smallskip

\begin{assumption}[Gaussian forecast model]
\label{ass:gaussian_forecast}
When explicit covariance and higher-moment identities are required, we specialize to the Gaussian case and assume
\[
P_{\mathcal{X}}=\mathcal{N}(\boldsymbol{\mu}_0^{f},\mathbf{P}^{f}),
\]
so that $\mathbb{E}[\mathbf{x}_0^{(j),f}]=\boldsymbol{\mu}_0^{f}$ and $\mathrm{Cov}(\mathbf{x}_0^{(j),f})=\mathbf{P}^{f}$.
This assumption is invoked only when necessary and does not otherwise restrict the general formulation.
\end{assumption}

\begin{definition}[Observation space and quantity-of-interest map]
\label{def:data_space}
Let $\mathcal{D}=\mathbb{R}^{d}$ with $d:=mL$ denote the stacked observation space, equipped with its Borel $\sigma$-algebra $\mathcal{B}_{\mathcal{D}}$ and Lebesgue measure $\mu_{\mathcal{D}}$. The stacked observation vector $\mathbf{z}^{(w)}$ defined in \eqref{eq:stacked_obs} is an element of $\mathcal{D}$. A quantity-of-interest (QoI) map is any measurable function $Q:\mathcal{X}\to\mathcal{D}$. Throughout the theoretical analysis, we work with a tangent-linear approximation of $Q$, represented by a stacked observation operator $\mathbf{H}^{(L)}\in\mathbb{R}^{d\times n}$:
\[
Q(\mathbf{x}_0)\approx\mathbf{H}^{(L)}\mathbf{x}_0
=
\begin{bmatrix}
\mathbf{H}\mathbf{M}_1\\
\mathbf{H}\mathbf{M}_2\\
\vdots\\
\mathbf{H}\mathbf{M}_L
\end{bmatrix}
\mathbf{x}_0
\in\mathbb{R}^{d},
\]
where $\mathbf{M}_\ell\in\mathbb{R}^{n\times n}$ is the tangent linear model (TLM) obtained by linearizing the nonlinear dynamics $\mathcal{M}$ about a reference trajectory $\{\bar{\mathbf{x}}_\ell\}_{\ell=0}^{L}$ (typically the ensemble mean trajectory), so that $\mathbf{M}_\ell := \frac{\partial}{\partial\mathbf{x}_0}\mathcal{M}^{(\ell)}(\mathbf{x}_0)\big|_{\mathbf{x}_0=\bar{\mathbf{x}}_0}$ maps perturbations at the window-initial time to perturbations at time $\ell$. 
\end{definition}

This tangent-linear formulation permits closed-form expressions for population covariance propagation and underpins the spectral analysis in Sections~\ref{subsec:practical_considerations}--\ref{sec:concentration_spectral}. The approximation is exact for linear dynamics and linear observation operators; for nonlinear dynamics, it is valid to first order in perturbations about the reference trajectory. The numerical experiments (Section~\ref{sec:motivation}) propagate the full nonlinear model $\mathcal{M}$ and then apply the linear observation operator $\mathbf{H}$; the theoretical results thus characterize algorithm behavior in the tangent-linear regime, which provides a useful approximation when ensemble spread is small relative to the scale of nonlinearity.
\begin{definition}[Convergence of approximate QoI maps]
\label{def:map_convergence}
A sequence of measurable maps $\{Q_s\}_{s\ge 1}$, with $Q_s:\mathcal{X}\to\mathcal{D}$, is said to converge to $Q$ in $L^2(P_{\mathcal{X}})$ if
\begin{align*}
\|Q_s-Q\|_{L^2(P_{\mathcal{X}})}^2
&:=
\mathbb{E}_{P_{\mathcal{X}}}\!\left[\|Q_s(\mathbf{x})-Q(\mathbf{x})\|^2\right]\\
&=
\int_{\mathcal{X}}\|Q_s(\mathbf{x})-Q(\mathbf{x})\|^2\,
\pi_{\mathcal{X}}^{f}(\mathbf{x})\,\mathrm{d}\mu_{\mathcal{X}}(\mathbf{x})
\longrightarrow 0
\quad\text{as } s\to\infty.
\end{align*}
\end{definition}
This notion of convergence will be used to formalize consistency of approximate observation operators and surrogate models. The index $s$ (for ``surrogate'' or ``approximation stage'') is distinct from the within-window time index $\ell\in\{1,\ldots,L\}$ used elsewhere in this paper.
The state space $\mathcal{X}$ and observation space $\mathcal{D}$ correspond to the parameter space $\Lambda$ and data space $D$, respectively, in~\cite{ButlerWildeyZhang2020}. Likewise, the forecast density $\pi_{\mathcal{X}}^{f}$ corresponds to $\pi_{\Lambda}$ in that notation, and the maps $Q_s:\mathcal{X}\to\mathcal{D}$ play the role of the parameter-to-QoI maps considered therein.




% -----------------------------------------------------------------------------

% ================================================================================

For approximate QoI maps $Q_s$ (where $s$ denotes the approximation stage, distinct from the within-window time index $\ell$), define the corresponding whitened residuals and their population moments by
\[
\mathbf{e}^{(j),s}:=(\mathbf{R}^{(L)})^{-1/2}\bigl(Q_s(\mathbf{x}^{(j),f})-\mathbf{z}^{(w)}\bigr),
\qquad
\boldsymbol{\mu}_E^{\,s}:=\mathbb{E}[\mathbf{e}^{(j),s}],
\qquad
\boldsymbol{\Sigma}_E^{\,s}:=\mathrm{Cov}(\mathbf{e}^{(j),s}).
\]
When $Q_s=Q$, these reduce to $\mathbf{e}^{(j)}$, $\boldsymbol{\mu}_E$, and $\boldsymbol{\Sigma}_E$.
For each fixed $s$, we also define the sample mean and sample covariance of the $s$-residuals by
\[
\bar{\mathbf{e}}^{\,s}:=\frac{1}{N}\sum_{j=1}^N \mathbf{e}^{(j),s},
\qquad
\mathbf{C}_E^{\,s}:=\frac{1}{N-1}\sum_{j=1}^N(\mathbf{e}^{(j),s}-\bar{\mathbf{e}}^{\,s})(\mathbf{e}^{(j),s}-\bar{\mathbf{e}}^{\,s})^\top.
\]
Let $\mathbf{P}_\kappa^{\,s}:=\mathbf{P}_\kappa(\boldsymbol{\Sigma}_E^{\,s})$ and $\hat{\mathbf{P}}_\kappa^{\,s}:=\mathbf{P}_\kappa(\mathbf{C}_E^{\,s})$ denote the corresponding population and empirical rank-$\kappa$ projectors (Definition~\ref{def:projectors_truncation}).

\subsection{Main theorem: bias--variance decomposition}
\label{subsec:main_bv}

We now combine the elementary bias--variance identity of Lemma~\ref{lem:mean_error_decomp} with the deterministic projector stability bounds of Lemma~\ref{lem:covariance_concentration} to obtain a decomposition and explicit estimates for the mean-squared error of the QPCA-EnDCF analysis mean. The resulting bounds isolate three effects: (i) truncation bias induced by restricting the correction to a rank-$\kappa$ subspace, (ii) approximation error arising from using an approximate QoI map $Q_s$, and (iii) sampling error associated with estimating the projector from a finite ensemble.

For an estimator $\bar{\mathbf{x}}^a$ of the true state $\mathbf{x}^{\mathrm{true}}$ at a fixed assimilation time, the MSE decomposes as
\begin{equation}
\mathrm{MSE}
=
\mathbb{E}\!\left[\|\bar{\mathbf{x}}^a-\mathbf{x}^{\mathrm{true}}\|^2\right]
=
\mathrm{Bias}^2+\mathrm{Variance},
\label{eq:mse_decomposition_empirical}
\end{equation}
where the expectation is taken over independent ensemble realizations (different initial ensemble draws) with the true trajectory and observations held fixed. The bias term
$\mathrm{Bias}^2=\|\mathbb{E}[\bar{\mathbf{x}}^a]-\mathbf{x}^{\mathrm{true}}\|^2$
captures systematic error, while
$\mathrm{Variance}=\mathbb{E}[\|\bar{\mathbf{x}}^a-\mathbb{E}[\bar{\mathbf{x}}^a]\|^2]$
quantifies sampling-induced fluctuations. Bias reflects deterministic algorithmic limitations; variance is stochastic and constrained by ensemble size.

In the sequential setting where truth and observations evolve over many assimilation cycles, the decomposition \eqref{eq:mse_decomposition_empirical} applies at each assimilation time (here, at each window endpoint). We therefore estimate bias and variance as functions of window index and then time-average over windows to obtain representative values for a method.
\begin{definition}[Bias Variance Hypotheses]
\label{def:bias_variance_hypotheses}
Consider QPCA-EnDCF with truncation level $\kappa$ and empirical gain
$\mathbf{K}^{\mathrm{DC}}=\mathbf{P}_{xz}\mathbf{P}_{zz}^{\dagger}$ as in Definition~\ref{def:endcf_gain_ensemble}. Let
\(
\bar{\mathbf{x}}^{a,s}=\frac{1}{N}\sum_{j=1}^N\mathbf{x}^{(j),a,s}
\)
denote the analysis ensemble mean (where the superscript $s$ indexes the approximation stage). Assume Assumptions~\ref{ass:regularity_bias_variance} and~\ref{ass:gain_moment}. Suppose that the population covariance $\boldsymbol{\Sigma}_E^{\,s}$ of the $s$-residuals has a cutoff gap
\begin{equation}
\label{eq:gap_cutoff_bv}
\lambda_\kappa^{\,s}-\lambda_{\kappa+1}^{\,s}\ge \delta_\kappa>0,
\qquad (\lambda_{d+1}:=0,\ d=mL),
\end{equation}
where $\{\lambda_i^{\,s}\}_{i=1}^d$ are the eigenvalues of $\boldsymbol{\Sigma}_E^{\,s}$ in decreasing order.
Then the mean-squared error admits the exact decomposition
\begin{equation}
\mathbb{E}\!\left[\|\bar{\mathbf{x}}^{a,s}-\mathbf{x}^{\mathrm{true}}\|^2\right]
=
\mathrm{Bias}^2(\kappa,s)+\mathrm{Var}(N,\kappa,s),
\end{equation}
where
\begin{equation}
\mathrm{Bias}^2(\kappa,s):=\bigl\|\mathbb{E}[\bar{\mathbf{x}}^{a,s}]-\mathbf{x}^{\mathrm{true}}\bigr\|^2,
\qquad
\mathrm{Var}(N,\kappa,s):=\mathbb{E}\!\left[\|\bar{\mathbf{x}}^{a,s}-\mathbb{E}[\bar{\mathbf{x}}^{a,s}]\|^2\right].
\end{equation}

\end{definition}

\begin{theorem}[Bias Decomposition with Spectral Truncation]
\label{thm:bias_variance}
	Under the Bias Variance Hypotheses~\ref{def:bias_variance_hypotheses}, the bias satisfies
	\[
		\mathrm{Bias}_{\mathrm{base}}^{2}(s)
		:=
		\bigl\|\boldsymbol{\mu}_L^f-\mathbf{x}^{\mathrm{true}}-\mathbb{E}[\mathbf{K}^{\mathrm{DC}}](\mathbf{R}^{(L)})^{1/2}\boldsymbol{\mu}_E^{\,s}\bigr\|^2,
		\]
	\begin{align}
	\mathrm{Bias}^2(\kappa,s)
	&\le
	2\,\mathrm{Bias}_{\mathrm{base}}^{2}(s)
	+
	4\,\mathbb{E}\!\left[\|\mathbf{K}^{\mathrm{DC}}\|_2^2\right]\,
	\|\mathbf{R}^{(L)}\|_2\,
	\bigl\|(\mathbf{I}-\mathbf{P}_\kappa^{\,s})\boldsymbol{\mu}_E^{\,s}\bigr\|^2
	\;+\;
	2\,C_{\mathrm{m}}\,
	\|Q_s-Q\|_{L^2(P_{\mathcal{X}})}^2
	\nonumber\\
		&\hspace{3.2em}
		+\;
		\frac{2\|\mathbf{R}^{(L)}\|_2}{N}\,
		\mathbb{E}\!\left[\|\mathbf{K}^{\mathrm{DC}}\|_2^2\right]\,
		\mathrm{tr}\!\bigl(\boldsymbol{\Sigma}_E^{\,s}\bigr)
		\nonumber\\
		&\hspace{3.2em}
		+\;
		2\,C_{\mathrm{p}}\,
		\frac{\kappa}{\delta_\kappa^2}\,
		\mathbb{E}\!\left[\|\mathbf{C}_E^{\,s}-\boldsymbol{\Sigma}_E^{\,s}\|_F^2\right],
		\label{eq:bias_bound_fixed_main}
		\end{align}
for constants $C_{\mathrm{m}},C_{\mathrm{p}}>0$ independent of $N$.
\end{theorem}

\begin{proof}
The exact decomposition follows directly from Lemma~\ref{lem:mean_error_decomp}. It remains to express the analysis mean in a form amenable to truncation and projector perturbation bounds and then to estimate the resulting bias and variance contributions.

\smallskip
\noindent By the QPCA-EnDCF construction (Definition~\ref{def:endcf_gain_ensemble} and Definition~\ref{def:projectors_truncation}), the analysis update acts on the whitened residuals through the empirical rank-$\kappa$ projector:
\[
\mathbf{x}^{(j),a,s}
=
\mathbf{x}^{(j),f}
-
\mathbf{K}^{\mathrm{DC}}(\mathbf{R}^{(L)})^{1/2}\hat{\mathbf{P}}_\kappa^{\,s}\,\mathbf{e}^{(j),s}.
\]
Averaging over $j$ gives the corresponding relation for the analysis mean,
\begin{equation}
\label{eq:mean_update_projector_main}
\bar{\mathbf{x}}^{a,s}
=
\bar{\mathbf{x}}_L^{f}
-
\mathbf{K}^{\mathrm{DC}}(\mathbf{R}^{(L)})^{1/2}\hat{\mathbf{P}}_\kappa^{\,s}\,\bar{\mathbf{e}}^{\,s}.
\end{equation}

\smallskip
	\noindent Insert and subtract the population-truncated correction $(\mathbf{R}^{(L)})^{1/2}\mathbf{P}_\kappa^{\,s}\boldsymbol{\mu}_E^{\,s}$ in \eqref{eq:mean_update_projector_main} and decompose $\bar{\mathbf{e}}^{\,s}=\boldsymbol{\mu}_E^{\,s}+(\bar{\mathbf{e}}^{\,s}-\boldsymbol{\mu}_E^{\,s})$ to obtain
	\begin{align*}
		\mathbb{E}[\bar{\mathbf{x}}^{a,s}]
		&=
		\boldsymbol{\mu}_L^f
		-
		\mathbb{E}\!\left[\mathbf{K}^{\mathrm{DC}}(\mathbf{R}^{(L)})^{1/2}\mathbf{P}_\kappa^{\,s}\,\boldsymbol{\mu}_E^{\,s}\right]\\
	&\quad-
	\mathbb{E}\!\left[\mathbf{K}^{\mathrm{DC}}(\mathbf{R}^{(L)})^{1/2}(\hat{\mathbf{P}}_\kappa^{\,s}-\mathbf{P}_\kappa^{\,s})\,\boldsymbol{\mu}_E^{\,s}\right]\\
	&\quad-
	\mathbb{E}\!\left[\mathbf{K}^{\mathrm{DC}}(\mathbf{R}^{(L)})^{1/2}\hat{\mathbf{P}}_\kappa^{\,s}\,(\bar{\mathbf{e}}^{\,s}-\boldsymbol{\mu}_E^{\,s})\right].
	\end{align*}
	Add and subtract the untruncated population correction $\mathbb{E}[\mathbf{K}^{\mathrm{DC}}](\mathbf{R}^{(L)})^{1/2}\boldsymbol{\mu}_E^{\,s}$ to isolate the residual (``base'') analysis bias:
	\[
		\mathbf{b}_{\mathrm{base}}^{\,s}
		:=
		\boldsymbol{\mu}_L^f-\mathbf{x}^{\mathrm{true}}-\mathbb{E}[\mathbf{K}^{\mathrm{DC}}](\mathbf{R}^{(L)})^{1/2}\boldsymbol{\mu}_E^{\,s}.
		\]
	Using $\mathbb{E}[\mathbf{K}^{\mathrm{DC}}(\mathbf{R}^{(L)})^{1/2}\boldsymbol{\mu}_E^{\,s}]=\mathbb{E}[\mathbf{K}^{\mathrm{DC}}](\mathbf{R}^{(L)})^{1/2}\boldsymbol{\mu}_E^{\,s}$, we obtain
	\begin{align*}
	\mathbb{E}[\bar{\mathbf{x}}^{a,s}]-\mathbf{x}^{\mathrm{true}}
	&=
	\mathbf{b}_{\mathrm{base}}^{\,s}
	+
	\mathbb{E}\!\left[\mathbf{K}^{\mathrm{DC}}(\mathbf{R}^{(L)})^{1/2}(\mathbf{I}-\mathbf{P}_\kappa^{\,s})\boldsymbol{\mu}_E^{\,s}\right]\\
	&\qquad
	-
	\mathbb{E}\!\left[\mathbf{K}^{\mathrm{DC}}(\mathbf{R}^{(L)})^{1/2}(\hat{\mathbf{P}}_\kappa^{\,s}-\mathbf{P}_\kappa^{\,s})\,\boldsymbol{\mu}_E^{\,s}\right]
	-
	\mathbb{E}\!\left[\mathbf{K}^{\mathrm{DC}}(\mathbf{R}^{(L)})^{1/2}\hat{\mathbf{P}}_\kappa^{\,s}\,(\bar{\mathbf{e}}^{\,s}-\boldsymbol{\mu}_E^{\,s})\right].
	\end{align*}
	Consequently,
	\[
	\mathrm{Bias}^2(\kappa,s)
	\le
	2\|\mathbf{b}_{\mathrm{base}}^{\,s}\|^2
	+
	2\left\|\mathbb{E}[\bar{\mathbf{x}}^{a,s}]-\mathbf{x}^{\mathrm{true}}-\mathbf{b}_{\mathrm{base}}^{\,s}\right\|^2.
	\]
	The truncation remainder relative to the full population correction is governed by $(\mathbf{I}-\mathbf{P}_\kappa^{\,s})\boldsymbol{\mu}_E^{\,s}$. By Jensen's inequality and submultiplicativity of the spectral norm,
	\[
	\left\|\mathbb{E}\!\left[\mathbf{K}^{\mathrm{DC}}(\mathbf{R}^{(L)})^{1/2}(\mathbf{I}-\mathbf{P}_\kappa^{\,s})\boldsymbol{\mu}_E^{\,s}\right]\right\|^2
	\le
	\mathbb{E}\|\mathbf{K}^{\mathrm{DC}}\|_2^2\,\|\mathbf{R}^{(L)}\|_2\,\|(\mathbf{I}-\mathbf{P}_\kappa^{\,s})\boldsymbol{\mu}_E^{\,s}\|^2,
	\]
	which yields the truncation contribution in \eqref{eq:bias_bound_fixed_main}.
	\end{proof}


\begin{theorem}[Variance Decomposition with Spectral Truncation]
\label{thm:variance_decomposition}Under the Bias Variance Hypotheses~\ref{def:bias_variance_hypotheses}, the variance satisfies
\begin{align}
\mathrm{Var}(N,\kappa,s)
&\le
\frac{2}{N}\,\mathbb{E}\!\left[\|\mathbf{x}_L^{(1),f}-\boldsymbol{\mu}_L^f\|^2\right]
+
\frac{2\|\mathbf{R}^{(L)}\|_2}{N}\,
\mathbb{E}\!\left[\|\mathbf{K}^{\mathrm{DC}}\|_2^2\right]\,
\mathrm{tr}\!\bigl(\boldsymbol{\Sigma}_E^{\,s}\bigr)
\nonumber\\
&\hspace{3.2em}
+\;
2\|\mathbf{R}^{(L)}\|_2\,
\mathbb{E}\!\left[\|\mathbf{K}^{\mathrm{DC}}\|_2^2\,
\|(\hat{\mathbf{P}}_\kappa^{\,s}-\mathbf{P}_\kappa^{\,s})\,(\bar{\mathbf{e}}^{\,s}-\boldsymbol{\mu}_E^{\,s})\|^2\right],
\label{eq:var_bound_fixed_main_raw}
\end{align}
Moreover, the projector factor admits the marginal bound (by Cauchy--Schwarz and Lemma~\ref{lem:covariance_concentration}(iii) applied with $(\boldsymbol{\Sigma}_E,\mathbf{C}_E)$ replaced by $(\boldsymbol{\Sigma}_E^{\,s},\mathbf{C}_E^{\,s})$):
\begin{align}
\mathbb{E}\!\left[\|(\hat{\mathbf{P}}_\kappa^{\,s}-\mathbf{P}_\kappa^{\,s})\,(\bar{\mathbf{e}}^{\,s}-\boldsymbol{\mu}_E^{\,s})\|^2\right]
&\le
\Bigl(\mathbb{E}\|\hat{\mathbf{P}}_\kappa^{\,s}-\mathbf{P}_\kappa^{\,s}\|_F^4\Bigr)^{1/2}
\Bigl(\mathbb{E}\|\bar{\mathbf{e}}^{\,s}-\boldsymbol{\mu}_E^{\,s}\|^4\Bigr)^{1/2}
\nonumber\\
&\le
\frac{8\kappa}{\delta_\kappa^2}\,
\Bigl(\mathbb{E}\|\mathbf{C}_E^{\,s}-\boldsymbol{\Sigma}_E^{\,s}\|_F^4\Bigr)^{1/2}
\Bigl(\mathbb{E}\|\bar{\mathbf{e}}^{\,s}-\boldsymbol{\mu}_E^{\,s}\|^4\Bigr)^{1/2}.
\label{eq:var_projector_term_bound_main}
\end{align}
We record \eqref{eq:var_projector_term_bound_main} separately because it does not, by itself, control the gain-weighted joint term in \eqref{eq:var_bound_fixed_main_raw} without additional assumptions on $\mathbf{K}^{\mathrm{DC}}$ beyond second-moment finiteness.
\end{theorem}




\begin{proof}
The sample-mean fluctuation term is controlled without any independence assumption. Using $\|\hat{\mathbf{P}}_\kappa^{\,s}\|_2\le 1$, Jensen's inequality, submultiplicativity, and Cauchy--Schwarz,
\begin{multline*}
\left\|\mathbb{E}\!\left[\mathbf{K}^{\mathrm{DC}}(\mathbf{R}^{(L)})^{1/2}\hat{\mathbf{P}}_\kappa^{\,s}\,(\bar{\mathbf{e}}^{\,s}-\boldsymbol{\mu}_E^{\,s})\right]\right\|^2
\le
\|\mathbf{R}^{(L)}\|_2
\left(\mathbb{E}\!\left[\|\mathbf{K}^{\mathrm{DC}}\|_2\,\|\bar{\mathbf{e}}^{\,s}-\boldsymbol{\mu}_E^{\,s}\|\right]\right)^2\\
\le
\|\mathbf{R}^{(L)}\|_2\,\mathbb{E}\|\mathbf{K}^{\mathrm{DC}}\|_2^2\,
\mathbb{E}\|\bar{\mathbf{e}}^{\,s}-\boldsymbol{\mu}_E^{\,s}\|^2.
\end{multline*}
Since $\mathbb{E}\|\bar{\mathbf{e}}^{\,s}-\boldsymbol{\mu}_E^{\,s}\|^2=N^{-1}\mathrm{tr}(\boldsymbol{\Sigma}_E^{\,s})$, this yields the $\mathcal{O}(N^{-1})$ contribution in \eqref{eq:bias_bound_fixed_main}.

	The projector-estimation term is bounded similarly. Indeed, by Jensen's inequality, submultiplicativity, $\|(\mathbf{R}^{(L)})^{1/2}\|_2^2=\|\mathbf{R}^{(L)}\|_2$, $\|\hat{\mathbf{P}}_\kappa^{\,s}-\mathbf{P}_\kappa^{\,s}\|_2\le \|\hat{\mathbf{P}}_\kappa^{\,s}-\mathbf{P}_\kappa^{\,s}\|_F$, and Cauchy--Schwarz (no independence assumption),
	\begin{align*}
	\left\|\mathbb{E}\!\left[\mathbf{K}^{\mathrm{DC}}(\mathbf{R}^{(L)})^{1/2}(\hat{\mathbf{P}}_\kappa^{\,s}-\mathbf{P}_\kappa^{\,s})\boldsymbol{\mu}_E^{\,s}\right]\right\|^2
	&\le
	\|\mathbf{R}^{(L)}\|_2\,\|\boldsymbol{\mu}_E^{\,s}\|^2
	\left(\mathbb{E}\!\left[\|\mathbf{K}^{\mathrm{DC}}\|_2\,\|\hat{\mathbf{P}}_\kappa^{\,s}-\mathbf{P}_\kappa^{\,s}\|_F\right]\right)^2\\
	&\le
	\mathbb{E}\|\mathbf{K}^{\mathrm{DC}}\|_2^2\,\|\mathbf{R}^{(L)}\|_2\,
	\mathbb{E}\|\hat{\mathbf{P}}_\kappa^{\,s}-\mathbf{P}_\kappa^{\,s}\|_F^2\,\|\boldsymbol{\mu}_E^{\,s}\|^2.
	\end{align*}
	Invoking Lemma~\ref{lem:covariance_concentration}(iii) gives
	\(
	\mathbb{E}\|\hat{\mathbf{P}}_\kappa^{\,s}-\mathbf{P}_\kappa^{\,s}\|_F^2
\lesssim
\frac{\kappa}{\delta_\kappa^2}\mathbb{E}\|\mathbf{C}_E^{\,s}-\boldsymbol{\Sigma}_E^{\,s}\|_F^2,
\)
and absorbing $\|\boldsymbol{\mu}_E^{\,s}\|^2$ into the constant yields the projector-estimation term in \eqref{eq:bias_bound_fixed_main}.

Finally, the map-approximation term $C_{\mathrm{m}}\|Q_s-Q\|_{L^2(P_{\mathcal{X}})}^2$ follows from stability of the mean update with respect to $Q_s\to Q$ in $L^2(P_{\mathcal{X}})$ (as in the continuity framework in~\cite{ButlerWildeyZhang2020}); we record it separately because it is orthogonal to the sampling and truncation mechanisms.

\smallskip
\noindent Subtract $\mathbb{E}[\bar{\mathbf{x}}^{a,s}]$ from \eqref{eq:mean_update_projector_main} and apply Young's inequality (i.e., $2ab\le a^2+b^2$) to separate (i) the fluctuation of the forecast mean and (ii) the fluctuation of the projected residual term. Using the standard identities
\[
\mathbb{E}\|\bar{\mathbf{x}}_L^f-\boldsymbol{\mu}_L^f\|^2=\frac{1}{N}\mathbb{E}\|\mathbf{x}_L^{(1),f}-\boldsymbol{\mu}_L^f\|^2,
\qquad
\mathbb{E}\|\bar{\mathbf{e}}^{\,s}-\boldsymbol{\mu}_E^{\,s}\|^2=\frac{1}{N}\mathrm{tr}\!\bigl(\boldsymbol{\Sigma}_E^{\,s}\bigr),
\]
we obtain \eqref{eq:var_bound_fixed_main_raw}. The remaining projector term reflects the fact that $\hat{\mathbf{P}}_\kappa^{\,s}-\mathbf{P}_\kappa^{\,s}$ and $\bar{\mathbf{e}}^{\,s}-\boldsymbol{\mu}_E^{\,s}$ are generally dependent, since both are computed from the same ensemble. The estimate \eqref{eq:var_projector_term_bound_main} follows by applying Cauchy--Schwarz to fourth moments and then invoking Lemma~\ref{lem:covariance_concentration}(iii) applied to $(\boldsymbol{\Sigma}_E^{\,s},\mathbf{C}_E^{\,s})$.
\end{proof}

\par\smallskip\noindent\refstepcounter{remark}\textbf{Remark~\theremark}
\label{rem:truncation_bias_tail}
The leading truncation contribution in \eqref{eq:bias_bound_fixed_main} depends on
\(
\|(\mathbf{I}-\mathbf{P}_\kappa^{\,s})\boldsymbol{\mu}_E^{\,s}\|^2,
\)
that is, on the component of the mean whitened innovation discarded by truncation. Bounds of the form
\(
\|(\mathbf{I}-\mathbf{P}_\kappa^{\,s})\boldsymbol{\mu}_E^{\,s}\|^2\lesssim \sum_{i>\kappa}\lambda_i^{\,s}
\)
do not hold in general; they require additional structure linking $\boldsymbol{\mu}_E^{\,s}$ to $\boldsymbol{\Sigma}_E^{\,s}$, for example a source/RKHS-type condition such as $\|(\boldsymbol{\Sigma}_E^{\,s})^{-1/2}\boldsymbol{\mu}_E^{\,s}\|<\infty$.
Theorem~\ref{thm:bias_variance} decomposes the mean-squared error into: (i) a base analysis-bias term $\mathrm{Bias}_{\mathrm{base}}^{2}(s)$ reflecting residual error after the (untruncated) population mean correction is applied; (ii) additional truncation bias from restricting the correction to a $\kappa$-dimensional subspace; (iii) approximation error from using $Q_s$ in place of $Q$; and (iv) sampling variability, including the stability factor $\kappa/\delta_\kappa^2$ associated with estimating a rank-$\kappa$ projector from $N$ samples. The pseudoinverse gain is essential in the small-ensemble regime $N-1<d$, and Assumption~\ref{ass:gain_moment} is the only probabilistic requirement needed to justify the gain-dependent mean-square bounds.

\FloatBarrier



% ================================================================================


\subsection{Corollaries and Practical Implications}
\label{subsec:corollaries_implications}

Theorem~\ref{thm:bias_variance} yields two immediate and practically relevant messages. First, the variance contribution $\mathrm{Var}(N,\kappa,s)$ decays at the canonical Monte Carlo rate $N^{-1}$, with a prefactor that depends on the effective dimension $\kappa$ and the cutoff gap $\delta_\kappa$ through the stability of the empirical projector. Second, apart from the base analysis-bias term $\mathrm{Bias}_{\mathrm{base}}^{2}(s)$, the additional bias separates into (i) a truncation component governed by the alignment of the mean innovation $\boldsymbol{\mu}_E^{\,s}$ with the discarded subspace and (ii) a model-approximation component controlled by $\|Q_s-Q\|_{L^2(P_{\mathcal{X}})}$.

%--------------------------------------------------------------------
\begin{corollary}[Scaling of the projector-estimation term (Gaussian case)]
\label{cor:projector_term_scaling_gaussian}
Under the Bias Variance Hypothesis~\ref{ass:regularity_bias_variance}, in addition, assume that the centered whitened residuals are Gaussian,
\[
\mathbf{e}^{(j),s}-\boldsymbol{\mu}_E^{\,s}\stackrel{\mathrm{i.i.d.}}{\sim}\mathcal{N}(\mathbf{0},\boldsymbol{\Sigma}_E^{\,s}).
\]
Then there exists a constant $C>0$, depending only on the observation-space dimension $d=mL$, such that
\begin{multline*}
\mathbb{E}\!\left[\|(\hat{\mathbf{P}}_\kappa^{\,s}-\mathbf{P}_\kappa^{\,s})\,(\bar{\mathbf{e}}^{\,s}-\boldsymbol{\mu}_E^{\,s})\|^2\right]
\le
\frac{C\,\kappa}{\delta_\kappa^2}\,
\frac{1}{N(N-1)}\,
\mathrm{tr}\!\bigl((\boldsymbol{\Sigma}_E^{\,s})^2\bigr)\\
\times\left(\mathbb{E}\|\mathbf{e}^{(1),s}-\boldsymbol{\mu}_E^{\,s}\|^4+\bigl(\mathbb{E}\|\mathbf{e}^{(1),s}-\boldsymbol{\mu}_E^{\,s}\|^2\bigr)^2\right)^{1/2}.
\end{multline*}
In particular, for fixed $d$ this term is $\mathcal{O}(\kappa/(\delta_\kappa^2N^2))$.
\end{corollary}

\begin{proof}
Combine \eqref{eq:var_projector_term_bound_main} with Lemma~\ref{lem:cov_frob_fourth_moment} applied to $(\mathbf{C}_E^{\,s},\boldsymbol{\Sigma}_E^{\,s})$ under the stated Gaussian hypothesis, and with Lemma~\ref{lem:sample_mean_fourth_moment} applied to $\bar{\mathbf{e}}^{\,s}$. The resulting bound is of order $(N-1)^{-1}\times N^{-1}$, and the constants can be absorbed into a single $C>0$ depending only on $d$.
\end{proof}

%--------------------------------------------------------------------
\begin{corollary}[Comparison to stochastic EnKF observation perturbations]
\label{cor:stochastic_comparison_new}
Under the Bias Variance Hypothesis~\ref{ass:regularity_bias_variance}, consider a stochastic windowed EnKF update with perturbed observations
\[
\mathbf{x}^{(j),a}_{\mathrm{stoch}}
=
\mathbf{x}^{(j),f}
+
\mathbf{K}^{(w)}\Bigl(\mathbf{z}^{(w)}+\boldsymbol{\epsilon}_{w}^{(j)}-\mathbf{H}^{(L)}\mathbf{x}^{(j),f}\Bigr),
\qquad
\boldsymbol{\epsilon}_{w}^{(j)}\stackrel{\mathrm{i.i.d.}}{\sim}\mathcal{N}(\mathbf{0},\mathbf{R}^{(L)}),
\]
where $\mathbf{K}^{(w)}$ is deterministic.
Let $\bar{\mathbf{x}}^{a}_{\mathrm{stoch}}$ denote the stochastic analysis mean and let $\bar{\mathbf{x}}^{a}_{\mathrm{det}}$
denote the corresponding mean obtained by setting $\boldsymbol{\epsilon}_{w}^{(j)}\equiv 0$. Then
\[
\bar{\mathbf{x}}^{a}_{\mathrm{stoch}}
=
\bar{\mathbf{x}}^{a}_{\mathrm{det}}
+
\frac{1}{N}\sum_{j=1}^N \mathbf{K}^{(w)}\boldsymbol{\epsilon}_{w}^{(j)},
\]
and hence
\[
\mathrm{Cov}\!\left(\bar{\mathbf{x}}^{a}_{\mathrm{stoch}}\right)
=
\mathrm{Cov}\!\left(\bar{\mathbf{x}}^{a}_{\mathrm{det}}\right)
+
\frac{1}{N}\mathbf{K}^{(w)}\mathbf{R}^{(L)}(\mathbf{K}^{(w)})^\top.
\]
In particular,
\begin{equation}
\label{eq:stoch_var_bound_new}
\begin{split}
\mathbb{E}\!\left[\left\|\bar{\mathbf{x}}^{a}_{\mathrm{stoch}}-\mathbb{E}[\bar{\mathbf{x}}^{a}_{\mathrm{stoch}}]\right\|^2\right]
&\;\ge\;
\mathbb{E}\!\left[\left\|\frac{1}{N}\sum_{j=1}^N \mathbf{K}^{(w)}\boldsymbol{\epsilon}_{w}^{(j)}\right\|^2\right]\\
&=
\frac{1}{N}\,\mathrm{tr}\!\bigl(\mathbf{K}^{(w)}\mathbf{R}^{(L)}(\mathbf{K}^{(w)})^\top\bigr)
\;\le\;
\frac{\|\mathbf{K}^{(w)}\|_2^2\,\mathrm{tr}(\mathbf{R}^{(L)})}{N}.
\end{split}
\end{equation}
If $\mathbf{R}^{(L)}=\bar r\,\mathbf{I}$, then $\mathrm{tr}(\mathbf{R}^{(L)})=d\,\bar r$ and the variance scales as
$\mathcal{O}(d/N)$.

The QPCA-EnDCF variance satisfies
$\mathrm{Var}(N,\kappa,s)=\mathcal{O}(1/N)$ (for fixed approximation stage $s$), with an additional stability factor $\kappa/\delta_\kappa^2$ arising from estimation of the rank-$\kappa$ projector (as in~\eqref{eq:var_projector_term_bound_main}).
\end{corollary}

\begin{proof}
The identity for $\bar{\mathbf{x}}^{a}_{\mathrm{stoch}}$ follows by averaging the memberwise update and collecting the perturbation terms. The covariance relation then follows from independence and
\[
\mathrm{Cov}\!\left(\frac{1}{N}\sum_{j=1}^N \mathbf{K}^{(w)}\boldsymbol{\epsilon}_{w}^{(j)}\right)
=
\frac{1}{N}\mathbf{K}^{(w)}\mathbf{R}^{(L)}(\mathbf{K}^{(w)})^\top.
\]
The final comparison for QPCA-EnDCF is a direct restatement of Theorem~\ref{thm:bias_variance} together with \eqref{eq:var_projector_term_bound_main}.
\end{proof}

Equation~\eqref{eq:stoch_var_bound_new} isolates an irreducible $\mathcal{O}(N^{-1})$ contribution to the variance of the stochastic analysis mean arising solely from observation perturbations. In empirical implementations, $\mathbf{K}^{(w)}$ is typically computed from sample covariances and is therefore random; the covariance identity holds conditionally on $\mathbf{K}^{(w)}$ (or on the forecast ensemble), and the unconditional variance term becomes $\frac{1}{N}\,\mathbb{E}\!\left[\mathbf{K}^{(w)}\mathbf{R}^{(L)}(\mathbf{K}^{(w)})^\top\right]$. QPCA-EnDCF eliminates this specific source of sampling noise, but introduces variability through estimation of the low-rank projector; the latter contribution is controlled by the cutoff gap $\delta_\kappa$ and by the covariance concentration rate.
\par\smallskip


\FloatBarrier




% ================================================================================




% ================================================================================
% SECTION: Bias--Variance Tradeoff
% ================================================================================
\section{Bias--Variance Tradeoff}\label{sec:bias_variance_results}

The calibration diagnostics in Section~\ref{subsec:calibration} established that QPCA-EnDCF yields well-calibrated uncertainty, as measured by spread--skill relationships and rank histograms. These results, however, do not explain why calibration improves. We address this question by decomposing the mean-squared error (MSE) into bias and variance components. Theory predicts that spectral truncation reduces sampling variance, while any induced truncation bias is governed by the discarded component $\|(\mathbf{I}-\mathbf{P}_\kappa^{\,s})\boldsymbol{\mu}_E^{\,s}\|$ and can remain small when the mean innovation is well aligned with the leading modes; this contrasts with classical regularization strategies, where variance reduction typically comes with a more direct bias penalty. Here we provide empirical validation and quantify the relative contributions of bias and variance. For QPCA-EnDCF, Theorem~\ref{thm:bias_variance} yields $\mathrm{Variance}=\mathcal{O}(1/N)$ with $\kappa$-dependent prefactors arising through projector stability (via $\kappa/\delta_\kappa^2$) and through the observation-space term involving $\mathrm{tr}(\boldsymbol{\Sigma}_E)$. The bias contribution is governed by $\|(\mathbf{I}-\mathbf{P}_\kappa)\boldsymbol{\mu}_E\|^2$, the component of the mean whitened innovation orthogonal to the retained eigenspace (see Theorem~\ref{thm:bias_variance} and Remark~\ref{rem:truncation_bias_tail}). Importantly, this cannot be bounded simply by the discarded spectral tail $\sum_{i>\kappa}\lambda_i$ without additional structure linking $\boldsymbol{\mu}_E$ to $\boldsymbol{\Sigma}_E$ (e.g., a source condition). However, when the mean mismatch is well-aligned with the leading covariance directions, which occurs empirically when the dominant eigenmode captures systematic forecast--observation discrepancy. The aggressive truncation ($\kappa=1$) can reduce sampling variance without a commensurate bias penalty. By contrast, standard ensemble updates exhibit variance scaling as $\mathcal{O}(r/N)$, where $r$ approaches the effective rank of the update. We estimate bias and variance empirically using $N_{\mathrm{trial}}=5$ independent ensemble realizations per algorithm, differing only in initial ensemble draws. (Note: $N_{\mathrm{trial}}$ denotes the number of Monte Carlo trials, distinct from the ensemble size $N=10$ and the total rank samples $M_{\mathrm{rank}}$ used in Section~\ref{subsec:rank_histograms}.) Let $\bar{\mathbf{x}}_{t,w}^a$ denote the analysis ensemble mean at the end of window $w$ in trial $t$, and let $\mathbf{x}_{w}^{\mathrm{true}}$ denote the corresponding true state (both evaluated at window endpoints). For each window $w$, define the across-trial mean
\[
\bar{\mathbf{x}}_{\cdot,w}^a:=\frac{1}{N_{\mathrm{trial}}}\sum_{t=1}^{N_{\mathrm{trial}}}\bar{\mathbf{x}}_{t,w}^a.
\]
The window-wise estimators are
\begin{align}
\widehat{\mathrm{Bias}}_{w}^2
&=
\left\|
\bar{\mathbf{x}}_{\cdot,w}^a-\mathbf{x}_{w}^{\mathrm{true}}
\right\|^2,
\\
\widehat{\mathrm{Variance}}_{w}
&=
\frac{1}{N_{\mathrm{trial}}}\sum_{t=1}^{N_{\mathrm{trial}}}
\left\|
\bar{\mathbf{x}}_{t,w}^a-\bar{\mathbf{x}}_{\cdot,w}^a
\right\|^2,
\end{align}
with $\widehat{\mathrm{MSE}}_{w}=\widehat{\mathrm{Bias}}_{w}^2+\widehat{\mathrm{Variance}}_{w}$. All norms are total (summed over $n$ state components), so bias$^2$, variance, and MSE values reported below are $\mathcal{O}(n)$ in magnitude. To summarize performance over the full run, we time-average these window-wise quantities:
\[
\widehat{\mathrm{Bias}}^2:=\frac{1}{W}\sum_{w=1}^{W}\widehat{\mathrm{Bias}}_{w}^2,
\qquad
\widehat{\mathrm{Variance}}:=\frac{1}{W}\sum_{w=1}^{W}\widehat{\mathrm{Variance}}_{w},
\qquad
\widehat{\mathrm{MSE}}:=\frac{1}{W}\sum_{w=1}^{W}\widehat{\mathrm{MSE}}_{w}.
\]

Figure~\ref{fig:bias_variance_evolution} shows the window-wise bias--variance decomposition across the 50-window sequence for Sequential EnKF, 4D-EnKF, and QPCA-EnDCF. For each window $w$, the stacked areas depict $\widehat{\mathrm{Bias}}_{w}^2$ (solid) and $\widehat{\mathrm{Variance}}_{w}$ (hatched), whose sum equals $\widehat{\mathrm{MSE}}_{w}$. Estimates are computed from $N_{\mathrm{trial}}=10$ independent realizations.

\begin{figure}[htbp]
\centering
\includegraphics[width=\textwidth]{figures/fig07a_bias_variance_evolution.png}
\caption{Window-wise bias--variance decomposition over 50 assimilation windows (left to right: Sequential EnKF, 4D-EnKF, QPCA-EnDCF). Solid fill shows $\widehat{\mathrm{Bias}}_{w}^2$ and hatched fill shows $\widehat{\mathrm{Variance}}_{w}$; the total height is $\widehat{\mathrm{MSE}}_{w}=\widehat{\mathrm{Bias}}_{w}^2+\widehat{\mathrm{Variance}}_{w}$. Estimates use $N_{\mathrm{trial}}=10$ independent initial-ensemble realizations.}
\label{fig:bias_variance_evolution}
\end{figure}

QPCA-EnDCF attains MSE $\approx13$, compared to $\approx22$ for Sequential EnKF and $\approx21$ for 4D-EnKF. The decomposition reveals the mechanism: squared bias is comparable across methods (Bias$^2\approx10$), while variance differs sharply. QPCA-EnDCF exhibits variance $\approx2.3$, an $\sim80\%$ reduction relative to stochastic methods (variance $\approx11.6$). Thus, aggressive truncation removes variance-dominated modes while retaining the leading mode capturing systematic mismatch. The time-averaged Bias$^2$/MSE ratio (computed from the same window-wise estimates) further clarifies this distinction. For QPCA-EnDCF, the ratio is about $82\%$, indicating error dominated by systematic bias rather than sampling noise. In contrast, stochastic methods show ratios near $45$--$47\%$, implying that roughly half of their error arises from sampling variance at fixed ensemble size. Figure~\ref{fig:mse_decomposition_bars} summarizes these results via stacked bar charts.

\begin{figure}[htbp]
\centering
\includegraphics[width=\textwidth]{figures/fig07b_mse_decomposition_bars.png}
\caption{Stacked-bar MSE decomposition. Left: absolute contributions of squared bias and variance. Right: percentage contributions normalized to unit height. QPCA-EnDCF exhibits comparable bias but markedly reduced variance relative to stochastic methods.}
\label{fig:mse_decomposition_bars}
\end{figure}

QPCA-EnDCF achieves MSE $=12.96$ (Bias$^2=10.67$, Variance$=2.29$), while Sequential EnKF and 4D-EnKF exhibit MSE $\approx21$--$22$ with variance contributions near $11.6$. Normalized bars emphasize the qualitative shift: QPCA-EnDCF error is overwhelmingly systematic, whereas stochastic methods are variance-limited.  These findings validate the theoretical predictions and explain QPCA-EnDCF’s superior calibration. Variance reduction arises from two effects: (i) the deterministic update avoids observation-perturbation noise, whose contribution scales like $\|\mathbf{K}^{(w)}\|_2^2\,L\,\mathrm{tr}(\mathbf{R})/N$ (as in \eqref{eq:stoch_var_bound_new}) and is substantial under baseline parameters; and (ii) spectral truncation ($\kappa=1$) concentrates updates in the dominant eigenmode, avoiding estimation of noise-dominated directions. With variance scaling as $\mathcal{O}(1/N)$ (with $\kappa$-dependent prefactors) rather than $\mathcal{O}(r/N)$, the observed roughly fivefold reduction follows directly. Crucially, this variance reduction occurs without a bias penalty. Unlike classical penalty-based regularization, which trades variance for increased bias, QPCA-EnDCF maintains bias comparable to or lower than that of stochastic methods while sharply reducing variance. This favorable regime is enabled by rapid spectral decay, whereby leading modes encode signal and higher modes primarily represent sampling noise. Together with the calibration diagnostics, the bias--variance analysis completes the mechanistic picture. QPCA-EnDCF’s near-ideal spread--skill ratio arises because variance is suppressed without inflating bias, allowing ensemble spread to reflect true error. Stochastic methods suffer from both inflated variance and collapsed spread, yielding severe miscalibration. The dominance of systematic bias in QPCA-EnDCF suggests that further gains are attainable through model or operator improvements, whereas stochastic methods face fundamental variance limits at fixed ensemble size.
\FloatBarrier





% ================================================================================
% SECTION: Inflation Requirements and Variance Maintenance
% ================================================================================
\section{Covariance Inflation and Maintenance of Ensemble Spread}
\label{sec:inflation}

Ensemble data assimilation represents forecast uncertainty with a finite set of samples. When the ensemble size $N$ is small relative to the state dimension $n$, empirical forecast covariances systematically underestimate uncertainty. This deficiency is twofold: the covariance is rank deficient, spanning at most $N\!-\!1$ directions, and its nonzero eigenvalues are negatively biased. Left uncorrected, these effects accumulate over assimilation cycles, eroding ensemble spread and ultimately causing filter divergence. For this reason, operational ensemble methods routinely employ covariance inflation to restore variance.

This section examines whether QPCA-EnDCF mitigates variance loss intrinsically through spectral regularization, thereby reducing or eliminating the need for explicit inflation. The underlying hypothesis is that aggressive truncation concentrates updates onto signal-dominated modes while suppressing variance-depleting sampling noise, stabilizing ensemble spread without artificial augmentation. We test this hypothesis empirically by comparing inflation requirements across ensemble sizes and filtering methods.

Section~\ref{subsec:inflation_problem} reviews the origin of variance loss in finite ensembles and summarizes standard inflation mechanisms. Section~\ref{subsec:inflation_free} demonstrates that QPCA-EnDCF achieves optimal performance without inflation across all tested ensemble sizes, whereas stochastic methods require ensemble-size-dependent tuning. Eliminating inflation removes a major source of heuristic calibration and simplifies practical deployment.

% ================================================================================
\subsection{The Role of Inflation in Finite Ensemble Filtering}
\label{subsec:inflation_problem}

Consider an ensemble of size $N$ representing uncertainty in an $n$-dimensional state space, with $N \ll n$. The empirical forecast covariance
\[
\widehat{\mathbf{P}}^{\,f}
=
\frac{1}{N-1}\mathbf{A}^f(\mathbf{A}^f)^{\top},
\]
constructed from anomaly matrix $\mathbf{A}^f \in \mathbb{R}^{n\times N}$, has rank at most $N-1$ and therefore $n-(N-1)$ zero eigenvalues. Rank deficiency alone would be manageable if the remaining eigenvalues accurately captured uncertainty within the ensemble subspace. In practice, finite sampling induces systematic negative bias: even the nonzero eigenvalues underestimate the true variances.

This bias compounds in sequential assimilation. Each analysis step reduces variance along observed directions; if the incoming forecast covariance is already underestimated, subsequent updates further contract the ensemble. Over many cycles this feedback produces ensemble collapse, in which members concentrate in a narrow region of state space and fail to represent posterior uncertainty, preventing effective assimilation of new observations. Covariance inflation counteracts this variance loss by artificially augmenting ensemble spread. Two strategies dominate operational practice. Multiplicative inflation rescales anomalies by a factor $\lambda_{\mathrm{infl}} \ge 1$,
\[
\mathbf{X}^f_{\lambda_{\mathrm{infl}}}
=
\bar{\mathbf{x}}^f\mathbf{1}^{\top}
+
\lambda_{\mathrm{infl}} \mathbf{A}^f,
\]
where $\bar{\mathbf{x}}^f$ is the ensemble mean. The resulting covariance satisfies
$\widehat{\mathbf{P}}^{\,f}_{\lambda_{\mathrm{infl}}}=\lambda_{\mathrm{infl}}^2\widehat{\mathbf{P}}^{\,f}$,
uniformly amplifying all eigenvalues while preserving eigenvectors. Geometrically, this is a uniform dilation of the ensemble ellipsoid. Multiplicative inflation is inexpensive and preserves correlation structure, but it cannot address rank deficiency: the null space is unchanged. Additive inflation perturbs each ensemble member with independent noise,
\[
\mathbf{x}^{(j)}_{\alpha_{\mathrm{add}}}
=
\mathbf{x}^{(j),f}+\boldsymbol{\varepsilon}^{(j)},
\qquad
\boldsymbol{\varepsilon}^{(j)}\sim\mathcal{N}(\mathbf{0},\alpha_{\mathrm{add}}^2\mathbf{Q}_{\mathrm{add}}),
\]
where $\mathbf{Q}_{\mathrm{add}}$ prescribes the perturbation structure. Unlike multiplicative inflation, additive noise can inject variance into directions orthogonal to the ensemble subspace, partially alleviating rank deficiency. This flexibility comes at a cost: unless $\mathbf{Q}_{\mathrm{add}}$ reflects dynamical correlations, the injected variance may be physically implausible. For simplicity, $\mathbf{Q}_{\mathrm{add}}$ is often taken as $\mathbf{I}_n$, prioritizing convenience over fidelity.

Optimal inflation parameters depend sensitively on ensemble size, observation density, and system dynamics, making manual tuning difficult. Adaptive schemes reduce this burden but introduce additional complexity and overhead. These limitations motivate the present investigation: whether spectral regularization via QPCA internalizes variance maintenance, rendering explicit inflation unnecessary.





% ================================================================================



\subsection{Empirical Investigation of Inflation Requirements}
\label{subsec:inflation_free}

We perform a systematic parametric study to quantify inflation requirements across ensemble sizes and filtering methods. Ensemble size is varied over $N \in \{5,10,25,50\}$, multiplicative inflation over $\lambda_{\mathrm{infl}} \in \{1.0,1.05,1.10,1.15,1.20,1.25\}$, and additive inflation over $\alpha_{\mathrm{add}} \in \{0.0,0.01,0.05,0.10,0.15,0.20\}$, with all other parameters fixed at the baseline values of Section~\ref{sec:motivation}. For additive inflation we use the standard isotropic choice $\mathbf{Q}_{\mathrm{add}}=\mathbf{I}_n$. Each configuration is repeated with three independent random seeds. Performance is assessed by time-averaged RMSE over the final 50 assimilation windows, after transients have decayed.

Figure~\ref{fig:mult_inflation} reports RMSE as a function of the multiplicative inflation factor $\lambda_{\mathrm{infl}}$ for representative ensemble sizes $N \in \{5,10,25\}$. Each panel compares Sequential EnKF, 4D-EnKF, and QPCA-EnDCF; curves show mean RMSE and shaded bands denote $\pm1$ standard deviation across trials.

\begin{figure}[htbp]
\centering
\includegraphics[width=\textwidth]{figures/fig_inflation_multiplicative.png}
\caption{Sensitivity of RMSE to multiplicative inflation. Time-averaged RMSE versus inflation factor $\lambda_{\mathrm{infl}}$ for Sequential EnKF (blue circles), 4D-EnKF (orange squares), and QPCA-EnDCF (green triangles). Panels correspond to $N \in \{5,10,25\}$. Means are averaged over 50 assimilation windows and three trials; shading indicates $\pm1$ standard deviation across trials.}
\label{fig:mult_inflation}
\end{figure}

QPCA-EnDCF exhibits a roughly uniform response across ensemble sizes. Performance is insensitive to mild inflation ($\lambda_{\mathrm{infl}}\in[1.0,1.05]$) but slightly increases for larger values, indicating that spectral regularization alone maintains adequate ensemble spread. By contrast, Sequential EnKF and 4D-EnKF display ensemble-size-dependent optima. For $N=5$, minimum RMSE occurs near $\lambda_{\mathrm{infl}}=1.10$. This trend reflects diminishing finite-sample bias with growing ensemble size. The magnitude of this effect is substantial for small ensembles. At $N=5$, Sequential EnKF with optimal inflation ($\lambda_{\mathrm{infl}}=1.10$) attains RMSE $2.15\pm0.18$, compared with $1.68\pm0.12$ for QPCA-EnDCF at $\lambda_{\mathrm{infl}}=1.0$, a reduction of approximately 22\%. Figure~\ref{fig:add_inflation} shows the corresponding results for additive inflation with $\mathbf{Q}_{\mathrm{add}}=\mathbf{I}_n$.

\begin{figure}[htbp]
\centering
\includegraphics[width=\textwidth]{figures/fig_inflation_additive.png}
\caption{Sensitivity of RMSE to additive inflation. Time-averaged RMSE versus additive inflation magnitude $\alpha_{\mathrm{add}}$ for Sequential EnKF, 4D-EnKF, and QPCA-EnDCF, with panels for $N \in \{5,10,25\}$. Plotting conventions match Figure~\ref{fig:mult_inflation}; shading indicates $\pm1$ standard deviation across three trials.}
\label{fig:add_inflation}
\end{figure}

QPCA-EnDCF again achieves its minimum RMSE at $\alpha_{\mathrm{add}}=0$, with slight monotonic degradation as $\alpha_{\mathrm{add}}$ increases. Since spectral regularization already controls variance, additive perturbations inject spurious energy that disrupts dynamically consistent correlations without compensating benefits.

For stochastic methods, additive inflation provides limited improvement relative to no inflation, but consistently underperforms multiplicative inflation by 5--10\%. This deficit reflects the difficulty of specifying an appropriate perturbation structure: isotropic noise introduces variance in dynamically irrelevant directions. More structured choices of $\mathbf{Q}_{\mathrm{add}}$ could improve performance but at the cost of additional modeling effort. The defining feature of Figure~\ref{fig:mult_inflation} and Figure~\ref{fig:add_inflation} is the invariance of QPCA-EnDCF's optimal parameters: $(\lambda_{\mathrm{opt}},\alpha_{\mathrm{opt}}) \approx (1.0,0.0)$ for all $N$. This eliminates the need for manual tuning or adaptive inflation schemes. In contrast, stochastic methods require ensemble-size-dependent inflation, with $\lambda_{\mathrm{opt}}$ decreasing toward unity as $N$ increases. The persistent inferiority of additive inflation relative to multiplicative inflation across all methods and ensemble sizes (visible in the consistently higher RMSE (add.) values) confirms that variance control through ensemble-wide rescaling is more effective than introducing isotropic random perturbations, which lack dynamical structure and contaminate correlations.


% ================================================================================

% ================================================================================
% SECTION: Robustness to Non-Ideal Observation Errors
% ================================================================================
\section{Robustness to Non-Ideal Observation Errors}
\label{sec:robustness}

Previous sections evaluate performance under idealized observational assumptions: Gaussian errors with known diagonal covariance. These assumptions simplify analysis and isolate algorithmic effects, but they are routinely violated in practice. Real observing systems may exhibit outliers (heavy tails), saturation or retrieval nonlinearities (skewness), and shared contamination or systematic biases (spatial correlations), all of which undermine the classical modeling assumptions. The operational question is whether QPCA-EnDCF’s gains persist under such non-idealities, or whether they are specific to the Gaussian/diagonal covariance setting. We probe robustness by varying observation error characteristics along two axes that isolate distinct failure modes. Section~\ref{subsec:nongaussian} studies distributional departures by testing nine error distributions spanning symmetric heavy-tailed, right-skewed, and Gaussian cases while keeping the second-moment structure correctly specified. Section~\ref{subsec:correlated_errors} studies structural departures by introducing spatial correlations with condition numbers spanning five orders of magnitude, again under correctly specified covariances. Together, these experiments distinguish sensitivity to higher-order distributional features (skewness, kurtosis, tail behavior) from sensitivity to correlation structure. Unless otherwise stated, all experiments use the baseline configuration of Section~\ref{sec:motivation} ($N=10$, $L=5$, $n=40$), modifying only the targeted observation model component.

% ================================================================================
\subsection{Non-Gaussian Observation Error Distributions}
\label{subsec:nongaussian}

Observation errors often depart from Gaussianity due to outliers, saturation, and nonlinear retrieval maps. Although Kalman-type derivations assume Gaussian noise, ensemble methods are frequently argued to be more robust because uncertainty is represented by samples rather than a parametric likelihood. Here we test whether the performance advantages observed under Gaussian errors persist under non-Gaussian noise. The design isolates shape effects from covariance misspecification: every distribution is standardized to mean zero and variance $\sigma_{\mathrm{obs}}^2=2.25$, and all methods are supplied the correct covariance $\mathbf{R}=2.25\mathbf{I}_m$. Thus any performance changes reflect higher-order properties rather than incorrect second-moment information. We consider nine distributions. The Gaussian $\mathcal{N}(0,\sigma_{\mathrm{obs}}^2)$ serves as a reference. Symmetric heavy-tailed cases include Student-$t$ with $\nu\in\{3,5,10\}$ degrees of freedom and Laplace with scale $b=\sigma_{\mathrm{obs}}/\sqrt{2}$, each centered at zero and rescaled to variance $\sigma_{\mathrm{obs}}^2$. Right-skewed cases include Gamma($k,\theta$) with shape $k\in\{2,5\}$ and Log-normal($\mu,\sigma_{\mathrm{LN}}$) with parameters $\sigma_{\mathrm{LN}}\in\{0.5,0.8\}$. These positive-valued distributions are shifted (by subtracting their mean) and rescaled to achieve mean zero and variance $\sigma_{\mathrm{obs}}^2$ while preserving their characteristic tail behavior and asymmetry. Thus, we use shifted and rescaled Gamma/Log-normal errors rather than the original positive-valued distributions.

Figure~\ref{fig:noise_distributions} illustrates representative samples. Despite pronounced differences in tails and skewness, all distributions share the same variance by construction.

\begin{figure}[htbp]
\centering
\includegraphics[width=\textwidth]{figures/noise_distributions.png}
\caption{Representative non-Gaussian observation error distributions used in the robustness study (Gaussian reference shown as a dashed curve). Student-$t$ and Laplace errors are symmetric and heavy-tailed; Gamma and Log-normal errors are right-skewed with increasingly heavy upper tails. All distributions are standardized to mean zero and variance $\sigma_{\mathrm{obs}}^2=2.25$, so comparisons isolate sensitivity to distributional shape.}
\label{fig:noise_distributions}
\end{figure}

For each distribution we run three independent trials, each consisting of 50 assimilation windows, yielding 150 window-level samples per method--distribution pair. Table~\ref{tab:nongaussian_performance} reports mean RMSE (trial-averaged over the 50 windows) and the relative improvement of QPCA-EnDCF compared with the stochastic methods.

\begin{table}[htbp]
\centering
\caption{Reconstruction error under non-Gaussian observation errors. Entries are mean RMSE across three independent trials, each time-averaged over 50 assimilation windows. All distributions have mean zero and variance $\sigma_{\mathrm{obs}}^2=2.25$, and all methods use the correctly specified diagonal covariance $\mathbf{R}=2.25\mathbf{I}_m$. Rightmost columns report the percentage RMSE reduction of QPCA-EnDCF relative to the stochastic methods (negative values indicate improvement).}
\label{tab:nongaussian_performance}
\small
\begin{tabular}{llccccc}
\toprule
\textbf{Distribution} & \textbf{Characteristics} & \textbf{QPCA} & \textbf{Seq} & \textbf{4D} & \textbf{vs. Seq} & \textbf{vs. 4D} \\
\midrule
Gaussian          & Baseline (reference)      & $\mathbf{3.63}$ & $4.45$ & $4.51$ & $-18.3\%$ & $-19.4\%$ \\
\midrule
\multicolumn{7}{l}{\textit{Heavy-tailed symmetric distributions:}} \\
Student-$t$ ($\nu$=3)  & Extreme heavy tail & $\mathbf{3.67}$ & $4.68$ & $4.46$ & $-21.6\%$ & $-17.8\%$ \\
Student-$t$ ($\nu$=5)  & Heavy tail         & $\mathbf{3.55}$ & $4.48$ & $4.58$ & $-20.7\%$ & $-22.4\%$ \\
Student-$t$ ($\nu$=10) & Moderate tail      & $\mathbf{3.51}$ & $4.54$ & $4.40$ & $-22.8\%$ & $-20.3\%$ \\
Laplace           & Double exponential & $\mathbf{3.69}$ & $4.51$ & $4.49$ & $-18.3\%$ & $-18.0\%$ \\
\midrule
\multicolumn{7}{l}{\textit{Right-skewed distributions:}} \\
Gamma ($k$=2)     & Moderate skewness  & $\mathbf{3.56}$ & $4.72$ & $4.45$ & $-24.5\%$ & $-19.9\%$ \\
Gamma ($k$=5)     & Mild skewness      & $\mathbf{3.66}$ & $4.51$ & $4.43$ & $-18.8\%$ & $-17.3\%$ \\
Log-normal ($\sigma$=0.5) & Skewed with heavy tail & $\mathbf{3.60}$ & $4.44$ & $4.66$ & $-18.9\%$ & $-22.7\%$ \\
Log-normal ($\sigma$=0.8) & Strong skew and tail  & $\mathbf{3.61}$ & $4.49$ & $4.55$ & $-19.6\%$ & $-20.6\%$ \\
\midrule
Cross-distribution mean & & $\mathbf{3.61}$ & $4.54$ & $4.50$ & $-20.4\%$ & $-19.8\%$ \\
Cross-distribution std. & & $\mathbf{0.052}$ & $0.092$ & $0.131$ & --- & --- \\
\bottomrule
\end{tabular}
\end{table}

QPCA-EnDCF is essentially insensitive to distributional shape: RMSE remains in the narrow band 3.51--3.69 across all nine distributions, from extreme heavy tails (Student-$t$, $\nu=3$) to strong right skew (Log-normal, $\sigma_{\mathrm{LN}}=0.8$). The cross-distribution standard deviation is $0.052$ (coefficient of variation $\approx 1.4\%$), indicating near-constant accuracy despite substantial changes in skewness and kurtosis. In contrast, the stochastic methods are both less accurate and more variable across distributions (cross-distribution standard deviations $0.092$ for Sequential EnKF and $0.131$ for 4D-EnKF). Across all distributions, QPCA-EnDCF improves RMSE by 18.3--24.5\% relative to Sequential EnKF (mean 20.4\%) and by 17.3--22.7\% relative to 4D-EnKF (mean 19.8\%). Figure~\ref{fig:mean_rmse_comparison} summarizes these gains and highlights the comparatively tight spread of QPCA-EnDCF across distributional families.

\begin{figure}[htbp]
\centering
\includegraphics[width=\textwidth]{figures/mean_rmse_comparison_nongaussian.png}
\caption{Mean RMSE under non-Gaussian observation errors (bars: mean; error bars: $\pm1$ standard deviation across $n=3$ trials). QPCA-EnDCF maintains consistently lower RMSE and smaller variability across all nine distributions, with 18.3--24.5\% improvement over Sequential EnKF and 17.3--22.7\% improvement over 4D-EnKF.}
\label{fig:mean_rmse_comparison}
\end{figure}

Temporal behavior is equally consistent. Figure~\ref{fig:rmse_trajectories} shows RMSE trajectories over the 50 assimilation windows: QPCA-EnDCF remains uniformly lower than both stochastic methods for every distribution, without deterioration over time.

\begin{figure}[htbp]
\centering
\includegraphics[width=\textwidth]{figures/rmse_trajectories_nongaussian.png}
\caption{RMSE evolution over 50 assimilation windows (after 10-window initialization) under each non-Gaussian error distribution (curves: mean; shading: 95\% confidence intervals across $n=3$ trials). QPCA-EnDCF maintains consistently lower RMSE than both stochastic methods across all distributions and time steps.}
\label{fig:rmse_trajectories}
\end{figure}

These results are consistent with the fact that QPCA-EnDCF is driven by second-moment structure. The whitening transform depends only on $\mathbf{R}$, which is identical and correctly specified in all experiments, and the subsequent spectral truncation is performed on an empirical covariance of whitened residuals. In contrast, stochastic EnKF variants rely on sampling perturbed observations from an assumed noise distribution, so changes in tail behavior or asymmetry can propagate through the update and increase variability. The deterministic, covariance-based update in QPCA-EnDCF avoids this mechanism, yielding the observed distributional robustness.






% ================================================================================


\subsection{Correlated Observation Errors}
\label{subsec:correlated_errors}

Observation errors in operational systems are rarely independent. Spatial correlations arise, for example, in satellite retrievals through shared atmospheric-path contamination and algorithmic dependencies; temporal correlations occur in ground networks through drift and common environmental forcing; and cross-channel correlations appear through shared calibration. When such correlations are ignored, the effective information content of the observations is overstated and posterior uncertainty becomes overconfident: highly correlated measurements contribute less independent information, yet a diagonal model treats them as additive evidence. Since optimal weighting depends on $\mathbf{R}$, correlation structure can materially affect both estimation and uncertainty quantification.

This subsection tests whether QPCA-EnDCF's whitening step provides robustness to correlated observation errors, and whether its advantages persist as correlation strength and conditioning increase. Importantly, we avoid substantial covariance misspecification: in every experiment, the filter is provided a covariance matrix $\mathbf{R}$ that matches the true error correlation structure. For severely ill-conditioned cases, a small diagonal regularization is applied to ensure numerical stability (see below), which represents a minor departure from exact specification rather than a modeling assumption. The question is therefore algorithmic: given known correlations, does whitening effectively decorrelate the errors and recover performance close to the diagonal baseline? For non-diagonal $\mathbf{R}$, whitening uses
\[
\mathbf{W} = (\mathbf{R}^{(L)})^{-1/2},
\qquad
\mathbf{W}\mathbf{R}^{(L)}\mathbf{W}^{\top} = \mathbf{I},
\]
a bijective linear change of coordinates that standardizes and decorrelates the observation errors. Since $\mathbf{R}^{(L)}=\mathbf{I}_L\otimes\mathbf{R}$, it suffices to whiten each per-observation block $\mathbf{R}$ independently. We consider the standard Cholesky whitening implementation: $\mathbf{R}=\mathbf{L}\mathbf{L}^{\top}$ with $\mathbf{W}=\mathbf{L}^{-\top}$ applied to each block. This approach has $\mathcal{O}(m^3)$ cost; the comparison here focuses on its numerical behavior under severe ill-conditioning. We test three covariance structures spanning five orders of magnitude in condition number. The diagonal baseline $\mathbf{R}=\sigma_{\mathrm{obs}}^2\mathbf{I}_m$ with $\sigma_{\mathrm{obs}}=1.5$ has $\mathrm{cond}(\mathbf{R})=1$. The exponential correlation model
\[
[\mathbf{R}]_{ij}=\sigma_{\mathrm{obs}}^2\exp\!\left(-\frac{d_{ij}}{\ell_{\mathrm{corr}}}\right),
\qquad
d_{ij}=\min(|i-j|,\,m-|i-j|),
\]
with $\ell_{\mathrm{corr}}=4$ yields moderate conditioning, $\mathrm{cond}(\mathbf{R})=649$. The Gaussian correlation model
\[
[\mathbf{R}]_{ij}=\sigma_{\mathrm{obs}}^2\exp\!\left(-\frac{d_{ij}^2}{2\ell_{\mathrm{corr}}^2}\right),
\]
with $\ell_{\mathrm{corr}}=4$ is severely ill-conditioned, $\mathrm{cond}(\mathbf{R})\approx 3.7\times10^5$. For this severely ill-conditioned case, a small diagonal regularization ($\varepsilon_{\mathrm{reg}}=10^{-3}$ times the trace) is added to both the data-generating covariance and the filter's $\mathbf{R}$, ensuring numerical stability while keeping the filter's covariance consistent with the true error distribution. Figure~\ref{fig:correlation_structures} visualizes the resulting patterns.

\begin{figure}[htbp]
\centering
\includegraphics[width=\textwidth]{figures/correlation_structures.png}
\caption{Representative observation error covariance structures $\mathbf{R}$ used in the correlation robustness study (left to right: diagonal baseline, exponential correlation, and Gaussian correlation), constructed using periodic index distances to avoid boundary effects. All cases share the same marginal variance $\sigma_{\mathrm{obs}}^2=2.25$; colors indicate covariance magnitude.}
\label{fig:correlation_structures}
\end{figure}

For each structure we compare QPCA-EnDCF with Cholesky whitening against 4D-EnKF as a stochastic baseline. All runs use a single realization with 50 assimilation windows. This choice reflects the substantial cost of repeated experiments under severely ill-conditioned covariances; the long window sequence still provides a meaningful within-run comparison. Figure~\ref{fig:reconstruction_errors} reports mean RMSE over the 50-window sequence for each correlation structure, along with the relative improvement of QPCA-EnDCF over 4D-EnKF.

\begin{figure}[htbp]
\centering
\includegraphics[width=\textwidth]{figures/reconstruction_errors.png}
\caption{Reconstruction error under correlated observation errors. Points show mean RMSE over 50 assimilation windows from a single realization for QPCA-EnDCF with Cholesky whitening (green) and 4D-EnKF (red) under three covariance structures: diagonal baseline ($\mathrm{cond}(\mathbf{R})=1$), exponential correlation ($\mathrm{cond}(\mathbf{R})=649$), and Gaussian correlation ($\mathrm{cond}(\mathbf{R})\approx 3.7\times10^5$). Annotated percentages indicate the RMSE reduction of QPCA-EnDCF relative to 4D-EnKF; the advantage grows with correlation strength, reaching approximately $32\%$ under severe ill-conditioning.}
\label{fig:reconstruction_errors}
\end{figure}

QPCA-EnDCF is robust across condition numbers spanning five orders of magnitude. Relative to the diagonal baseline, RMSE increases only modestly under correlation: from 3.47 (diagonal) to 3.71 (exponential, $+6.9\%$) and 3.65 (Gaussian, $+5.1\%$). Even under $\mathrm{cond}(\mathbf{R})\approx 3.7\times10^5$, performance remains within 7\% of the uncorrelated reference, indicating that whitening effectively neutralizes the impact of correlation on the subsequent spectral filtering. The Cholesky whitening implementation is numerically stable at the level relevant for filtering, yielding consistent RMSE results across all cases. Given its favorable constant factors, Cholesky whitening is the natural default when $\mathbf{R}$ is positive definite. By contrast, 4D-EnKF exhibits substantially greater degradation as correlations strengthen, with the most pronounced effect occurring in the severely ill-conditioned Gaussian case, where the RMSE increases from $4.64$ to $5.35$, corresponding to a $15.3\%$ deterioration. Consequently, QPCA-EnDCF's advantage increases with correlation strength: the reduction relative to 4D-EnKF is 25.2\% (diagonal), 24.7\% (exponential), and 31.9\% (Gaussian). Thus the whitening-based deterministic update is most beneficial precisely in regimes where correlations reduce the effective independent information content and stochastic perturbations become less efficient.

\FloatBarrier








% ================================================================================
% SECTION: Operational Considerations
% ================================================================================


\section{Operational Considerations}
\label{sec:operational}

The preceding sections demonstrate that QPCA-EnDCF improves reconstruction accuracy and calibration (Section~\ref{subsec:calibration}), removes the need for inflation (Section~\ref{sec:inflation}), and remains robust to non-ideal observation errors (Section~\ref{sec:robustness}). For deployment, the central question is how these advantages trade against computational constraints. In operational settings, ensemble size $N$ is often severely limited (typically $N \ll n$), while window length $L$ affects both information content and cost, with matrix operations scaling as $\mathcal{O}(m^3L^3)$. Practical implementation therefore requires characterizing performance degradation as $N$ and $L$ move from computationally generous to resource-constrained regimes, and comparing these trends across methods. This section studies sensitivity to the two primary computational parameters. Section~\ref{subsec:operational_ensemble_size} varies $N \in \{5,10,15,20,30,50,100\}$, spanning extreme undersampling to well-sampled regimes, and evaluates both reconstruction accuracy and calibration. Section~\ref{subsec:operational_window_length} varies $L \in \{1,3,5,7,10,15\}$ to identify regimes where temporal coupling provides net benefit. The goal is to establish minimum viable operating points and provide actionable guidance for parameter selection.

% ================================================================================
\subsection{Ensemble Size Dependence}
\label{subsec:operational_ensemble_size}

Ensemble size controls the fundamental accuracy--cost tradeoff: forecast propagation costs scale approximately linearly in $N$, while statistical error increases as sampling noise corrupts empirical covariances. This tradeoff differs across methods. Stochastic EnKF variants incur sampling error in $\widehat{\mathbf{P}}^{\,f}$ and additional variance from perturbed observations, whereas deterministic methods avoid the latter. The key operational question is whether removing perturbation noise yields meaningful gains in small-ensemble regimes. A useful baseline comes from random matrix theory: in the undersampled regime $N \ll n$, for i.i.d. samples from a distribution with covariance $\mathbf{P}^f$,
\[
\mathbb{E}\!\left[\|\widehat{\mathbf{P}}^{\,f}-\mathbf{P}^f\|_F^2\right] = \mathcal{O}\!\left(\frac{\bigl(\mathrm{tr}(\mathbf{P}^f)\bigr)^2 + \|\mathbf{P}^f\|_F^2}{N}\right),
\qquad
\widehat{\mathbf{P}}^{\,f}=\frac{1}{N-1}\mathbf{A}^f(\mathbf{A}^f)^{\top}.
\]
When $\mathbf{P}^f$ has entries of comparable magnitude, this scales as $\mathcal{O}(n^2/N)$; when normalized per component, it scales as $\mathcal{O}(n/N)$.
Stochastic filters further introduce variance from observation perturbations, scaling like $\mathcal{O}(\mathrm{tr}(\mathbf{R})/N)$ per update. QPCA-EnDCF eliminates this perturbation component; the following experiments quantify the practical impact. We vary $N \in \{5,10,15,20,30,50,100\}$ while holding the baseline configuration fixed ($n=40$, $m=20$, $L=5$, $\sigma_{\mathrm{obs}}=1.5$). This spans extreme undersampling ($N=5$, rank-$4$ covariances, $N/n=0.125$) through moderate undersampling ($N=20$, $N/n=0.5$) to well-sampled conditions ($N=100$, $N/n=2.5$). Each setting uses three independent trials with 50 assimilation windows. Figure~\ref{fig:performance_degradation} summarizes reconstruction accuracy. Panel (a) reports absolute RMSE versus $N$; panel (b) reports degradation relative to the $N=100$ baseline.

\begin{figure}[htbp]
\centering
\includegraphics[width=\textwidth]{figures/performance_degradation.png}
\caption{Reconstruction accuracy versus ensemble size $N \in \{5,10,15,20,30,50,100\}$. Curves show mean RMSE across three trials; shading indicates $\pm1$ standard deviation. \textbf{(A)} Absolute RMSE versus $N$ (logarithmic horizontal axis). \textbf{(B)} Relative degradation: percent increase in RMSE relative to $N=100$. QPCA-EnDCF degrades modestly from 3.33 ($N=100$) to 4.27 ($N=5$), whereas Sequential EnKF and 4D-EnKF degrade much more severely at small $N$.}
\label{fig:performance_degradation}
\end{figure}

QPCA-EnDCF exhibits bounded degradation as $N$ decreases, increasing from RMSE 3.33 at $N=100$ to 4.27 at $N=5$ (28.2\%). Stochastic methods deteriorate much more sharply in the severely undersampled regime: at $N=5$, Sequential EnKF and 4D-EnKF reach RMSE 4.71 and 4.86, respectively. This behavior is consistent with the additional variance introduced by observation perturbations, which becomes non-negligible when $N$ is small and is compounded by repeated updates (particularly for sequential assimilation).

Two operational patterns emerge. First, QPCA-EnDCF remains viable even under extreme undersampling: at $N=5$ it achieves RMSE $4.27\pm0.03$, improving on both stochastic methods. Second, the gap narrows with increasing $N$ as sampling error decreases for all methods; performance becomes statistically indistinguishable by $N=100$. This scaling yields corresponding computational savings. QPCA-EnDCF with $N=10$ attains an RMSE of $3.70$, statistically comparable to Sequential EnKF with $N=30$, which attains an RMSE of $3.62$, and improves upon 4D-EnKF with $N=20$, which attains an RMSE of $3.79$. (Values may differ slightly from baseline summaries due to independent trials and evaluation protocols.) Thus, for a fixed accuracy target, QPCA-EnDCF can reduce ensemble size by roughly $2$--$3\times$, cutting forecast propagation cost by 50--67\%. Reconstruction accuracy alone is not sufficient for deployment; reliable uncertainty quantification is equally important. We therefore extend the calibration diagnostics of Section~\ref{subsec:calibration} across $N \in \{5,10,15,20,30,50,100\}$ using spread--skill ratio $\bar{\gamma}$ (averaged over three trials and 50 windows per trial). Figure~\ref{fig:ensemble_calibration} reports $\bar{\gamma}$ versus $N$; the target $\bar{\gamma}=1$ indicates calibrated spread.

\begin{figure}[htbp]
\centering
\includegraphics[width=\textwidth]{figures/fig_ensemble_size_spread_skill.png}
\caption{Spread--skill ratio $\bar{\gamma}$ versus ensemble size $N \in \{5,10,15,20,30,50,100\}$. The dashed line marks the ideal target $\bar{\gamma}=1$. QPCA-EnDCF approaches ideal calibration as $N$ increases and is substantially closer to $\bar{\gamma}=1$ than the stochastic baselines at small and moderate ensemble sizes.}
\label{fig:ensemble_calibration}
\end{figure}

QPCA-EnDCF exhibits systematic calibration improvement with increasing $N$ ($\bar{\gamma}=0.38$ at $N=5$, $0.77$ at $N=10$, and $0.95$ at $N=50$), consistent with reduced covariance sampling error. Stochastic methods remain severely underdispersed throughout the tested range: even at $N=50$, $\bar{\gamma}\approx 0.12$--$0.14$, indicating persistent ensemble collapse. Operationally, this gap is decisive: QPCA-EnDCF at $N=10$ provides substantially better calibration than stochastic methods at $N=50$, implying comparable (or better) uncertainty quantification at roughly one-fifth the forecast propagation cost.

In summary, ensemble-size scaling shows that QPCA-EnDCF (i) maintains lower RMSE across all $N$ and (ii) achieves markedly better calibration in the small- to moderate-ensemble regime that dominates operational practice. These gains translate into concrete budget advantages: for a fixed accuracy or calibration target, QPCA-EnDCF can operate with substantially smaller ensembles than stochastic baselines.




% ================================================================================


\subsection{Window Length}
\label{subsec:operational_window_length}

Window length $L$ controls how many observation times are assimilated jointly, and therefore trades information content against computational and modeling costs. As $L$ grows, the number of assimilated observations increases linearly ($d$), which can improve estimation by exploiting temporal constraints. The costs are also clear: matrix sizes grow from $m\times m$ at $L=1$ to $d\times d$ for $L>1$, with dense linear algebra scaling as $\mathcal{O}(d^3)$. Longer windows propagate forecasts further, increasing susceptibility to accumulated model error, and once $L$ exceeds the decorrelation time, additional observations become largely redundant, yielding diminishing returns. For QPCA-EnDCF, $L$ directly changes the dimension of the whitened residual covariance $\mathbf{C}_E\in\mathbb{R}^{d\times d}$. Increasing $L$ provides a richer temporal residual structure from which principal modes can be identified, but gains should saturate beyond the decorrelation scale and may degrade if extra observations primarily contribute noise. In stochastic 4D-EnKF, longer windows can help through improved conditioning, but the impact is tempered by the accumulation of perturbation noise across multiple observation times. We vary $L\in\{1,3,5,7,10,15\}$ and adjust the number of windows $W$ so that the total number of assimilated observations satisfies $LW\approx 400$. This keeps the overall information budget approximately fixed, so differences reflect algorithmic response to temporal coupling rather than increased sample size. Each configuration uses three independent trials. Figure~\ref{fig:rmse_vs_L} reports reconstruction accuracy and the relative advantage of QPCA-EnDCF as $L$ varies. Panel (A) shows RMSE versus $L$ for all methods; panel (B) shows the percent improvement of QPCA-EnDCF relative to Sequential EnKF and 4D-EnKF.

\begin{figure}[htbp]
\centering
\includegraphics[width=\textwidth]{figures/combined_window_rmse_analysis.png}
\caption{Window-length sensitivity for $L\in\{1,3,5,7,10,15\}$ with $LW\approx 400$ held approximately constant. Shading denotes $\pm1$ standard deviation across three trials. \textbf{(A)} RMSE versus $L$ for Sequential EnKF (blue circles), 4D-EnKF (purple squares), and QPCA-EnDCF (orange triangles). \textbf{(B)} Percent RMSE improvement of QPCA-EnDCF relative to Sequential EnKF (green bars) and 4D-EnKF (orange bars) versus $L$.}
\label{fig:rmse_vs_L}
\end{figure}

QPCA-EnDCF shows a pronounced transition between $L=1$ and $L=3$. At $L=1$, RMSE is $4.87\pm0.10$, slightly worse than 4D-EnKF ($4.64\pm0.19$). At this window length, spectral regularization is based on a comparatively small residual covariance $\mathbf{C}_E\in\mathbb{R}^{m\times m}$, and the resulting eigenspectrum is not reliably informative. Increasing to $L=3$ reduces RMSE to $3.69\pm0.08$ (a 24.2\% decrease), indicating that joint assimilation over a short window is sufficient to enter a regime where dominant residual modes are stably identified and exploited. For $L\ge 5$, QPCA-EnDCF stabilizes (RMSE $\approx 3.64$--$3.68$ for $L\in\{5,7,10\}$) and increases only slightly at $L=15$. This plateau is consistent with diminishing returns once the window exceeds the temporal decorrelation scale: additional observations contribute limited independent information, while larger windows increase cost and may amplify model error or introduce weak spurious modes. Stochastic methods are far less sensitive to $L$. Sequential EnKF is essentially flat across the tested range (RMSE 4.58--4.72). 4D-EnKF improves modestly from $4.64$ at $L=1$ to $4.40$ at $L=15$ (5.1\%), suggesting that the benefit from temporal coupling is modest relative to the improvement achieved by QPCA-EnDCF over the range $L=1$ to $L=3$. Panel (B) emphasizes the regime change: QPCA-EnDCF is slightly worse than 4D-EnKF at $L=1$ (improvement $-5.1\%$), but for every $L\ge 3$ it outperforms 4D-EnKF by 16.3--21.0\%, with the largest gain at $L=3$. Thus, temporal coupling is essential for QPCA-EnDCF: once $L$ is large enough to expose coherent cross-time residual structure, spectral regularization delivers a substantial and stable advantage. Operationally, these results suggest (i) QPCA-EnDCF should be used with $L\ge 3$, and (ii) $L\in[5,10]$ provides a good balance between accuracy and cost. Conversely, if strict sequential processing is required ($L=1$), 4D-EnKF is competitive and may be preferable.

\FloatBarrier


% ================================================================================
\section{Conclusion}\label{sec:conclusion}
% ================================================================================

This work developed a theoretical framework for deterministic spectral regularization in ensemble data assimilation and demonstrated through controlled experiments that such methods achieve superior probabilistic calibration compared to conventional stochastic approaches. Three theoretical results characterize the mechanism: (i) deterministic spectral projection avoids variance injection from stochastic observation perturbations and leaves residual components in noise-dominated orthogonal directions unmodified; (ii) decomposition of estimation error yields a canonical $\mathcal{O}(1/N)$ sampling variance and identifies a $\kappa$- and $\delta_\kappa$-dependent stability prefactor associated with estimating the rank-$\kappa$ projector, which in low effective-dimension regimes can produce rank-driven sampling behavior relative to the observation dimension $d$; 
(iii) the convergence analysis shows projector-estimation error decays as $N$ grows (with effective dimension $\kappa$), and ensemble-size scaling experiments show calibration approaches ideal values, supporting deterministic spectral projection as a practical low-rank route to improved UQ under severe undersampling.

Numerical experiments on the Lorenz--96 system confirm theoretical predictions. QPCA-EnDCF achieves spread--skill ratios approaching unity at the baseline $N=10$ while stochastic methods remain severely underdispersed despite covariance inflation. The method requires no inflation, exhibits robustness to non-Gaussian errors and spatial correlations, and maintains advantages from extreme undersampling ($N = 5$) through well-sampled regimes. Geometric interpretation reveals the mechanism: spectral projection confines corrections to low-rank signal subspaces while preserving ensemble diversity exactly in orthogonal directions. This contrasts with stochastic perturbations where finite sampling induces spurious correlations, causing variance compression along intended preservation directions.

Operational studies establish practical deployment guidelines. 
QPCA-EnDCF at $N = 10$ achieves accuracy comparable to stochastic methods at $N = 20$--$30$ (2--3$\times$ savings) and calibration superior to stochastic methods at $N = 50$ (5$\times$ savings). 
Window length analysis reveals a sharp transition: sequential operation ($L = 1$) eliminates QPCA-EnDCF's advantage over stochastic baselines, but windowed operation with $L \geq 3$ restores 16--21\% accuracy improvements over 4D-EnKF. These findings establish minimum viable operating points and quantify computational savings in resource-constrained settings.

Several idealizations limit the generality of the present study. All experiments are conducted in a perfect-model setting, thereby excluding structural model error that is ubiquitous in operational data assimilation. The observing system is likewise simplified relative to real networks, which involve asynchronous data arrival, nonlinear observation operators, and heterogeneous (often correlated) error characteristics. In addition, the state dimension considered here ($n=40$) is far below operational regimes ($n\sim 10^6$--$10^9$), where both computational constraints and error structure can change qualitative behavior. Future work will therefore focus on assessing performance under controlled model error, developing adaptive procedures for selecting the truncation rank, validating the methodology on intermediate-complexity models, and evaluating scalability and robustness in operational-scale settings.
Determining whether the calibration and efficiency gains observed in idealized chaotic systems persist at operational scales, and whether analogous benefits arise in other high-dimensional inference problems, remains an important open direction.


\FloatBarrier




% ================================================================================
% APPENDICES
% ================================================================================

\appendix

% ================================================================================
% APPENDIX: Algorithm Specifications
% ================================================================================


\section{Algorithm Specifications}\label{app:algorithms}

We provide complete algorithmic specifications for the three ensemble data assimilation methods compared in this study: Sequential Ensemble Kalman Filter, Four-Dimensional Ensemble Kalman Filter, and QPCA Ensemble Data Consistency Filter.
% ================================================================================

\subsection{Sequential Stochastic Ensemble Kalman Filter}\label{app:seq_enkf}

The sequential stochastic ensemble Kalman filter processes observations individually, updating the ensemble through perturbed Kalman corrections at each observation time.

\begin{algorithm}[H]
\caption{Sequential Stochastic Ensemble Kalman Filter}
\label{alg:seq_enkf_appendix}
\begin{algorithmic}[1]
\Require Initial ensemble $\mathbf{X}_0 \in \mathbb{R}^{n \times N}$, observation operator $\mathbf{H}$, observation error covariance $\mathbf{R}$, observation sequence $\{\mathbf{z}_k\}_{k=1}^K$, forward model $\mathcal{M}$, multiplicative inflation factor $\lambda_{\mathrm{infl}} \geq 1$
\Ensure Analysis ensemble $\mathbf{X}_K^a$
\State $\mathbf{X} \gets \mathbf{X}_0$
\For{$k = 1$ to $K$}
    \For{$j = 1$ to $N$}
        \State $\mathbf{x}^{(j)} \gets \mathcal{M}(\mathbf{x}^{(j)})$ \Comment{Forecast propagation}
    \EndFor
    \State $\bar{\mathbf{x}}^f \gets \frac{1}{N}\sum_{j=1}^{N}\mathbf{x}^{(j)}$, \quad $\mathbf{A}^f \gets \mathbf{X} - \bar{\mathbf{x}}^f\mathbf{1}^{\top}$ \Comment{Forecast anomalies}
    \State $\widehat{\mathbf{P}}^{\,f} \gets \frac{1}{N-1}\mathbf{A}^f(\mathbf{A}^f)^{\top}$, \quad $\mathbf{Y} \gets \mathbf{H}\mathbf{X}$
    \State $\bar{\mathbf{y}} \gets \frac{1}{N}\mathbf{Y}\mathbf{1}$ \Comment{Mean predicted observation}
    \State $\mathbf{S} \gets \frac{1}{N-1}(\mathbf{Y} - \bar{\mathbf{y}}\mathbf{1}^{\top})(\mathbf{Y} - \bar{\mathbf{y}}\mathbf{1}^{\top})^{\top} + \mathbf{R}$
    \State $\mathbf{K}_k \gets \widehat{\mathbf{P}}^{\,f}\mathbf{H}^{\top}\mathbf{S}^{-1}$
    \For{$j = 1$ to $N$}
        \State $\boldsymbol{\epsilon}_{k}^{(j)} \sim \mathcal{N}(\mathbf{0}, \mathbf{R})$
        \State $\mathbf{x}^{(j)} \gets \mathbf{x}^{(j)} + \mathbf{K}_k(\mathbf{z}_k + \boldsymbol{\epsilon}_{k}^{(j)} - \mathbf{H}\mathbf{x}^{(j)})$ \Comment{Perturbed update}
    \EndFor
    \State $\bar{\mathbf{x}}^a \gets \frac{1}{N}\sum_{j=1}^{N}\mathbf{x}^{(j)}$
    \State $\mathbf{x}^{(j)} \gets \bar{\mathbf{x}}^a + \lambda_{\mathrm{infl}}(\mathbf{x}^{(j)} - \bar{\mathbf{x}}^a)$ for $j=1,\ldots,N$ \Comment{Post-update inflation}
\EndFor
\State \Return $\mathbf{X}$
\end{algorithmic}
\end{algorithm}

Observation perturbations $\boldsymbol{\epsilon}_{k}^{(j)} \sim \mathcal{N}(\mathbf{0}, \mathbf{R})$ are introduced to preserve ensemble variance in expectation. Multiplicative inflation $\lambda_{\mathrm{infl}} > 1$ is applied after each analysis update, inflating analysis anomalies about the analysis mean to counteract ensemble collapse.
% ================================================================================

\subsection{Four-Dimensional Stochastic Ensemble Kalman Filter}\label{app:4d_enkf}

The four-dimensional stochastic ensemble Kalman filter accumulates information from $L$ consecutive observations within temporal windows, performing joint analysis updates.

\begin{algorithm}[H]
\caption{Four-Dimensional Stochastic Ensemble Kalman Filter}
\label{alg:4d_enkf_appendix}
\begin{algorithmic}[1]
\Require Initial ensemble $\mathbf{X}_0 \in \mathbb{R}^{n \times N}$, observation operator $\mathbf{H}$, observation error covariance $\mathbf{R}$, window length $L$, number of windows $W$, observation sequence $\{\mathbf{z}_k\}_{k=1}^{K}$, forward model $\mathcal{M}$, multiplicative inflation factor $\lambda_{\mathrm{infl}} \geq 1$
\Ensure Analysis ensemble at window ends
\State $\mathbf{X} \gets \mathbf{X}_0$, \quad $\mathbf{R}^{(L)} \gets \mathbf{I}_L \otimes \mathbf{R}$
\For{$w = 1$ to $W$}
    \State Initialize $\{\mathbf{X}_\ell\}_{\ell=0}^{L}$ with $\mathbf{X}_0 \gets \mathbf{X}$
    \For{$\ell = 1$ to $L$}
        \For{$j = 1$ to $N$}
            \State $\mathbf{x}^{(j)} \gets \mathcal{M}(\mathbf{x}^{(j)})$
        \EndFor
        \State $\mathbf{X}_\ell \gets \mathbf{X}$, \quad $\mathbf{Y}_\ell \gets \mathbf{H}\mathbf{X}_\ell$
    \EndFor
    \State $\mathbf{z}^{(w)} \gets [\mathbf{z}_{(w-1)L+1}^{\top}, \ldots, \mathbf{z}_{wL}^{\top}]^{\top}$, \quad $\mathbf{Y}_{\mathrm{stack}} \gets [\mathbf{Y}_1^{\top}, \ldots, \mathbf{Y}_L^{\top}]^{\top}$
    \State $\bar{\mathbf{x}}_L^f \gets \frac{1}{N}\mathbf{X}_L\mathbf{1}$, \quad $\bar{\mathbf{y}}_{\mathrm{stack}} \gets \frac{1}{N}\mathbf{Y}_{\mathrm{stack}}\mathbf{1}$ \Comment{Ensemble means}
    \State $\mathbf{A}_x \gets \mathbf{X}_L - \bar{\mathbf{x}}_L^f\mathbf{1}^{\top}$, \quad $\mathbf{A}_z \gets \mathbf{Y}_{\mathrm{stack}} - \bar{\mathbf{y}}_{\mathrm{stack}}\mathbf{1}^{\top}$ \Comment{Forecast anomalies}
    \State $\mathbf{P}_{xz} \gets \frac{1}{N-1}\mathbf{A}_x\mathbf{A}_z^{\top}$, \quad $\mathbf{P}_{zz} \gets \frac{1}{N-1}\mathbf{A}_z\mathbf{A}_z^{\top}$
    \State $\mathbf{K}^{(w)} \gets \mathbf{P}_{xz}(\mathbf{P}_{zz} + \mathbf{R}^{(L)})^{-1}$
    \For{$j = 1$ to $N$}
        \State $\boldsymbol{\epsilon}_{w}^{(j)} \sim \mathcal{N}(\mathbf{0}, \mathbf{R}^{(L)})$
        \State $\mathbf{x}_L^{(j)} \gets \mathbf{x}_L^{(j)} + \mathbf{K}^{(w)}(\mathbf{z}^{(w)} + \boldsymbol{\epsilon}_{w}^{(j)} - \mathbf{Y}_{\mathrm{stack}}^{(j)})$
    \EndFor
    \State $\mathbf{X} \gets \mathbf{X}_L$
    \State $\bar{\mathbf{x}}^a \gets \frac{1}{N}\mathbf{X}\mathbf{1}$
    \State $\mathbf{x}^{(j)} \gets \bar{\mathbf{x}}^a + \lambda_{\mathrm{infl}}(\mathbf{x}^{(j)} - \bar{\mathbf{x}}^a)$ for $j=1,\ldots,N$ \Comment{Post-update inflation}
\EndFor
\State \Return $\mathbf{X}$
\end{algorithmic}
\end{algorithm}

Joint processing of $L$ observations exploits temporal correlations through extended cross-covariance $\mathbf{P}_{xz} \in \mathbb{R}^{n \times d}$. Perturbations are sampled from $\mathcal{N}(\mathbf{0}, \mathbf{R}^{(L)})$ with block-diagonal covariance structure. Multiplicative inflation is applied after the analysis update, consistent with the sequential variant.
% ================================================================================

\subsection{QPCA Ensemble Data Consistency Filter}\label{app:qpca_endcf}

The QPCA ensemble data consistency filter implements deterministic residual filtering through whitening, spectral decomposition, and low-rank projection.

\begin{algorithm}[H]
\caption{QPCA Ensemble Data Consistency Filter}
\label{alg:qpca_endcf}
\begin{algorithmic}[1]
\Require Initial ensemble $\mathbf{X}_0 \in \mathbb{R}^{n \times N}$, observation operator $\mathbf{H}$, observation error covariance $\mathbf{R}$, window length $L$, number of windows $W$, observation sequence $\{\mathbf{z}_k\}_{k=1}^{K}$, forward model $\mathcal{M}$, spectral truncation $\kappa$
\Ensure Analysis ensemble at window ends
\State $\mathbf{X} \gets \mathbf{X}_0$, \quad $\mathbf{R}^{(L)} \gets \mathbf{I}_L \otimes \mathbf{R}$
\For{$w = 1$ to $W$}
    \State \textit{// Forecast phase}
    \State Initialize $\{\mathbf{X}_\ell\}_{\ell=0}^{L}$ with $\mathbf{X}_0 \gets \mathbf{X}$
    \For{$\ell = 1$ to $L$}
        \For{$j = 1$ to $N$}
            \State $\mathbf{x}^{(j)} \gets \mathcal{M}(\mathbf{x}^{(j)})$
        \EndFor
        \State $\mathbf{X}_\ell \gets \mathbf{X}$, \quad $\mathbf{Y}_\ell \gets \mathbf{H}\mathbf{X}_\ell$
    \EndFor
    \State
    \State \textit{// Whitened residual computation}
    \State $\mathbf{z}^{(w)} \gets [\mathbf{z}_{(w-1)L+1}^{\top}, \ldots, \mathbf{z}_{wL}^{\top}]^{\top}$, \quad $\mathbf{Y}_{\mathrm{stack}} \gets [\mathbf{Y}_1^{\top}, \ldots, \mathbf{Y}_L^{\top}]^{\top}$
    \State $\bar{\mathbf{x}}_L^f \gets \frac{1}{N}\mathbf{X}_L\mathbf{1}$, \quad $\bar{\mathbf{y}}_{\mathrm{stack}} \gets \frac{1}{N}\mathbf{Y}_{\mathrm{stack}}\mathbf{1}$ \Comment{Ensemble means}
    \State $\mathbf{D} \gets \mathbf{Y}_{\mathrm{stack}} - \mathbf{z}^{(w)}\mathbf{1}^{\top}$
    \State $\mathbf{E} \gets (\mathbf{R}^{(L)})^{-1/2}\mathbf{D}$ \Comment{Whitening}
    \State $\mathbf{E}_c \gets \mathbf{E} - \frac{1}{N}(\mathbf{E}\mathbf{1})\mathbf{1}^{\top}$ \Comment{Centering}
    \State
	    \State \textit{// Spectral decomposition}
	    \State $\mathbf{C}_E \gets \frac{1}{N-1}\mathbf{E}_c\mathbf{E}_c^{\top}$
	    \State $[\hat{\mathbf{V}}, \hat{\boldsymbol{\Lambda}}] \gets \mathrm{eig}(\mathbf{C}_E)$ \Comment{$\hat{\mathbf{V}}=[\hat{\mathbf{v}}_1,\ldots,\hat{\mathbf{v}}_d]$, $\hat{\boldsymbol{\Lambda}}=\mathrm{diag}(\hat{\lambda}_1,\ldots,\hat{\lambda}_d)$}
	    \State $r \gets \min(d, N-1)$ \Comment{Maximum rank of $\mathbf{C}_E$, where $d=mL$}
	    \State Sort indices so $\hat{\lambda}_1 \geq \hat{\lambda}_2 \geq \cdots \geq \hat{\lambda}_{r} \geq 0$; permute $\hat{\mathbf{v}}_i$ consistently
	    \State $\hat{\mathbf{V}}_{\kappa} \gets [\hat{\mathbf{v}}_1, \ldots, \hat{\mathbf{v}}_{\kappa}]$ \Comment{Leading $\kappa$ eigenvectors}
	    \State
	    \State \textit{// Regularized correction}
	    \State $\mathbf{Q}_{\mathrm{PCA}} \gets -\hat{\mathbf{V}}_{\kappa}\hat{\mathbf{V}}_{\kappa}^{\top}\mathbf{E}$
	    \State $\boldsymbol{\Delta}_{\mathrm{obs}} \gets (\mathbf{R}^{(L)})^{1/2}\mathbf{Q}_{\mathrm{PCA}}$
    \State
    \State \textit{// Update}
	    \State $\mathbf{A}_x \gets \mathbf{X}_L - \bar{\mathbf{x}}_L^f\mathbf{1}^{\top}$
	    \State $\mathbf{A}_z \gets \mathbf{Y}_{\mathrm{stack}} - \bar{\mathbf{y}}_{\mathrm{stack}}\mathbf{1}^{\top}$
	    \State $\mathbf{P}_{xz} \gets \frac{1}{N-1}\mathbf{A}_x\mathbf{A}_z^{\top}$, \quad $\mathbf{P}_{zz} \gets \frac{1}{N-1}\mathbf{A}_z\mathbf{A}_z^{\top}$
	    \State $\mathbf{K}^{\mathrm{DC}} \gets \mathbf{P}_{xz}(\mathbf{P}_{zz})^{\dagger}$
	    \State $\mathbf{X}_L \gets \mathbf{X}_L + \mathbf{K}^{\mathrm{DC}}\boldsymbol{\Delta}_{\mathrm{obs}}$
    \State $\mathbf{X} \gets \mathbf{X}_L$
\EndFor
\State \Return $\mathbf{X}$
\end{algorithmic}
\end{algorithm}

\subsubsection{Implementation Notes}

The whitening transformation $\mathbf{E} \gets (\mathbf{R}^{(L)})^{-1/2}\mathbf{D}$ normalizes residuals so observation errors have identity covariance in transformed space. The algorithm accepts general $\mathbf{R}$. For the diagonal special case $\mathbf{R}=\sigma_{\mathrm{obs}}^2\mathbf{I}_m$, this reduces to scalar division by $\sigma_{\mathrm{obs}}$ and can be implemented efficiently. For correlated errors, matrix square root requires Cholesky decomposition or symmetric eigendecomposition.

Spectral truncation implements hard thresholding: $\mathbf{Q}_{\mathrm{PCA}} = -\hat{\mathbf{V}}_{\kappa}\hat{\mathbf{V}}_{\kappa}^{\top}\mathbf{E}$ retains only leading $\kappa$ eigenmodes. The negative sign ensures corrections shift ensemble toward observations.

The algorithm applies a deterministic correction to the ensemble without stochastic observation perturbations. As written, $\boldsymbol{\Delta}_{\mathrm{obs}}$ is a matrix of memberwise observation-space increments derived from the residual matrix, so corrections are deterministic but generally member dependent. State anomalies are computed without inflation factor, in contrast to explicit inflation in stochastic methods.

\FloatBarrier


\section{Numerical Experiment Details}\label{app:experiment_details}

\subsection{Rank Histogram Analysis}\label{subsec:rank_histograms}


For each assimilation cycle $k$ and state component $i$, we compute the rank of the truth within the analysis ensemble,
\begin{equation}
r_{k,i}
=
\left|\{j : x_{k,i}^{(j),a} < x_{k,i}^{\mathrm{true}}\}\right| + 1,
\label{eq:rank}
\end{equation}
where $x_{k,i}^{(j),a}$ is the $i$th component of analysis ensemble member $j$. Ties (which occur with probability zero under continuous distributions and are empirically rare) are broken randomly. Under perfect calibration, ranks are uniformly distributed on $\{1,\ldots,N+1\}$ with probability $1/(N+1)$. Systematic departures indicate miscalibration: dome-shaped histograms reflect overdispersion, U-shaped histograms indicate underdispersion, and asymmetry signals bias.

Uniformity is assessed using a chi-squared goodness-of-fit statistic,
\begin{equation}
\chi^2 = \sum_{b=1}^{N+1}\frac{(O_b - E_b)^2}{E_b},
\label{eq:chi_squared}
\end{equation}
where $O_b$ denotes the observed count in bin $b$ and $E_b = M_{\mathrm{rank}}/(N+1)$ is the expected count under uniformity, with $M_{\mathrm{rank}}$ being the total number of rank samples. Because large sample sizes make $\chi^2$ highly sensitive to minor deviations, we also report a flatness metric
\begin{equation}
\text{Flatness} = \frac{\sqrt{\frac{1}{N+1}\sum_{b=1}^{N+1}(f_b - \bar{f})^2}}{\bar{f}},
\label{eq:flatness}
\end{equation}
where $f_b = O_b / M_{\mathrm{rank}}$ is the observed frequency of bin $b$ and $\bar{f} = 1/(N+1)$ is the expected uniform frequency. This normalized standard deviation quantifies the practical magnitude of nonuniformity: perfectly uniform histograms yield Flatness $= 0$, while severe departures produce values $\gg 1$.

Ranks are computed at each algorithm's analysis times. Aggregating across all state components ($n=40$) and trials (5) yields $M_{\mathrm{rank}}=n \times K \times 5 = 50{,}000$ rank samples for the Sequential EnKF (evaluated at all $K=WL=250$ observation times), and $M_{\mathrm{rank}}=n \times W \times 5 = 10{,}000$ samples for the windowed methods (evaluated at the $W=50$ window endpoints). For $N=10$, there are $N+1=11$ bins, yielding an expected count of $E_b=M_{\mathrm{rank}}/11$ per bin under uniformity.

\begin{figure}[htbp]
\centering
\includegraphics[width=\textwidth]{figures/fig04d_rank_histograms.png}
\caption{Rank histogram analysis aggregated over all state components, analysis times, and trials. The dashed line indicates the expected uniform count $E_b=M_{\mathrm{rank}}/(N+1)$ in each panel (here $E_b\approx 4545$ for Seq-EnKF and $E_b\approx 909$ for 4D-EnKF and QPCA-EnDCF, corresponding to uniform relative frequency $1/11$), and the shaded band denotes the 95\% confidence interval under the uniformity null. Left: Sequential EnKF exhibits severe U-shaped distortion ($\chi^2=173{,}195$, flatness $1.861$). Center: 4D-EnKF shows substantial but reduced departure from uniformity ($\chi^2=32{,}269$, flatness $1.796$). Right: QPCA-EnDCF shows near-uniformity ($\chi^2=396.8$, flatness $0.199$).}
\label{fig:rank_histograms}
\end{figure}

QPCA-EnDCF provides the closest approximation to uniformity, with a flatness metric of $0.199$ and rank frequencies fluctuating mildly around the expected value. Although $\chi^2=396.8$ formally rejects uniformity ($p<0.001$), this reflects the extreme power of the test with $10{,}000$ samples. In relative terms, the departure is nearly two orders of magnitude smaller than for 4D-EnKF and three orders of magnitude smaller than for the Sequential EnKF, indicating near-distributional consistency in practice.

Sequential EnKF is strongly miscalibrated ($\chi^2=173{,}195$, flatness $1.861$), with a pronounced U shape: extreme ranks occur far more frequently than expected, while central ranks are depleted. This is the canonical signature of severe underdispersion caused by ensemble collapse. The 4D-EnKF shows improvement ($\chi^2=32{,}269$, flatness $1.796$) but retains a U-shaped structure, indicating persistent underdispersion despite temporal coupling. While increasing $L$ provides more observational information per window, it simultaneously increases the effective observation dimension $mL$, exacerbating undersampling when $N \ll mL$. For 4D-EnKF with $N = 10$ and $mL = 100$, the latter effect dominates.

\FloatBarrier




% % =============================================================================
% % APPENDIX: Notation and Dimension Reference
% % =============================================================================
% \section{Notation and Dimension Reference}
% \label{sec:notation_reference}

% This section provides a consolidated reference for the mathematical objects and dimensions used in the analysis of QPCA-EnDCF.
% % We distinguish population quantities (deterministic objects defined from the forecast distribution) from their \emph{sample}
% % counterparts (random objects computed from a finite ensemble). In particular, the EnDCF gain used in implementation is
% % \emph{empirical} and therefore random; throughout the analysis we either work conditionally on the ensemble or impose explicit
% % moment assumptions to justify bounds in expectation. Moreover, in the small-ensemble regime $(N-1)<mL$ we define the gain using a
% % Moore--Penrose pseudoinverse, so the gain is always well-defined.

% % -----------------------------------------------------------------------------

% % ================================================================================

% \subsection{Scalar Parameters and Dimensions}

% \begin{table}[H]
% \centering
% \begin{tabular}{cl}
% \toprule
% \textbf{Symbol} & \textbf{Description} \\
% \midrule
% $n$ & State dimension (Lorenz--96: $n=40$) \\
% $m$ & Observation dimension at a single time ($m=20$) \\
% $L$ & Assimilation window length (typically $L=5$) \\
% $d$ & Stacked observation dimension, $d:=mL$ \\
% $N$ & Ensemble size (typically $N=10$--$15$) \\
% $r$ & Maximum sample-covariance rank, $r:=\min(d,N-1)$ \\
% $\kappa$ & Spectral truncation level ($1\le \kappa \le r$) \\
% $r_\Sigma$ & Rank of population covariance, $r_\Sigma:=\mathrm{rank}(\boldsymbol{\Sigma}_E)\le d$ \\
% $r_C$ & Rank of sample covariance, $r_C:=\mathrm{rank}(\mathbf{C}_E)\le \min(r_\Sigma, N-1)$ \\
% $\delta_\kappa$ & Cutoff spectral gap, $\delta_\kappa:=\lambda_\kappa-\lambda_{\kappa+1}$ \\
% $F$ & Forcing parameter for Lorenz--96 dynamics ($F=8.0$) \\
% $\sigma_{\mathrm{obs}}$ & Observation-noise standard deviation ($\sigma_{\mathrm{obs}}=1.5$) \\
% $\tau_{\mathrm{Lyap}}$ & Lyapunov timescale ($\tau_{\mathrm{Lyap}}\approx 0.6$ time units) \\
% $T_{\mathrm{obs}}$ & Time interval between observations ($T_{\mathrm{obs}}=0.1$ time units) \\
% \bottomrule
% \end{tabular}
% \end{table}

% % -----------------------------------------------------------------------------

% % ================================================================================

% \subsection{State and Observation Vectors}

% \begin{table}[H]
% \centering
% \small
% \begin{tabular}{ccp{7.5cm}}
% \toprule
% \textbf{Symbol} & \textbf{Dimension} & \textbf{Description} \\
% \midrule
% $\mathbf{x}$ & $\mathbb{R}^n$ & State vector \\
% $\mathbf{x}^{(j)}$ & $\mathbb{R}^n$ & Ensemble member $j$ \\
% $\mathbf{x}^{(j),f}$ & $\mathbb{R}^n$ & Forecast ensemble member $j$ \\
% $\mathbf{x}^{(j),a}$ & $\mathbb{R}^n$ & Analysis ensemble member $j$ \\
% $\bar{\mathbf{x}}$ & $\mathbb{R}^n$ & Ensemble mean, $\bar{\mathbf{x}}=\frac{1}{N}\sum_{j=1}^N \mathbf{x}^{(j)}$ \\
% $\bar{\mathbf{x}}^{f}$ & $\mathbb{R}^n$ & Forecast ensemble mean \\
% $\bar{\mathbf{x}}^{a}$ & $\mathbb{R}^n$ & Analysis ensemble mean \\
% $\boldsymbol{\mu}^f$ & $\mathbb{R}^n$ & Population forecast mean \\
% $\mathbf{x}^{\mathrm{true}}$ & $\mathbb{R}^n$ & True state (time subscript $\mathbf{x}^{\mathrm{true}}_k$ or $\mathbf{x}_{t_k}$ used when emphasizing time dependence) \\
% $\mathbf{z}_k$ & $\mathbb{R}^m$ & Observation at time $k$ (global time index) \\
% 	$\mathbf{z}^{(w)}$ & $\mathbb{R}^{d}$ & Stacked observations over window $w$, $\mathbf{z}^{(w)}=[\mathbf{z}_{(w-1)L+1}^\top,\ldots,\mathbf{z}_{wL}^\top]^\top$ \\
% 	$\boldsymbol{\eta}_k$ & $\mathbb{R}^m$ & Observation noise (physical), $\boldsymbol{\eta}_k\sim \mathcal{N}(\mathbf{0},\mathbf{R})$ \\
% 	$\boldsymbol{\epsilon}_{k}^{(j)}$ & $\mathbb{R}^{m}$ & \textbf{Algorithmic} perturbed observation for member $j$ at time $k$, $\boldsymbol{\epsilon}_{k}^{(j)}\sim \mathcal{N}(\mathbf{0},\mathbf{R})$ \\
% 	$\boldsymbol{\epsilon}_{w}^{(j)}$ & $\mathbb{R}^{d}$ & \textbf{Algorithmic} perturbed observation for member $j$ in window $w$, $\boldsymbol{\epsilon}_{w}^{(j)}\sim \mathcal{N}(\mathbf{0},\mathbf{R}^{(L)})$ \\
% 	$\mathbf{e}^{(j)}$ & $\mathbb{R}^{d}$ & Whitened residual for member $j$ \\
% 	$\bar{\mathbf{e}}$ & $\mathbb{R}^{d}$ & Sample mean of whitened residuals \\
% 	$\boldsymbol{\mu}_E$ & $\mathbb{R}^{d}$ & Population mean of whitened residuals \\
% 	$\boldsymbol{\mu}_E^{\,\ell}$ & $\mathbb{R}^{d}$ & Population mean with approximate QoI map $Q_\ell$ \\
% $\mathbf{v}_i$ & $\mathbb{R}^{d}$ & Population eigenvector $i$ of $\boldsymbol{\Sigma}_E$ \\
% $\hat{\mathbf{v}}_i$ & $\mathbb{R}^{d}$ & Sample eigenvector $i$ of $\mathbf{C}_E$ \\
% $\mathbf{1}$ & $\mathbb{R}^N$ & Vector of ones \\
% \bottomrule
% \end{tabular}
% \end{table}

% \begin{note}
% \textbf{Distinction between $\boldsymbol{\eta}_k$ and $\boldsymbol{\epsilon}_k^{(j)}$:}
% $\boldsymbol{\eta}_k$ denotes the \emph{physical} observation noise present in the data-generating process, whereas $\boldsymbol{\epsilon}_k^{(j)}$ and $\boldsymbol{\epsilon}_w^{(j)}$ denote \emph{artificial} observation perturbations introduced algorithmically in stochastic ensemble Kalman filters to preserve ensemble variance. Both are drawn from the same distribution, but serve fundamentally different purposes.
% \end{note}

% % -----------------------------------------------------------------------------

% % ================================================================================

\subsection{Matrices and Operators}

\begin{table}[H]
\centering
\footnotesize
\resizebox{\textwidth}{!}{%
\begin{tabular}{ccl}
\toprule
\textbf{Symbol} & \textbf{Dimension} & \textbf{Description} \\
\midrule
$\mathbf{X}$ & $\mathbb{R}^{n\times N}$ & Ensemble matrix, $\mathbf{X}=[\mathbf{x}^{(1)},\ldots,\mathbf{x}^{(N)}]$ \\
$\mathbf{A}$ & $\mathbb{R}^{n\times N}$ & Anomaly matrix, $\mathbf{A}=\mathbf{X}-\bar{\mathbf{x}}\mathbf{1}^\top$ \\
$\mathbf{H}$ & $\mathbb{R}^{m\times n}$ & Observation operator at a single time \\
$\mathbf{H}^{(L)}$ & $\mathbb{R}^{d\times n}$ & Stacked observation operator over a window \\
$\mathbf{R}$ & $\mathbb{R}^{m\times m}$ & Observation-error covariance (single time), $\mathbf{R}=\sigma_{\mathrm{obs}}^2\mathbf{I}_m$ \\
$\mathbf{R}^{(L)}$ & $\mathbb{R}^{d\times d}$ & Stacked observation-error covariance, $\mathbf{R}^{(L)}=\mathbf{I}_L\otimes \mathbf{R}$ \\
$\mathbf{P}^f$ & $\mathbb{R}^{n\times n}$ & \textbf{Population} forecast covariance (deterministic) \\
$\widehat{\mathbf{P}}^{\,f}$ & $\mathbb{R}^{n\times n}$ & \textbf{Sample} forecast covariance, $\widehat{\mathbf{P}}^{\,f}=\frac{1}{N-1}\mathbf{A}\mathbf{A}^\top$ \\
$\mathbf{P}^a$ & $\mathbb{R}^{n\times n}$ & Population analysis covariance \\
$\widehat{\mathbf{P}}^{\,a}$ & $\mathbb{R}^{n\times n}$ & Sample analysis covariance \\
$\mathbf{E}$ & $\mathbb{R}^{d\times N}$ & Whitened residual matrix, $\mathbf{E}=[\mathbf{e}^{(1)},\ldots,\mathbf{e}^{(N)}]$ \\
$\mathbf{C}_E$ & $\mathbb{R}^{d\times d}$ & Sample covariance of whitened residuals \\
$\boldsymbol{\Sigma}_E$ & $\mathbb{R}^{d\times d}$ & Population covariance of whitened residuals \\
$\boldsymbol{\Sigma}_E^{\,\ell}$ & $\mathbb{R}^{d\times d}$ & Population covariance under approximate map $Q_\ell$ \\
$\boldsymbol{\Sigma}_E^\kappa$ & $\mathbb{R}^{d\times d}$ & Rank-$\kappa$ truncation of $\boldsymbol{\Sigma}_E$ \\
$\mathbf{C}_E^\kappa$ & $\mathbb{R}^{d\times d}$ & Rank-$\kappa$ truncation of $\mathbf{C}_E$ \\
$\mathbf{P}_\kappa$ & $\mathbb{R}^{d\times d}$ & Population rank-$\kappa$ projector, $\mathbf{P}_\kappa=\sum_{i=1}^\kappa \mathbf{v}_i\mathbf{v}_i^\top$ \\
$\hat{\mathbf{P}}_\kappa$ & $\mathbb{R}^{d\times d}$ & Sample rank-$\kappa$ projector, $\hat{\mathbf{P}}_\kappa=\sum_{i=1}^\kappa \hat{\mathbf{v}}_i\hat{\mathbf{v}}_i^\top$ \\
$\mathbf{K}_k$ & $\mathbb{R}^{n\times m}$ & Sequential Kalman gain at time $k$ (using sample covariance $\widehat{\mathbf{P}}^{\,f}$) \\
$\mathbf{K}^{(w)}$ & $\mathbb{R}^{n\times d}$ & Windowed Kalman gain for window $w$ ($d=mL$, from sample covariances) \\
$\mathbf{K}^{\mathrm{DC}}$ & $\mathbb{R}^{n\times d}$ & \textbf{Empirical} EnDCF gain (random), defined via a pseudoinverse \\
$\mathbf{X}_L$ & $\mathbb{R}^{n\times N}$ & Forecast ensemble at assimilation time / window endpoint \\
$\mathbf{Y}_{\mathrm{stack}}$ & $\mathbb{R}^{d\times N}$ & Stacked forecast observations, columns $Q(\mathbf{x}^{(j),f})=\mathbf{H}^{(L)}\mathbf{x}^{(j),f}$ \\
$\mathbf{A}_x$ & $\mathbb{R}^{n\times N}$ & State anomalies at endpoint, $\mathbf{A}_x=\mathbf{X}_L-\bar{\mathbf{x}}_L^f\mathbf{1}^\top$ \\
$\mathbf{A}_z$ & $\mathbb{R}^{d\times N}$ & Obs anomalies, $\mathbf{A}_z=\mathbf{Y}_{\mathrm{stack}}-\bar{\mathbf{y}}_{\mathrm{stack}}\mathbf{1}^\top$ \\
$\mathbf{P}_{xz}$ & $\mathbb{R}^{n\times d}$ & Cross-covariance, $\mathbf{P}_{xz}=\frac{1}{N-1}\mathbf{A}_x\mathbf{A}_z^\top$ \\
$\mathbf{P}_{zz}$ & $\mathbb{R}^{d\times d}$ & Auto-covariance, $\mathbf{P}_{zz}=\frac{1}{N-1}\mathbf{A}_z\mathbf{A}_z^\top$ \\
$\hat{\mathbf{V}}_\kappa$ & $\mathbb{R}^{d\times \kappa}$ & Matrix of first $\kappa$ sample eigenvectors, $\hat{\mathbf{V}}_\kappa=[\hat{\mathbf{v}}_1,\ldots,\hat{\mathbf{v}}_\kappa]$ \\
\bottomrule
\end{tabular}%
}
\end{table}

% % -----------------------------------------------------------------------------

% % ================================================================================

% \subsection{Spectral Quantities}

% \begin{table}[H]
% \centering
% \begin{tabular}{ccl}
% \toprule
% \textbf{Symbol} & \textbf{Type} & \textbf{Description} \\
% \midrule
% $\lambda_i$ & Scalar & Population eigenvalue $i$ of $\boldsymbol{\Sigma}_E$ (deterministic) \\
% & & Ordered: $\lambda_1\ge \lambda_2\ge \cdots \ge \lambda_{d}\ge 0$, with $\lambda_i=0$ for $i>r_\Sigma$ \\
% $\hat{\lambda}_i$ & Scalar & Sample eigenvalue $i$ of $\mathbf{C}_E$ (random) \\
% & & Ordered: $\hat{\lambda}_1\ge \hat{\lambda}_2\ge \cdots \ge \hat{\lambda}_{d}\ge 0$, with $\hat{\lambda}_i=0$ for $i>r_C$ \\
% $\delta_\kappa$ & Scalar & Cutoff gap, $\delta_\kappa:=\lambda_\kappa-\lambda_{\kappa+1}$ (when $\lambda_\kappa>\lambda_{\kappa+1}$) \\
% $\mathrm{gap}_i$ & Scalar & Spectral gap around $\lambda_i$, $\mathrm{gap}_i:=\min\{\lambda_{i-1}-\lambda_i,\ \lambda_i-\lambda_{i+1}\}$ \\
% \bottomrule
% \end{tabular}
% \end{table}

% % -----------------------------------------------------------------------------

% % ================================================================================

% \subsection{Probability Spaces and Measures}

% \begin{table}[H]
% \centering
% \begin{tabular}{cl}
% \toprule
% \textbf{Symbol} & \textbf{Description} \\
% \midrule
% $\mathcal{X}$ & State space, $\mathcal{X}=\mathbb{R}^n$ \\
% $\mathcal{D}$ & Stacked observation space, $\mathcal{D}=\mathbb{R}^{d}$ \\
% $\mathcal{B}_{\mathcal{X}}$ & Borel $\sigma$-algebra on $\mathcal{X}$ \\
% $\mathcal{B}_{\mathcal{D}}$ & Borel $\sigma$-algebra on $\mathcal{D}$ \\
% $\mu_{\mathcal{X}}$ & Lebesgue measure on $\mathcal{X}$ \\
% $\mu_{\mathcal{D}}$ & Lebesgue measure on $\mathcal{D}$ \\
% $\pi^{f}_{\mathcal{X}}$ & Forecast (prior) density on $\mathcal{X}$ \\
% 	$P_{\mathcal{X}}$ & Forecast distribution on $\mathcal{X}$; under Gaussian modeling $P_{\mathcal{X}}=\mathcal{N}(\boldsymbol{\mu}^f,\mathbf{P}^f)$ \\
% 	$Q:\mathcal{X}\to\mathcal{D}$ & QoI map; in QPCA-EnDCF, $Q(\mathbf{x})=\mathbf{H}^{(L)}\mathbf{x}$ \\
% 	$Q_\ell:\mathcal{X}\to\mathcal{D}$ & Approximate QoI map at approximation level $\ell$; a sequence $\{Q_\ell\}_{\ell\geq 1}$ \\
% 	& converging to $Q$ in $L^2(P_{\mathcal{X}})$ (e.g., reduced-order or surrogate observation operators) \\
% \bottomrule
% \end{tabular}
% \end{table}

% % -----------------------------------------------------------------------------

% % ================================================================================

% \subsection{Key Relationships and Definitions}

% \begin{itemize}
% \item \textbf{Whitened residuals:}
% \[
% \mathbf{e}^{(j)}
% :=
% (\mathbf{R}^{(L)})^{-1/2}\bigl(\mathbf{H}^{(L)}\mathbf{x}^{(j),f}-\mathbf{z}^{(w)}\bigr)
% \in \mathbb{R}^{d}.
% \]

% \item \textbf{Population mean and covariance of residuals:}
% \[
% \boldsymbol{\mu}_E:=\mathbb{E}[\mathbf{e}^{(j)}],
% \qquad
% \boldsymbol{\Sigma}_E:=\mathrm{Cov}(\mathbf{e}^{(j)})
% =
% (\mathbf{R}^{(L)})^{-1/2}\mathbf{H}^{(L)}\mathbf{P}^f(\mathbf{H}^{(L)})^\top(\mathbf{R}^{(L)})^{-1/2}.
% \]

% \item \textbf{Sample residual mean and covariance:}
% \[
% \bar{\mathbf{e}}:=\frac{1}{N}\sum_{j=1}^N \mathbf{e}^{(j)},
% \qquad
% \mathbf{C}_E:=\frac{1}{N-1}\sum_{j=1}^N (\mathbf{e}^{(j)}-\bar{\mathbf{e}})(\mathbf{e}^{(j)}-\bar{\mathbf{e}})^\top.
% \]

% \item \textbf{Stacked observation operator (linearized form):} $\mathbf{H}^{(L)}\in\mathbb{R}^{d\times n}$; for example,
% \[
% \mathbf{H}^{(L)}=
% \begin{bmatrix}
% \mathbf{H}\mathbf{M}_0\\
% \mathbf{H}\mathbf{M}_1\\
% \vdots\\
% \mathbf{H}\mathbf{M}_{L-1}
% \end{bmatrix},
% \]
% where $\mathbf{M}_\ell\in\mathbb{R}^{n\times n}$ is the tangent linear map (Jacobian of $\mathcal{M}^{(\ell)}$ evaluated about a reference state) mapping the initial state to time $\ell$ (static case: $\mathbf{M}_\ell=\mathbf{I}$). This linear form approximates the nonlinear QoI map $Q$ from Definition~\ref{def:data_space}.

% \item \textbf{Block-diagonal observation covariance:}
% \[
% \mathbf{R}^{(L)}=\mathbf{I}_L\otimes \mathbf{R}=\mathrm{diag}(\mathbf{R},\ldots,\mathbf{R}).
% \]

% \item \textbf{Ensemble mean and anomalies:}
% \[
% \bar{\mathbf{x}}=\frac{1}{N}\mathbf{X}\mathbf{1},
% \qquad
% \mathbf{A}=\mathbf{X}-\bar{\mathbf{x}}\mathbf{1}^\top.
% \]

% \item \textbf{Sample forecast covariance (distinct from population):}
% \[
% \widehat{\mathbf{P}}^{\,f}:=\frac{1}{N-1}\mathbf{A}\mathbf{A}^\top,
% \qquad
% \text{while }\mathbf{P}^f:=\mathrm{Cov}_{P_{\mathcal{X}}}(\mathbf{x}) \text{ denotes the population covariance.}
% \]

% \item \textbf{Rank constraints:}
% \[
% \mathrm{rank}(\mathbf{A})\le \min(n,N-1),
% \qquad
% \mathrm{rank}(\widehat{\mathbf{P}}^{\,f})\le N-1,
% \qquad
% \mathrm{rank}(\mathbf{C}_E)\le r=\min(d,N-1).
% \]

% \item \textbf{Empirical EnDCF covariances and gain (pseudoinverse form):}
% Let $\mathbf{X}_L\in\mathbb{R}^{n\times N}$ be the forecast ensemble at the assimilation time (or window endpoint) and
% $\mathbf{Y}_{\mathrm{stack}}\in\mathbb{R}^{d\times N}$ the stacked forecast observations. Define
% \[
% \mathbf{A}_x := \mathbf{X}_L-\bar{\mathbf{x}}_L^f\mathbf{1}^\top,
% \qquad
% \mathbf{A}_z := \mathbf{Y}_{\mathrm{stack}}-\bar{\mathbf{y}}_{\mathrm{stack}}\mathbf{1}^\top,
% \]
% and
% \[
% \mathbf{P}_{xz} := \frac{1}{N-1}\mathbf{A}_x\mathbf{A}_z^\top,
% \qquad
% \mathbf{P}_{zz} := \frac{1}{N-1}\mathbf{A}_z\mathbf{A}_z^\top.
% \]
% The EnDCF gain used in implementation is the \emph{empirical} gain
% \[
% \mathbf{K}^{\mathrm{DC}}
% :=
% \mathbf{P}_{xz}\,\mathbf{P}_{zz}^{\dagger},
% \]
% where $(\cdot)^{\dagger}$ denotes the Moore--Penrose pseudoinverse (so $\mathbf{K}^{\mathrm{DC}}$ is always well-defined, even when $N-1<d$).

% \item \textbf{Moment assumption for the (random) empirical gain:}
% When bounding mean-square errors in expectation, we assume the gain has finite second moment,
% \[
% \mathbb{E}\bigl[\|\mathbf{K}^{\mathrm{DC}}\|_2^2\bigr]<\infty.
% \]
% (Any stronger moment condition used later will be stated explicitly.)

% \item \textbf{QPCA truncation uses projectors (not eigenvalue tails by default):}
% For $1\le \kappa\le r$, the population and sample rank-$\kappa$ projectors $\mathbf{P}_\kappa$ and $\hat{\mathbf{P}}_\kappa$ (defined in the notation table) project onto the leading $\kappa$ eigenmodes.
% Truncation effects in the bias are governed by the alignment of the \emph{mean innovation} with the retained subspace
% (e.g.\ $\|(\mathbf{I}-\mathbf{P}_\kappa)\boldsymbol{\mu}_E\|^2$), and do \emph{not} generally reduce to an eigenvalue tail
% $\sum_{i>\kappa}\lambda_i$ without additional structural assumptions.
% \end{itemize}

% \FloatBarrier







\FloatBarrier
% ================================================================================
% BIBLIOGRAPHY
% ================================================================================
% Bibliography commented out until citations are added
\bibliographystyle{plainnat}
\bibliography{references}

\end{document}
